\providecommand{\spnHS}{\mathrm{HS}}
\providecommand{\spncodim}{\operatorname{codim}}
\providecommand{\spnim}{\operatorname{im}}
\providecommand{\spnE}{\mathbb{E}}
\providecommand{\spnP}{\mathbb{P}}
\providecommand{\spnTV}{d_{\mathrm{TV}}}
\providecommand{\spnStab}{\operatorname{Stab}}
\providecommand{\spnMat}{\operatorname{Mat}}
\providecommand{\spnDD}{\mathfrak{D}}
\providecommand{\spnT}{\mathsf{T}}
\providecommand{\spnqbin}[3]{\genfrac{[}{]}{0pt}{}{#1}{#2}_{#3}}
\providecommand{\spnOrth}{\mathcal{O}}
\providecommand{\spnTame}{\mathcal{T}}
\providecommand{\spnEmb}{\mathcal{E}}
\providecommand{\spnAlt}{\operatorname{Alt}}
\providecommand{\spncoker}{\operatorname{coker}}
\section{Compressions of random symplectic and orthogonal matrices}\label{sec:cmp}

\subsection{The positive-type Fourier estimate and matrix compressions}\label{spn:sec:positive-type}
\sources{grh_sharp_torsion_and_symplectic_blocks; grh_beyond_dirichlet_2026-11-25}

The spanning argument of the source contains a finite-field estimate for a
positive-type weight on a finite group (\code{card_mul_sum_det_le} in the
compact edition). It bounds one additive character sum through the masses of
two stabilizers and is then used to estimate the mass of singular square
compressions. Passing
from vectors to ordered frames controls the law of an entire rectangular
compression and the factorial moments of its kernel dimension.

Let $V=\F_q^n$, let $G$ be a finite group, and let $\pi\colon G\to\GL(V)$ be a
representation. The finite field may have any characteristic. A probability
mass function $\mu\colon G\to\R_{\ge0}$ is \emph{of positive type} if the
matrix $\bigl(\mu(g^{-1}h)\bigr)_{g,h\in G}$ is positive semidefinite over
$\C$. This is an assumption in addition to nonnegativity and total mass one;
uniform measure is one example. Write $g\cdot x=\pi(g)x$ and
$g\cdot y=y\circ\pi(g)^{-1}$ for the actions on $V$ and $V^*$, and let both act
diagonally on tuples. Fix a nontrivial additive character
$\psi\colon\F_q\to\C^\times$. Expectations $\spnE_\mu$ and probabilities
$\spnP_\mu$ refer to $g$ with law $\mu$.

\begin{lemma}[Positive-type coefficient bound; \stU]\label{spn:lem:positive-coefficient}
For $X=(x_1,\dots,x_k)\in V^k$ and $Y=(y_1,\dots,y_k)\in(V^*)^k$,
\begin{equation}\label{spn:eq:frame-coefficient}
 \abs[\Big]{\spnE_\mu\psi\Bigl(\sum_{i=1}^ky_i(\pi(g)x_i)\Bigr)}
 \le q^{nk/2}\sqrt{\mu(\spnStab X)\,\mu(\spnStab Y)} ,
\end{equation}
where the stabilizers fix every member of the tuple. No transitivity or
irreducibility is needed.
\end{lemma}

\begin{proof}
Applying the case $k=1$ to the representation on $V^k$ gives the general case,
so let $k=1$ and write $x=x_1$, $y=y_1$. Choose vectors $u_g$ in a complex
Hilbert space with $\langle u_g,u_h\rangle=\mu(g^{-1}h)$; a finite positive
semidefinite matrix always has such a Gram representation. Inner products are
conjugate-linear in the first argument. In $\ell^2(V)$ let $e_x$ be the point
masses and put $c_y(z)=q^{-n/2}\psi(-y(z))$. Additive-character orthogonality
makes both indexed families orthonormal. Consider the two finite tensor sums
\[
 P=\sum_{g\in G}c_{g\cdot y}\otimes u_g,\qquad
 Q=\sum_{h\in G}e_{h\cdot x}\otimes u_h .
\]
Setting $t=g^{-1}h$ in their inner product gives
\[
 \langle P,Q\rangle
 =\abs{G}q^{-n/2}\sum_{t\in G}\mu(t)\psi\bigl(y(\pi(t)x)\bigr).
\]
Orthonormality similarly gives $\norm{P}^2=\abs{G}\mu(\spnStab y)$ and
$\norm{Q}^2=\abs{G}\mu(\spnStab x)$. Cauchy--Schwarz proves the result.
\end{proof}

Let $A\colon\F_q^a\to V$ and $B\colon\F_q^b\to V^*$ be linear maps. Define the
$b\times a$ matrix $M(g)$ and its nullity by
\[
 M(g)_{ij}=(Be_i)(\pi(g)Ae_j),\qquad\nu(g)=\dim_{\F_q}\ker M(g);
\]
in coordinate identifications $M(g)=B^{\spnT}\pi(g)A$. For nonnegative
integers $d,j$, let
\[
 F_q(d,j)=\prod_{i=0}^{j-1}(q^d-q^i)\quad(0\le j\le d),\qquad
 F_q(d,j)=0\quad(j>d),
\]
the number of ordered independent $j$-frames in $\F_q^d$, with $F_q(d,0)=1$.
The number of $b\times j$ matrices of rank $k$ is
\begin{equation}\label{spn:eq:rank-matrix-count}
 R_q(b,j;k)=\frac{F_q(b,k)F_q(j,k)}{F_q(k,k)} .
\end{equation}
Indeed, factor a rank-$k$ matrix as a full-column-rank $b\times k$ matrix
times a full-row-rank $k\times j$ matrix, and divide by the
$\abs{\GL_k(\F_q)}$ changes of intermediate basis. When $A$ and $B$ are
injective, let $\alpha_k$ be the maximum $\mu$-mass of the pointwise
stabilizer of an independent $k$-frame in $\spnim A$, and let $\beta_k$ be
the corresponding maximum for independent $k$-frames in $\spnim B$ under the
dual action.

\begin{theorem}[Distribution of a compression; \stU]\label{spn:thm:tv}
Let $A,B$ be injective, let $g$ have a law $\mu$ of positive type, and let $Z$
be uniform on all $b\times a$ matrices over $\F_q$. Then
\begin{equation}\label{spn:eq:inherited-TV}
 \spnTV\bigl(M(g),Z\bigr)\le\frac12\Bigl(\sum_{k=1}^{\min(a,b)}
 R_q(b,a;k)\,q^{nk}\alpha_k\beta_k\Bigr)^{1/2}.
\end{equation}
\end{theorem}

\begin{proof}
A rank-$k$ matrix $W$ defining a nontrivial Fourier character
$M\mapsto\psi(\tr(W^{\spnT}M))$ factors as a sum of $k$ rank-one matrices
with independent left and independent right vectors. Its value on $M(g)$
becomes $\psi\bigl(\sum_iy_i(\pi(g)x_i)\bigr)$ for a corresponding independent
frame $X$ in $\spnim A$ and covector frame $Y$ in $\spnim B$.
\Cref{spn:lem:positive-coefficient} bounds the squared modulus of the Fourier
coefficient by $q^{nk}\alpha_k\beta_k$. The $\ell^1$--$\ell^2$ inequality and
Parseval give the displayed estimate, since the uniform distribution has zero
Fourier coefficient at every nontrivial character.
\end{proof}

The next two results extract a different observable, the kernel dimension.
They give direct rank-tail certificates which can remain useful when
\eqref{spn:eq:inherited-TV} is vacuous. Put
\[
 \alpha(x)=\mu(\spnStab x),\quad\beta(y)=\mu(\spnStab y),\qquad
 P_A=\sum_{v\in\F_q^a\setminus\set{0}}\sqrt{\alpha(Av)},\quad
 P_B=\sum_{y\in\F_q^b\setminus\set{0}}\sqrt{\beta(By)} .
\]

\begin{proposition}[Kernel-size moment; \stU]\label{spn:prop:kernel-first-moment}
Let $A,B$ be linear, not necessarily injective, and $\mu$ of positive type.
Then
\begin{equation}\label{spn:eq:first-kernel-moment}
 \abs[\Big]{\spnE_\mu q^{\nu(g)}-\Bigl(1+\frac{q^a-1}{q^b}\Bigr)}
 \le q^{n/2-b}P_AP_B .
\end{equation}
Consequently, for $1\le r\le a$,
\begin{equation}\label{spn:eq:first-kernel-tail}
 \spnP_\mu(\nu\ge r)\le\min\Bigl\{1,\
 \frac{(q^a-1)q^{-b}+q^{n/2-b}P_AP_B}{q^r-1}\Bigr\}.
\end{equation}
\end{proposition}

\begin{proof}
Count the kernel by additive-character orthogonality:
\[
 \spnE_\mu q^\nu=q^{-b}\sum_{v\in\F_q^a}\sum_{y\in\F_q^b}
 \spnE_\mu\psi\bigl((By)(\pi(g)Av)\bigr).
\]
Pairs for which $v=0$ or $y=0$ contribute exactly $q^{-b}(q^a+q^b-1)$. Bound
each remaining coefficient with \cref{spn:lem:positive-coefficient}, and
factor the two finite sums of square roots. This proves the two-sided moment
estimate. The nonnegative variable $q^\nu-1$ is at least $q^r-1$ on
$\set{\nu\ge r}$; Markov's inequality gives the tail estimate.
\end{proof}

For square compressions $a=b=m$, if every relevant nonzero-vector stabilizer
mass is at most $Y$, the moment estimate becomes
$\abs{\spnE_\mu q^\nu-(2-q^{-m})}\le q^{m+n/2}(1-q^{-m})^2Y$. The source's
singular-minor conclusion instead has the form
$q\,\spnP(\det M=0)\le2+q^{m+n/2}Y$. Keeping the zero-character contribution
exact and retaining the kernel-size moment gives information about every
deficiency threshold, rather than just singularity. The denominator $q^r-1$
is best possible if only the expectation of $q^\nu-1$ is known: an abstract
variable supported on $\set{0,r}$ attains equality in the Markov step.

Assume henceforth that $A$ and $B$ are injective, with $a,b\le n$. For
$1\le j\le a$ define
\begin{equation}\label{spn:eq:moment-main}
 U_j=\frac{F_q(a,j)}{q^{bj}},\qquad
 \eps_j=\sum_{k=1}^{\min(b,j)}R_q(b,j;k)\,q^{nk/2}\sqrt{\alpha_k\beta_k} .
\end{equation}

\begin{theorem}[Factorial kernel moments; \stU]\label{spn:thm:factorial-kernel}
For every $1\le j\le a$,
\begin{equation}\label{spn:eq:factorial-moment}
 \abs[\big]{\spnE_\mu F_q(\nu(g),j)-U_j}\le U_j\eps_j .
\end{equation}
For every $1\le j\le r\le a$ it follows that
\begin{equation}\label{spn:eq:factorial-tail}
 \spnP_\mu(\nu\ge r)\le\min\Bigl\{1,\ \frac{U_j(1+\eps_j)}{F_q(r,j)}\Bigr\},
\end{equation}
and one may minimize the right-hand side over $j$.
\end{theorem}

\begin{proof}
For each $g$, the quantity $F_q(\nu(g),j)$ counts full-column-rank matrices
$T\in\spnMat_{a\times j}(\F_q)$ for which $M(g)T=0$. For a fixed such $T$,
character orthogonality on all $b\times j$ matrices gives
\[
 \mathbf 1_{M(g)T=0}=q^{-bj}\sum_{W\in\spnMat_{b\times j}(\F_q)}
 \psi\bigl(\tr(W^{\spnT}M(g)T)\bigr).
\]
Take expectations and sum over the $F_q(a,j)$ choices of $T$. The term $W=0$
contributes exactly $U_j$.

For a nonzero $W$ of rank $k$, choose a factorization $W=CD$, where $C$ has
size $b\times k$ and rank $k$, and $D$ has size $k\times j$ and rank $k$.
Cyclicity of the trace yields
$\tr(W^{\spnT}M(g)T)=\tr(C^{\spnT}M(g)TD^{\spnT})$. The columns of
$ATD^{\spnT}$ form an independent $k$-frame in $\spnim A$: each of
$D^{\spnT}$, $T$ and $A$ is injective. The columns of $BC$ form an
independent $k$-frame in $\spnim B$. Inequality
\eqref{spn:eq:frame-coefficient} therefore bounds the absolute value of the
expected character by $q^{nk/2}\sqrt{\alpha_k\beta_k}$. This bound is
independent of $T$ and of the chosen factorization.

There are $R_q(b,j;k)$ possibilities for $W$ of rank $k$. The triangle
inequality, summed over $W$ and $T$, proves \eqref{spn:eq:factorial-moment}.
Finally, $F_q(\nu,j)$ is nonnegative and is at least $F_q(r,j)$ when
$\nu\ge r$. Apply Markov's inequality and the upper moment bound to obtain
\eqref{spn:eq:factorial-tail}.
\end{proof}

The normalization has an exact probabilistic meaning. If $Z$ is uniform on all
$b\times a$ matrices, annihilating any fixed independent $j$-frame imposes
$bj$ independent linear equations, hence
$\spnE F_q(\dim\ker Z,j)=F_q(a,j)q^{-bj}=U_j$. These are precisely the
factorial kernel moments of the independent-entry model. Counting frames is
equivalent, by duality, to counting surjections to $\F_q^j$; that moment
viewpoint has well-developed precedents \cite{Wood2019Random}, and exact
finite-field rank laws and total-variation comparisons have an established
literature \cite{Fulman2015Stein}.

\begin{remark}[The two certificates are incomparable]\label{spn:rem:tv-vs-moments}
The bound \eqref{spn:eq:inherited-TV} controls all matrix events
simultaneously. The moment bound \eqref{spn:eq:factorial-tail} instead sums
over frequency matrices of size $b\times j$, not $b\times a$; it requires
frame stabilizers only through rank $\min(b,j)$ and can focus on a particular
rank tail. The two bounds use different norms and quantities, and neither
dominates the other in all regimes. A small total-variation bound can be far
stronger. Conversely, take $G=\GL_{32}(\F_2)$ with uniform probability,
$a=31$, $b=4$, and coordinate injections $A,B$. The event $\nu\ge28$ is
exactly failure of full row rank. Transitivity on independent frames, also
for the dual action, gives $\alpha_k=\beta_k=1/F_2(32,k)$. Exact rational
arithmetic gives $\eps_8<0.058894$ and the following certificates (decimal
upper bounds rounded upward).
\begin{center}
\begin{tabular}{ll}
\toprule
Quantity & Value or certificate\\
\midrule
First moment, $j=1$ & $\spnP(\nu\ge28)\le0.500115$\\
Eighth moment, $j=8$ & $\spnP(\nu\ge28)\le0.004137$\\
Uncapped right side of \eqref{spn:eq:inherited-TV} & $>2.132$, hence vacuous\\
Exact probability in this example & $15/(2^{32}-1)$\\
\bottomrule
\end{tabular}
\end{center}
The exact value holds because $\pi(g)\spnim A$ is a uniform hyperplane in
$\F_2^{32}$, and projection to $\F_2^4$ loses one row of rank exactly when
that hyperplane contains the fixed $28$-dimensional projection kernel; there
are $2^{32}-1$ hyperplanes in total and $2^4-1$ containing this kernel, and
no larger loss is possible. So the example does not improve the specialized
exact $\GL_n$ calculation; the comparison is between general certificates
using the same stabilizer information.
\end{remark}

\begin{example}[Positive type is essential]\label{spn:ex:positive-type}
Take $G=\GL_2(\F_2)$, $a=b=1$, $A=B=e_1$, and put all probability on
$s=\left(\begin{smallmatrix}0&1\\1&0\end{smallmatrix}\right)$. This weight is
nonnegative, normalized and inverse-symmetric. Both relevant stabilizer
masses are zero, but $M(s)=0$ with probability one. Thus $\spnE2^\nu=2$,
whereas the center in \eqref{spn:eq:first-kernel-moment} is $3/2$ and its
asserted error would be zero. The law is not of positive type: the principal
Gram matrix indexed by $\set{1,s}$ is
$\left(\begin{smallmatrix}0&1\\1&0\end{smallmatrix}\right)$, with determinant
$-1$. On the other hand, for any subgroup $H$ the kernel
$\mathbf 1_H(g^{-1}h)/\abs{H}$ is a positive semidefinite block matrix on
cosets, and convex combinations preserve positivity; for instance
$\tfrac12\mathrm{Unif}(\GL_4(\F_2))+\tfrac12\mathrm{Unif}(\Sp_4(\F_2))$ is
of positive type, with weights $29/40320$ on $\Sp_4$ and $1/40320$
elsewhere. For a square $2\times2$ compression its nullity law is
$(52/105,\,16/35,\,1/21)$ on $\nu=0,1,2$, against $(16/35,\,18/35,\,1/35)$
for uniform $\GL_4(\F_2)$. All $636$ relevant ordered-frame Fourier
inequalities for these two laws, the main terms and every stated finite
moment and tail are verified in exact
arithmetic.\footnote{\artifact{grh_beyond_dirichlet_2026-11-25/data/rank_moments_results.json}}
\end{example}

%======================================================================




\subsection{Which blocks of a random symplectic matrix look uniform?}\label{spn:sec:symplectic}
\sources{grh_sharp_torsion_and_symplectic_blocks}

In a symplectic space a change of basis preserves alternating pairings. Over
a fixed finite field, \cref{spn:thm:tv} gives a uniform range for all
rectangular compressions, and for isotropic blocks at a fixed distance from
the boundary the limiting total variation distance can be determined
exactly. Throughout, $n$ is even, $g$ is uniform in $\Sp_n(\F_q)$,
$A\colon\F_q^a\to V$ and $B\colon\F_q^b\to V^*$ are injective, $M(g)$ is the
compression of \cref{spn:sec:positive-type}, and $Z$ is a uniform $b\times a$
matrix. Write $\spnqbin{m}{i}{q}$ for the Gaussian binomial coefficient.

\begin{theorem}[Symplectic compression; \stU]\label{spn:thm:symplectic}
For arbitrary $A,B$ as above with $a+b\le n$, an absolute constant $C$ gives
\begin{equation}\label{spn:eq:symplectic-upper}
 \spnTV(M(g),Z)\le C\,q^{-(n+1-a-b)/2} .
\end{equation}
Thus $n-a-b\to\infty$ suffices at fixed $q$, uniformly in both subspaces. For
every fixed integer $c\ge0$, the dimensions $n=2m$, $a=m$, $b=m-c$ admit
isotropic choices such that the distance stays bounded away from zero as
$m\to\infty$.
\end{theorem}

\begin{proof}[Proof of the upper bound]
The symplectic form identifies the dual action with the given one. By Witt
extension, the orbit of an independent $j$-frame consists of all independent
frames with its prescribed alternating Gram matrix. Once $i-1$ independent
vectors have been specified, the pairings with the next vector give $i-1$
independent linear constraints. Of the $q^{n-i+1}$ solutions at most
$q^{i-1}$ are in the previous span. Since $j\le n/2$, every orbit has size at
least
\[
 q^{nj-j(j-1)/2}\prod_{i=1}^j(1-q^{2i-2-n})
 \ \ge\ q^{nj-j(j-1)/2}\prod_{k\ge1}(1-2^{-2k}).
\]
Each stabilizer mass is the reciprocal of its orbit size. By
\eqref{spn:eq:rank-matrix-count},
$R_q(b,a;j)\le q^{j(a+b-j)}/\prod_{k\ge1}(1-2^{-k})$. Substitution in
\eqref{spn:eq:inherited-TV} leaves a geometric sum
$\sum_{j\ge1}q^{-j(n+1-a-b)}$, as claimed.
\end{proof}

\begin{proof}[Proof of the obstruction]
Under the symplectic identification take $\spnim A$ Lagrangian and
$\spnim B=W_0$ isotropic of dimension $m-c$. Then $L=g\spnim A$ is a uniform
Lagrangian. Full row rank of $M(g)$ means $L\cap W_0=0$, since $L^\perp=L$.
The number of Lagrangians in dimension $2m$ is $\prod_{k=1}^m(q^k+1)$, as
follows by counting pairs of a Lagrangian and a line in it and passing to the
line's symplectic quotient. A fixed $i$-dimensional subspace of $W_0$ lies in
$L$ with probability $1/\prod_{k=m-i+1}^{m}(q^k+1)$, by the same quotient
argument. Subspace-lattice M\"obius inversion gives the full-rank probability
\begin{equation}\label{spn:eq:Pmc}
 P_{m,c}=\sum_{i=0}^{m-c}
 \frac{(-1)^iq^{i(i-1)/2}\spnqbin{m-c}{i}{q}}{\prod_{k=m-i+1}^{m}(1+q^k)},
\end{equation}
where the Gaussian coefficient counts $i$-subspaces of $W_0$ and the M\"obius
function in dimension $i$ is $(-1)^iq^{i(i-1)/2}$. For $c=0$ the probability
also has the elementary form $\prod_{k=1}^m(1+q^{-k})^{-1}$, since there are
$q^{m(m+1)/2}$ Lagrangian complements, parametrized by symmetric matrices.

Set $Q=q^{-1}$. The $i$th term converges to
$(-Q^{c+1})^i/\prod_{k=1}^i(1-Q^k)$ and is bounded in absolute value by
$Q^{(c+1)i}/\prod_{k\ge1}(1-Q^k)$; these facts follow by extracting the
highest powers of $q$ in the Gaussian coefficient and the denominator.
Dominated convergence and the power-series identity
$\sum_{i\ge0}z^i/\prod_{k=1}^i(1-Q^k)=\prod_{k\ge0}(1-zQ^k)^{-1}$ for
$\abs{z}<1$ yield
\[
 \lim_{m\to\infty}P_{m,c}=\prod_{k=c+1}^\infty(1+q^{-k})^{-1}.
\]
For the uniform rectangular matrix the full-rank probability is
$\prod_{k=c+1}^{m}(1-q^{-k})$. The two positive limits have ratio
$\prod_{k=c+1}^\infty(1-q^{-2k})^{-1}>1$. Their difference is a lower bound
for the limit inferior of the total variation distance.
\end{proof}

\begin{theorem}[Exact limiting distance for the isotropic family; \stU]\label{spn:thm:symplectic-exact}
Let $n=2m$, $a=m$, $b=m-c$ with $\spnim A$ Lagrangian and $\spnim B$
isotropic, and fix $q$ and $c$. As $m\to\infty$, the total variation distance
between $M(g)$ and $Z$ converges to
\begin{equation}\label{spn:eq:symplectic-limit}
 \frac{Q^{c+1}}{1-Q}\Bigl(\prod_{j=c+2}^{\infty}(1-Q^j)
 -\frac{1}{\prod_{j=c+1}^{\infty}(1+Q^j)}\Bigr),\qquad Q=q^{-1}.
\end{equation}
At fixed $q$, this is asymptotic to $Q^{2c+2}/(1-Q)$ as $c\to\infty$. The
event that the matrix loses exactly one row of rank accounts for the entire
distance in the limit.
\end{theorem}

\begin{proof}
Witt extension realizes every change of coordinates in each of the two
isotropic subspaces by a symplectic automorphism. Left and right invariance
of uniform $g$ therefore make its compression uniform conditional on its
rank, since left and right general linear transformations act transitively
on each rank. The uniform rectangular model has the same conditional laws.
Hence total variation of the full matrices equals total variation of their
row deficiencies $\mathrm{def}=b-\rank M$.

Let $s_m(k)$ and $u_m(k)$ be the two probabilities of deficiency $k$. Each
possible intersection $D=L\cap W_0$ of dimension $k$ gives a symplectic
quotient $D^\perp/D$ of dimension $2(m-k)$, with the image of $W_0$ isotropic
of dimension $m-k-c$. Conditional on $D\subseteq L$, the Lagrangian $L/D$ is
uniform there. In the rectangular model, specifying a $k$-dimensional left
kernel first requires the random map $\F_q^b\to\F_q^m$ to vanish on it, then
to be injective on its quotient. Therefore, for $0\le k\le m-c$,
\[
 s_m(k)=\spnqbin{m-c}{k}{q}\frac{P_{m-k,c}}{\prod_{j=m-k+1}^{m}(1+q^j)},
 \qquad
 u_m(k)=\spnqbin{m-c}{k}{q}q^{-mk}\prod_{j=c+k+1}^{m}(1-Q^j).
\]
Both families are tight on the nonnegative integers. The expected number of
lines in the left kernel is $\spnqbin{m-c}{1}{q}/(q^m+1)$ in the symplectic
model, by the one-line quotient count, and $\spnqbin{m-c}{1}{q}q^{-m}$ for a
uniform matrix. Both remain bounded as $m\to\infty$, while a $k$-dimensional
kernel contains $(q^k-1)/(q-1)$ lines. Pointwise convergence of both mass
functions, which follows from the displayed formulas, thus implies
convergence in total variation to their respective limiting laws.

Write $s(k),u(k)$ for the limits and $R_k=s(k)/u(k)$, with $u(k)>0$. Taking
ratios gives
\[
 R_0=\prod_{j=c+1}^{\infty}(1-Q^{2j})^{-1}>1,\qquad
 \frac{R_{k+1}}{R_k}=Q^{-k}(1-Q^{c+k+1}).
\]
Set $t=Q^{c+1}$. Retaining the first two factors of $R_0$ yields
$R_2\ge[Q(1+t)(1+Qt)]^{-1}>1$, since $Q\le1/2$ and $t\le Q$. For $k\ge1$ the
successive ratio is at least $Q^{-1}(1-Q^2)>1$. All deficiencies except
possibly $k=1$ therefore have excess mass in the symplectic law, with strict
excess at $k=0$. As both limiting laws have total mass one, the distance is
exactly $u(1)-s(1)$. Substitution gives \eqref{spn:eq:symplectic-limit}.
Finally, $\prod_{j=c+2}^\infty(1-Q^j)=1-Qt/(1-Q)+O_q(t^2)$ and
$\prod_{j=c+1}^\infty(1+Q^j)^{-1}=1-t/(1-Q)+O_q(t^2)$ by the convergent
logarithmic expansions. Their difference is $t+O_q(t^2)$, proving the
asymptotic.
\end{proof}

\begin{example}\label{spn:ex:symplectic-F2}
For a square block over $\F_2$ with isotropic test spaces ($q=2$, $c=0$), the
limiting probability of losing exactly one row of rank is $0.419422$ in the
symplectic model and $0.577576$ for a uniform binary matrix. Their difference
$0.158154$ is the limiting total variation distance for the entire matrix.
Small cases of the isotropic frame counts are checked by exact
enumeration.\footnote{\artifact{grh_sharp_torsion_and_symplectic_blocks/checks/results.txt}}
\end{example}

The Lagrangian intersection count used when $c=0$ also appears in
\cite[Corollary~2]{Kueng2015Qubit}. For random $\GL_n$ blocks,
\cite[Theorem~1.8]{Gorodetsky2026Statistics} shows that
$n-\min(a,b)\to\infty$ suffices. Along the symplectic family just constructed
this quantity grows without bound, yet the full-row-rank test detects a
nonvanishing discrepancy. The boundary calculation quantifies the discrepancy
at each fixed deficit. For arbitrary restricted forms on both subspaces, and
for any number of frames, the law is determined exactly in the next
subsection. When one of the two subspaces is not isotropic, the law
conditional on rank is no longer uniform, and rank can miss almost all of the
discrepancy (\cref{spn:sec:lagrangian,spn:sec:nondegenerate}); many frames are
compared in \cref{spn:sec:many}, and symplectic frames are lifted to every
congruence precision in \cref{spn:sec:padic}.

%======================================================================

\subsection{The complete law of an arbitrary symplectic compression}\label{spn:sec:complete-law}
\sources{symplectic_compression_2026-11-26; prime_endpoints_and_symplectic_compressions; exact_symplectic_compressions_2026-11-26; symplectic_compression_2026-11-26_2; symplectic_joint_observations_2026-11-27}

The two preceding theorems bound or compute the discrepancy for special pairs
of subspaces. Counting fibres instead of estimating Fourier coefficients gives
the probability of every matrix of pairings between one fixed frame and a
uniform symplectic translate of another, for arbitrary restricted alternating
forms on both frames and arbitrary rectangular sizes, and more generally of all
pairings among any number of independent frames. In this subsection $q$
is any prime power, $V$ has dimension $2n$ with a nondegenerate alternating
form $\langle\ ,\ \rangle$ of matrix $J$, and an alternating form is a bilinear
form $B$ with $B(x,x)=0$ for all $x$; its matrix has zero diagonal as well as
$H^{\spnT}=-H$, and both conditions matter in characteristic two. Fix injective
linear maps $Y\colon\F_q^a\to V$ and $X_0\colon\F_q^b\to V$ with Gram matrices
$H_1=Y^{\spnT}JY$ and $H_0=X_0^{\spnT}JX_0$ of ranks $2s_1$ and $2s_0$. For
uniform $g\in\Sp(V)$ the observation is
\begin{equation}\label{spn:eq:complete-observation}
 M(g)=Y^{\spnT}JgX_0\in\spnMat_{a,b}(\F_q).
\end{equation}
Neither form need be zero or nondegenerate. Two such frames exist exactly when
$a-s_1\le n$ and $b-s_0\le n$. Every independent family of $a$ output
functionals has the form $Y^{\spnT}J$, so \eqref{spn:eq:complete-observation}
is the compression of \cref{spn:sec:positive-type} between arbitrary
subspaces.

For integers $d,s\ge0$ put
\begin{equation}\label{spn:eq:E-count}
 \spnEmb_n(d,s)=q^{2nd-\binom d2}\prod_{j=n-d+s+1}^{n}(1-q^{-2j})
 \quad\text{if }0\le2s\le d\text{ and }d-s\le n,
\end{equation}
and $\spnEmb_n(d,s)=0$ otherwise; empty products are one, so
$\spnEmb_n(0,0)=1$.

\begin{lemma}[Prescribed-Gram frames; \stU]\label{spn:lem:gram-frames}
Let $E$ be a $d$-dimensional alternating space of rank $2s$. The number of
isometric injections $E\to V$ is $\spnEmb_n(d,s)$, and $\Sp(V)$ acts
transitively on them whenever there is one.
\end{lemma}

\begin{proof}
Write $E=S\oplus R$ with $S$ nondegenerate of dimension $2s$ and $R$ the
radical, of dimension $h=d-2s$. A symplectic basis of $S$ embeds in
$\abs{\Sp_{2n}(\F_q)}/\abs{\Sp_{2(n-s)}(\F_q)}$ ways. Once its image is fixed,
an ordered basis of $R$ must go to an independent isotropic $h$-frame in the
orthogonal complement, which is symplectic of dimension $2(n-s)$; after $i$
radical vectors have been chosen, the next lies in their orthogonal complement
and outside their span, which leaves $q^{2(n-s)-i}-q^i$ choices. So the count is
nonzero exactly when $h\le n-s$, and
$\abs{\Sp_{2u}(\F_q)}=q^{u^2}\prod_{j=1}^u(q^{2j}-1)$ turns it into
\eqref{spn:eq:E-count}, the remaining factors having the consecutive indices
$n-d+s+1,\dots,n$. For transitivity, extend the image of the symplectic basis
of $S$ to a symplectic decomposition of $V$. In its complement, extend the
ordered isotropic basis $z_1,\dots,z_h$ of the image of $R$ to a symplectic
basis: independence of the functionals $\langle z_i,\cdot\rangle$ permits
choosing a partner of $z_1$ pairing to one with $z_1$ and to zero with the
later $z_i$; remove that hyperbolic plane, continue, and complete the residual
complement by symplectic pairs. The change of symplectic basis between two
such completions carries one prescribed frame to the other. This is the
alternating case of Witt's extension theorem, and the argument works in
characteristic two.
\end{proof}

\paragraph{Any number of frames.}
The law of one compression is the two-frame case of an exact law for the complete set of pairings
among any number of independent frames. For $1\le i\le k$ let $E_i$ be a space of dimension
$d_i\ge1$ carrying an alternating form $H_i$ of rank $2s_i$, with $d_i-s_i\le n$, and let
$X_i\colon E_i\to V$ be independent uniform isometric injections. These are ordered frames, not
subspaces, and their images may meet. Put $E=\bigoplus_iE_i$ and $D=\sum_id_i$. The observation
\[
 M_{ij}(u,v)=\langle X_iu,X_jv\rangle\qquad(i<j)
\]
is equivalently the alternating form $\Gamma$ on $E$ with diagonal blocks $H_i$ and upper
off-diagonal blocks $M_{ij}$; it has at most $q^B$ values, where
$B=\sum_{i<j}d_id_j=\binom D2-\sum_i\binom{d_i}2$. For a candidate $\Gamma$ write
\[
 \rank\Gamma=2s,\quad R=\operatorname{rad}\Gamma,\quad r=\dim R=D-2s,\quad R_i=R\cap E_i,\quad
 t_i=\dim R_i .
\]
The spaces $R_i$ form a direct sum, which need not be all of $R$. Let $T_\ell$ be the number of
$\ell$-dimensional subspaces $L\le R$ with $L\cap R_i=0$ for every $i$.

\begin{theorem}[Exact joint likelihood; \stU]\label{spn:thm:joint-law}
For every finite field and every candidate $\Gamma$,
\begin{equation}\label{spn:eq:joint-mobius}
 T_\ell=\sum_{0\le j_i\le t_i}(-1)^Jq^{\sum_i\binom{j_i}2}\Bigl(\prod_i\spnqbin{t_i}{j_i}q\Bigr)
 \spnqbin{r-J}{\ell-J}q,\qquad J=\sum_ij_i,
\end{equation}
and
\begin{equation}\label{spn:eq:joint-law}
 \spnP(\text{the observation is }\Gamma)=
 \frac{\sum_{\ell=0}^rT_\ell\,\spnEmb_n(D-\ell,s)}{\prod_i\spnEmb_n(d_i,s_i)} ;
\end{equation}
the $\ell$th summand divided by the same denominator is the probability of $\Gamma$ together with an
$\ell$-dimensional space of linear relations among the combined frame columns. If
$\prod_{i=1}^k\prod_{h=0}^{t_i-1}(1-q^hz)=\sum_Jb_Jz^J$, then
$T_\ell=\sum_{J\le\ell}b_J\spnqbin{r-J}{\ell-J}q$, so once Gaussian elimination has produced the
invariants of $\Gamma$ the likelihood costs $O(D^2+\log(n+1))$ further arithmetic operations.
\end{theorem}

\begin{proof}
For a finite vector space $W$ of dimension $t$, the sum
$\sum_{U\le W}(-1)^{\dim U}q^{\binom{\dim U}2}$ is one if $t=0$ and zero otherwise: it is the value at
$z=1$ of $\prod_{h=0}^{t-1}(1-q^hz)=\sum_j(-1)^jq^{\binom j2}\spnqbin tjqz^j$, which follows by
induction from the recurrence of the Gaussian coefficients. Apply this to each intersection
$L\cap R_i$ and multiply the indicators. On summing over the $\ell$-subspaces $L\le R$, subspaces
$U_i\le R_i$ of dimensions $j_i$ span a direct sum of dimension $J$, which lies in
$\spnqbin{r-J}{\ell-J}q$ of them; this is \eqref{spn:eq:joint-mobius}, and collecting the terms of
total index $J$ gives the polynomial form.

A tuple of frames is a linear map $F\colon E\to V$, $F(v_1,\dots,v_k)=\sum_iX_iv_i$, with Gram form
$\Gamma$ and injective restrictions to every $E_i$. Its kernel $L$ lies in $R$, since
$\Gamma(z,w)=\langle Fz,Fw\rangle$ vanishes for all $w$ when $Fz=0$, and injectivity on the summands
says $L\cap R_i=0$ for all $i$. Conversely each such $L$ of dimension $\ell$ and each isometric
injection of $E/L$, which has dimension $D-\ell$ and rank $2s$, determine exactly one admissible
tuple with kernel $L$. \Cref{spn:lem:gram-frames} counts these injections, and division by the
number $\prod_i\spnEmb_n(d_i,s_i)$ of frame tuples gives \eqref{spn:eq:joint-law} and its
refinement. For the cost, the product has at most $r$ linear factors; it, the table of Gaussian
coefficients and the triangular sums take $O(D^2)$ operations each, and after cancellation the
ratio of embedding counts has at most $D$ factors, computed from $q^{2n}$ and a table of powers of
$q$.
\end{proof}

For fixed individual forms the likelihood is therefore a function of the $k+1$ integers
$(s,t_1,\dots,t_k)$, which retain all the information for a likelihood test against independent
uniform cross blocks; the separate ranks of the cross blocks need not.

\begin{theorem}[Simultaneous realization and least ambient dimension; \stU]\label{spn:thm:joint-support}
An alternating form $\Gamma$ on $E=\bigoplus_iE_i$, with arbitrary diagonal blocks, is the Gram form
of a map into a symplectic space of dimension $2n$ that is injective on every $E_i$ if and only if
\begin{equation}\label{spn:eq:joint-support}
 s+\max_it_i\le n .
\end{equation}
The least such dimension is $2(s+\max_it_i)$. The statement does not presuppose that the individual
forms are realizable.
\end{theorem}

\begin{proof}
If $L\le R$ is the kernel of a realization and $\ell=\dim L$, then $L\cap R_i=0$ gives
$\ell\le r-t_i$ for every $i$, and the quotient, of dimension $D-\ell$ and rank $2s$, embeds only if
$n\ge D-\ell-s\ge D-r+\max_it_i-s=s+\max_it_i$ (\cref{spn:lem:gram-frames}). Conversely put
$t=\max_it_i$, choose injections $f_i\colon R_i\to\F_q^t$, and extend their sum to
$f\colon R\to\F_q^t$ by zero on a complement of $\bigoplus_iR_i$; $f$ is onto because one $f_i$ is,
also when $t=0$. Its kernel $L$ has dimension $r-t$ and meets every $R_i$ trivially. The quotient
$E/L$ has rank $2s$ and a radical of dimension $t$: send its nondegenerate part to $s$ standard
hyperbolic pairs and its radical to one vector from each of $t$ further hyperbolic pairs, and compose
with the quotient map. Every summand stays injective because $L\cap E_i=L\cap R_i=0$.
\end{proof}

The proof constructs a smallest realization by linear algebra, without enumerating a symplectic
group. If every $H_i$ is nondegenerate, all $R_i$ vanish and feasibility is simply
$\rank\Gamma\le2n$; for degenerate forms the radical intersections matter. A two-frame implementation
of the construction, which outputs and checks the frames, passes $805$ exhaustive cross-matrix
cases.\footnote{\artifact{symplectic_compression_2026-11-26/results/minimal_realization.json}}
Enumeration of $733{,}531$ frame tuples in nine configurations with three or four frames over $\F_2$
and $\F_3$, all $753$ possible observations included, agrees with \cref{spn:thm:joint-law} and
\eqref{spn:eq:joint-support}, and a second implementation by the polynomial form agrees at every
observation.\footnote{\raggedright\artifact{symplectic_joint_observations_2026-11-27/results/exhaustive_checks.json},
\artifact{symplectic_joint_observations_2026-11-27/results/verification_summary.json}}

\paragraph{Two frames.}
The compression \eqref{spn:eq:complete-observation} is the case of two frames, with $E_1=\F_q^a$
carrying $H_1$ and $E_0=\F_q^b$ carrying $H_0$. For a proposed value $M\in\spnMat_{a,b}(\F_q)$ form the alternating
matrix
\begin{equation}\label{spn:eq:combined-gram}
 \Gamma_M=\begin{pmatrix}H_1&M\\-M^{\spnT}&H_0\end{pmatrix}
\end{equation}
on $E=E_1\oplus E_0=\F_q^a\oplus\F_q^b$, write $\rank\Gamma_M=2s$ and
$r=a+b-2s$, let $R=\operatorname{rad}\Gamma_M$, and put
\[
 t_1=\dim(R\cap E_1)=\dim(\ker H_1\cap\ker M^{\spnT}),\qquad
 t_0=\dim(R\cap E_0)=\dim(\ker H_0\cap\ker M),
\]
and $c=r-t_0-t_1$. For two frames the relation count of \cref{spn:thm:joint-law} also has a direct
evaluation. With $\spnqbin uvq$ the number of $v$-subspaces of $\F_q^u$ (zero outside
$0\le v\le u$) and $\abs{\GL_0(\F_q)}=1$, define for $0\le\ell\le r$ the
nonnegative integer
\begin{equation}\label{spn:eq:T-count}
 T_\ell(r,t_0,t_1)=\sum_{k=\max(0,\ell-c)}^{\min(\ell,t_0,t_1)}
 \spnqbin{t_0}kq\spnqbin{t_1}kq\abs{\GL_k(\F_q)}\spnqbin{c}{\ell-k}q
 q^{(\ell-k)(t_0+t_1-k)} .
\end{equation}

\begin{lemma}[Relations meeting neither frame; \stU]\label{spn:lem:transverse}
The number of $\ell$-dimensional subspaces $L\le R$ with
$L\cap E_0=L\cap E_1=0$ is $T_\ell(r,t_0,t_1)$. Such a subspace exists exactly
when $0\le\ell\le r-\max(t_0,t_1)=c+\min(t_0,t_1)$.
\end{lemma}

\begin{proof}
Put $R_i=R\cap E_i$, choose a complement $C$ of $R_0\oplus R_1$ in $R$, and set
$J_L=L\cap(R_0\oplus R_1)$. Both coordinate projections of $J_L$ are
injective, so if $\dim J_L=k$ it is the graph of an isomorphism between a
$k$-subspace of $R_0$ and a $k$-subspace of $R_1$; there are
$\spnqbin{t_0}kq\spnqbin{t_1}kq\abs{\GL_k(\F_q)}$ choices. The projection $D$
of $L$ to $C$ has dimension $\ell-k$, with $\spnqbin c{\ell-k}q$ choices. For
fixed $J_L$ and $D$, the possible $L/J_L$ are the graphs of the linear maps
$D\to(R_0\oplus R_1)/J_L$, of which there are $q^{(\ell-k)(t_0+t_1-k)}$; all of
them meet $R_0\oplus R_1$ in $J_L$ and meet neither $R_i$. The inequality
$\ell+t_i\le r$ is necessary. It is sufficient: take the graph of an
isomorphism between subspaces of $R_0$ and $R_1$ of dimension
$\min(t_0,t_1)$, add $C$, and pass to any smaller subspace.
\end{proof}

\begin{corollary}[Complete compression law; \stU]\label{spn:thm:complete-law}
Over every finite field and for every $M\in\spnMat_{a,b}(\F_q)$,
\begin{equation}\label{spn:eq:complete-law}
 \spnP\bigl(M(g)=M\bigr)=
 \frac{\sum_{\ell=0}^{r}T_\ell(r,t_0,t_1)\,\spnEmb_n(a+b-\ell,s)}
 {\spnEmb_n(a,s_1)\,\spnEmb_n(b,s_0)} ,
\end{equation}
and the $\ell$th summand divided by the same denominator is
$\spnP\bigl(M(g)=M,\ \dim(\spnim Y\cap\spnim gX_0)=\ell\bigr)$. For fixed
$n,q,H_0,H_1$ the likelihood is a function of the triple $(s,t_0,t_1)$. The
support is exactly
\begin{equation}\label{spn:eq:complete-support}
 \spnP\bigl(M(g)=M\bigr)>0\iff s+\max(t_0,t_1)\le n .
\end{equation}
More generally, for arbitrary alternating $H_1,H_0$ and an arbitrary cross
matrix $M$, with no frames assumed to exist, the least dimension of a
nondegenerate alternating space containing injective realizations of both
forms with cross pairing $M$ is $2(s+\max(t_0,t_1))$.
\end{corollary}

\begin{proof}
By \cref{spn:lem:gram-frames}, $gX_0$ is uniform among its
$\spnEmb_n(b,s_0)$ possible values. A symplectic change of coordinates
carrying one target frame to another carries the source frames giving a
prescribed $M$, and a prescribed intersection dimension, bijectively onto each
other. So the target may also be taken uniform among its $\spnEmb_n(a,s_1)$
values, independently, and the law is the case of two frames of
\cref{spn:thm:joint-law}, whose relation count is evaluated by
\cref{spn:lem:transverse}. The map $(u,v)\mapsto Yu=-Xv$ identifies the kernel of
$F(u,v)=Yu+Xv$ with $\spnim Y\cap\spnim X$, which gives the refinement. The support
and the least dimension are the case of two frames of \cref{spn:thm:joint-support}.
\end{proof}

Both the symplectic law and the uniform law are constant on each class of
matrices sharing a triple $(s,t_0,t_1)$, so the triple is sufficient for every
likelihood test, and the total variation distance between the two matrix laws
equals the distance between the induced laws of the triple; different triples
may still give equal likelihoods. Conditional on an observed $M$ of positive
probability, the intersection dimension has law proportional to
$T_\ell\spnEmb_n(a+b-\ell,s)$, and given also $\ell$ all admissible relation
spaces $L$ are equally likely. Three special cases are worth recording.
\begin{itemize}
\item If $H_0$ and $H_1$ are nondegenerate, then $t_0=t_1=0$,
$T_\ell=\spnqbin r\ell q$, and the likelihood depends only on
$\rank\Gamma_M=a+\rank(H_0+M^{\spnT}H_1^{-1}M)$, the rank of an alternating
residual form rather than the rank of $M$.
\item If both frames are isotropic, then $\rank\Gamma_M=2\rank M$,
$t_1=a-\rank M$ and $t_0=b-\rank M$, so the likelihood depends only on
$\rank M$, as used in the proof of \cref{spn:thm:symplectic-exact}.
\item If $a=b=n$ and $H_1=0$, the target is Lagrangian; with $K=\ker M$ one has
$\rank\Gamma_M=2\rank M+\rank(H_0|_{K\times K})$, $t_1=n-\rank M$ and
$t_0\le\dim K=t_1$, so \eqref{spn:eq:complete-support} says that $M$ occurs
exactly when $K$ is $H_0$-isotropic. \Cref{spn:thm:lagrangian} gives the
likelihood in closed form in this case.
\end{itemize}
Evaluating one likelihood needs Gaussian elimination for the invariants and
$O((a+b)^2)$ further integer operations, without enumerating the group. The
law was compared with direct enumeration of prescribed-Gram source frames in
$33$ configurations, covering $2{,}305$ matrix values and $4{,}279$ joint
cells, including all frame types through dimension three in a
four-dimensional binary space, cases over $\F_3$ and genuine extension-field
cases over $\F_4$; an independent implementation checks $1{,}213$ further
matrix values, a separate checker verifies \cref{spn:lem:transverse} on
$3{,}614$ subspaces, and a second exhaustive check over $\F_2$ and $\F_3$ in
ambient dimensions two to six finds, for instance, exactly $174$ of the $512$
binary $3\times3$ matrices supported for two three-dimensional frames of
rank two in dimension four. Two further independent enumerations agree: one
of $329{,}957$ frame-pair observations, with $35{,}230$ matrix outcomes and
$130{,}773$ joint matrix and intersection outcomes, over $\F_2$, $\F_3$ and
$\F_4$ in ambient dimensions up to eight, and one of $125{,}232$ complete
matrix outcomes in $87$ form types over $\F_2$ and
$\F_3$.\footnote{\raggedright\artifact{symplectic_compression_2026-11-26/results/complete_law.json},
\artifact{symplectic_compression_2026-11-26/results/independent_frame_audit.json},
\artifact{symplectic_compression_2026-11-26/results/radical_transversality.json},
\artifact{prime_endpoints_and_symplectic_compressions/compression/verification.txt},
\artifact{exact_symplectic_compressions_2026-11-26/results/exhaustive_checks.json},
\artifact{symplectic_compression_2026-11-26_2/results.json}}
The orbit classification of a tuple by its kernel and induced form is
classical; see \cite[Lemma~2.6]{GuralnickLarsenTiep2024Character} and, for
related counts of pairs of subspaces with prescribed intersection and
nondegeneracy, \cite{DeBoeckVandeVoorde2024Anzahl}. Counts of
representations of alternating forms by alternating forms, including
prescribed rank, go back to \cite{Wan1995Representations,WeiZhang2001Number},
and Witt's lemma for degenerate forms in every characteristic is
\cite[Theorem~3.4]{SprehnWahl2020Forms}.

\paragraph{The first deviation from uniformity.}
For fixed frames the exact law determines the correction to independent
uniform entries at every matrix. Put $d=a+b$, $\sigma_q(u)=\sum_{j=0}^{u-1}q^{2j}$
(so $\sigma_q(0)=0$), and note that the number of transverse lines is
$T_1=(q^r-q^{t_0}-q^{t_1}+1)/(q-1)$, the lines of $R$ minus those in $R_0$ or
$R_1$.

\begin{theorem}[First correction and its sharp rate; \stU]\label{spn:thm:first-correction}
Fix $q,a,b,H_1,H_0$ and let $n\to\infty$. For $M\in\spnMat_{a,b}(\F_q)$ with
invariants $s,r,t_0,t_1$ put
\[
 K(M)=q^{d-1}T_1-\sigma_q(d-s)+\sigma_q(a-s_1)+\sigma_q(b-s_0),
\]
with the first term read as zero if $d=0$. Then, uniformly in $M$,
\begin{gather*}
 q^{ab}\spnP\bigl(M(g)=M\bigr)=1+q^{-2n}K(M)+O(q^{-4n}),\\
 \spnTV\bigl(M(g),Z\bigr)=\frac{q^{-2n}}{2q^{ab}}\sum_M\abs{K(M)}+O(q^{-4n}),
\end{gather*}
where $Z$ is uniform on $\spnMat_{a,b}(\F_q)$. For every $n$ and all $a,b\ge1$,
\begin{equation}\label{spn:eq:scalar-lower}
 \spnTV\bigl(M(g),Z\bigr)\ge\frac{q-1}{q(q^{2n}-1)} ,
\end{equation}
so the coefficient of $q^{-2n}$ is at least $(q-1)/q$ and the rate $q^{-2n}$
is sharp for every pair of nonempty frames. If $a=0$ or $b=0$ the law is
trivially uniform.
\end{theorem}

\begin{proof}
Put $z=q^{-2n}$ and $C(u,z)=\prod_{j=0}^{u-1}(1-q^{2j}z)$. Extracting powers of
$q$ from \eqref{spn:eq:E-count}, and using
$\binom d2-\binom a2-\binom b2=ab$, the law becomes
\[
 q^{ab}\spnP\bigl(M(g)=M\bigr)=\frac{\sum_{\ell=0}^rT_\ell\,
 q^{d\ell-\ell(\ell+1)/2}z^\ell\,C(d-\ell-s,z)}{C(a-s_1,z)\,C(b-s_0,z)} ;
\]
this holds at every admissible $n$, since a quotient too large to embed
contributes the vanishing factor $j=n$. The denominator has constant term one
and $C(u,z)=1-\sigma_q(u)z+O(z^2)$. Only $\ell=0$ and $\ell=1$ contribute to the
linear coefficient, with $-\sigma_q(d-s)$ and $q^{d-1}T_1$, and inverting the
denominator adds $\sigma_q(a-s_1)+\sigma_q(b-s_0)$. There are finitely many
$M$, so the remainder is uniform, and $\bigl|\abs{x+y}-\abs x\bigr|\le\abs y$
gives the distance after summation.

For \eqref{spn:eq:scalar-lower} take nonzero $u\in\F_q^a$, $v\in\F_q^b$. The
scalar $u^{\spnT}M(g)v$ is the pairing of the fixed nonzero vector $Yu$ with
$gX_0v$, which is uniform on $\F_q^{2n}\setminus\set0$ because $\Sp(V)$ is
transitive on nonzero vectors. So it vanishes with probability
$(q^{2n-1}-1)/(q^{2n}-1)$ and takes each nonzero value with probability
$q^{2n-1}/(q^{2n}-1)$, while under $Z$ it is uniform; the distance between these
two laws is $(q-1)/(q(q^{2n}-1))$, and a function of an observation cannot
increase total variation.
\end{proof}

For $a=b=1$ and $H_0=H_1=0$ one finds $K(0)=1-q$ and $K(M)=1$ for $M\ne0$; the
observation is the scalar itself, so \eqref{spn:eq:scalar-lower} is an
equality and the coefficient $(q-1)/q$ is attained.

Every nontrivial additive character $\psi$ satisfies
\[
 \spnE\psi\bigl(u^{\spnT}M(g)v\bigr)=-\frac1{q^{2n}-1},\qquad
 q^{-ab}\sum_MK(M)\psi(u^{\spnT}Mv)=-1,
\]
so every nonzero rank-one character
test sees the same first-order bias whatever the two forms. The constants in
\cref{spn:thm:first-correction} depend on $q,a,b,H_0,H_1$. In ambient
dimension $2n$, \cref{spn:thm:symplectic} gives only the exponent
$(2n+1-a-b)/2$, but uniformly in the frame sizes, and
\cref{spn:thm:symplectic-exact} shows that the fixed-size rate cannot persist
when $a$ and $b$ grow with $n$. \Cref{spn:thm:complement-boundary,spn:thm:offset-profile}
identify positive limits when the target frame, or both frames, are
nondegenerate and grow, and \cref{spn:thm:padic-lifting} shows that the first
correction persists at every congruence precision over a local ring.

%======================================================================

\subsection{Lagrangian projections of an arbitrary alternating form}\label{spn:sec:lagrangian}
\sources{symplectic_rank_blindness_2026-11-26; grh_sharp_torsion_and_symplectic_blocks; grh_exact_extensions_2026-11-27; grh_extensions_2026-11-27}

In this subsection the ambient space is $\F_q^{2n}$, with
\[
 J_n=\begin{pmatrix}0&I_n\\-I_n&0\end{pmatrix},\qquad
 \Sp_{2n}(\F_q)=\set{g:g^{\spnT}J_ng=J_n},
\]
and $q$ is any prime power. An alternating matrix has zero diagonal and satisfies
$H^{\spnT}=-H$; the zero diagonal is a separate condition in characteristic two. Fix $n\ge1$,
$0\le2s\le n$ and an alternating $H\in\spnMat_n(\F_q)$ of rank $2s$, and put $d=n-2s$. Choose
an injective $X_0\in\spnMat_{2n,n}(\F_q)$ with $X_0^{\spnT}J_nX_0=H$; one exists, since a
symplectic $2s$-space and a $d$-dimensional isotropic subspace of its orthogonal complement
can be embedded. For uniform $g\in\Sp_{2n}(\F_q)$ write
\begin{equation}\label{spn:eq:lagrangian-model}
 gX_0=\begin{pmatrix}A\\B\end{pmatrix},\qquad A,B\in\spnMat_n(\F_q).
\end{equation}
In the notation of \cref{spn:sec:positive-type}, $A$ is the compression of $g$ between the
column space of $X_0$, on which the symplectic form restricts to $H$, and the first $n$
coordinate functionals; under the symplectic identification of $V$ with $V^*$ these functionals
span the Lagrangian $\langle f_1,\dots,f_n\rangle$. With the target frame $Y=(0;I_n)$ one has
$Y^{\spnT}J_ngX_0=-A$, so this is the case $a=b=n$, $H_1=0$, $H_0=H$ of
\cref{spn:thm:complete-law}, whose support condition there says that $\ker A$ is
$H$-isotropic; in this case the sum \eqref{spn:eq:complete-law} has a closed form, which a
direct fibre count gives. For $H=0$ both spaces are Lagrangian, which is the case $c=0$ of
\cref{spn:thm:symplectic-exact}. Let $U_n$ be uniform on
$\spnMat_n(\F_q)$. For a fixed matrix $M$ put
\[
 K=\ker M,\qquad k=\dim K,\qquad t=\dim(K\cap\operatorname{rad}H),\qquad
 C_{n,s}=\prod_{j=s+1}^n(1-q^{-2j}),
\]
and call $K$ \emph{$H$-isotropic} if $x^{\spnT}Hy=0$ for all $x,y\in K$. Empty products equal
one.

\begin{corollary}[Likelihood of every Lagrangian projection; \stU]\label{spn:thm:lagrangian}
Over every finite field,
\begin{equation}\label{spn:eq:lagrangian-density}
 \frac{\spnP(A=M)}{q^{-n^2}}=
 \begin{cases}
 \dfrac{q^{\binom k2}}{C_{n,s}}\displaystyle\prod_{j=k-t+1}^{k}(1-q^{-j}),
 &K\text{ is $H$-isotropic},\\[8pt]
 0,&\text{otherwise}.
 \end{cases}
\end{equation}
\end{corollary}

\begin{proof}
\emph{The frame orbit.} By \cref{spn:lem:gram-frames}, $\Sp_{2n}(\F_q)$ acts transitively on
\[
 \mathcal X_H=\set{X\in\spnMat_{2n,n}(\F_q):\rank X=n,\ X^{\spnT}J_nX=H},
\]
so $gX_0$ is uniform on $\mathcal X_H$ (see also the form-prescribed sampling framework of
\cite{Ishizuka2026Uniform}), and
\begin{equation}\label{spn:eq:frame-orbit}
 \abs{\mathcal X_H}=\spnEmb_n(n,s)=q^{n(3n+1)/2}C_{n,s}.
\end{equation}

\emph{The fibre over $M$.} The Gram condition on $(M;B)$ is
\begin{equation}\label{spn:eq:lagrangian-gram}
 M^{\spnT}B-B^{\spnT}M=H .
\end{equation}
Its restriction to $K\times K$ vanishes, so there is no solution unless $K$ is
$H$-isotropic. Put $r=n-k$ and choose invertible $P,T$ with
$PMT=\operatorname{diag}(I_r,0_k)$; replacing $B$ by $P^{-\spnT}BT$ and $H$ by $T^{\spnT}HT$
preserves \eqref{spn:eq:lagrangian-gram}, injectivity of the frame and the invariants $k,t$. In
these coordinates write $B=(B_{ij})$ and
$H=\left(\begin{smallmatrix}H_{11}&H_{12}\\-H_{12}^{\spnT}&0\end{smallmatrix}\right)$ in blocks
of sizes $r$ and $k$. Equation \eqref{spn:eq:lagrangian-gram} becomes
$B_{11}-B_{11}^{\spnT}=H_{11}$ and $B_{12}=H_{12}$. There are $q^{r(r+1)/2}$ choices of
$B_{11}$, choosing its diagonal and one entry of each off-diagonal pair, which needs no
division by two, and $q^{kr}$ free choices of $B_{21}$. On $K$ the kernel of $H_{12}$ is
$K\cap\operatorname{rad}H$, of dimension $t$, and the frame is injective exactly when
$v\mapsto(H_{12}v,B_{22}v)$ is injective on $K$, that is when $B_{22}$ is injective on that
$t$-dimensional space. This leaves $q^{k(k-t)}\prod_{i=0}^{t-1}(q^k-q^i)
=q^{k^2}\prod_{j=k-t+1}^k(1-q^{-j})$ choices of $B_{22}$. Since
$r(r+1)/2+kr+k^2=n(n+1)/2+\binom k2$, the fibre has
\[
 q^{n(n+1)/2+\binom k2}\prod_{j=k-t+1}^k(1-q^{-j})
\]
elements. Divide by \eqref{spn:eq:frame-orbit} and by $q^{-n^2}$.
\end{proof}

If $H$ is nonsingular then $t=0$ for every $K$, and every solution of
\eqref{spn:eq:lagrangian-gram} is automatically injective, since a kernel vector of $(M;B)$
lies in $\operatorname{rad}H$; every supported likelihood then exceeds one. For degenerate $H$
this can fail: for $q=2$ and $k=t=2$ the numerator in \eqref{spn:eq:lagrangian-density} is
$2(1-2^{-1})(1-2^{-2})=3/4$, while $C_{n,s}>\prod_{j\ge2}(1-2^{-2j})>11/12$ when $s\ge1$, so
this supported likelihood is less than one.

The likelihood converts into finite sums without enumerating the group. Let
$I_s(l)=\prod_{i=0}^{l-1}(q^{2s-i}-q^i)/(q^l-q^i)$ for $0\le l\le s$ and $I_s(l)=0$
otherwise; this is the number of isotropic $l$-subspaces of a nondegenerate alternating
$2s$-space, counted by ordered independent isotropic bases. The number of $H$-isotropic
$k$-subspaces $K$ with $\dim(K\cap\operatorname{rad}H)=t$ is
\begin{equation}\label{spn:eq:kernel-types}
 N_{n,s}(k,t)=\spnqbin{d}{t}{q}\,q^{(d-t)(k-t)}I_s(k-t):
\end{equation}
choose $K\cap\operatorname{rad}H$, then the isotropic image of $K$ in the nondegenerate
quotient by $\operatorname{rad}H$, and then $K$ as the graph of a linear map from that image to
$\operatorname{rad}H/(K\cap\operatorname{rad}H)$, after splitting the quotient. Conditional on
its nullity $k$, the kernel of $U_n$ is uniform among the $k$-subspaces, and
\begin{equation}\label{spn:eq:uniform-nullity}
 u_n(k)=\spnP(\dim\ker U_n=k)=q^{-k^2}\frac{a_n^2}{a_k^2a_{n-k}},\qquad
 a_j=\prod_{i=1}^j(1-q^{-i}),
\end{equation}
by choosing the kernel and then an injective map from its quotient. Put
\[
 w_{n,s}(k,t)=u_n(k)\,N_{n,s}(k,t)\Big/\spnqbin nkq ,
\]
and let $L_{n,s}(k,t)$ be the supported value of \eqref{spn:eq:lagrangian-density}. Summation
of the likelihood and of the overlap $\min\set{\spnP(A=M),\spnP(U_n=M)}$ gives
\begin{gather}
 \spnP(\dim\ker A=k)=\sum_{t=0}^kw_{n,s}(k,t)L_{n,s}(k,t),\label{spn:eq:lagrangian-laws}\\
 \spnTV(A,U_n)=1-\sum_{k=0}^n\sum_{t=0}^kw_{n,s}(k,t)\min\set{1,L_{n,s}(k,t)} .\notag
\end{gather}

\paragraph{The nondegenerate endpoint.}
Let $n=2m$ and let $H$ be nonsingular. Write $C_n=C_{n,m}=\prod_{j=m+1}^n(1-q^{-2j})$,
$a_\infty=\prod_{i\ge1}(1-q^{-i})$, and
\begin{equation}\label{spn:eq:deltaq}
 \delta_q=a_\infty\sum_{k=2}^\infty\frac{q^{-k^2}}{a_k^2}\Bigl(1-q^{-\binom k2}\Bigr).
\end{equation}

\begin{corollary}[Rank tests miss a persistent obstruction; \stU]\label{spn:thm:rank-blind}
For nonsingular $H$ and even $n$,
\begin{equation}\label{spn:eq:rank-blind}
 \spnTV\bigl(\dim\ker A,\dim\ker U_n\bigr)=(u_n(0)+u_n(1))(C_n^{-1}-1)
 =a_n\Bigl[1+\frac{q(1-q^{-n})}{(q-1)^2}\Bigr](C_n^{-1}-1),
\end{equation}
which equals
$\dfrac{a_\infty}{q^2-1}\Bigl(1+\dfrac{q}{(q-1)^2}\Bigr)q^{-n}+O_q(q^{-2n})$ as $n\to\infty$;
the optimal rank-only event is $\set{\rank M\ge n-1}$. In contrast,
\begin{equation}\label{spn:eq:full-blind}
 \spnTV(A,U_n)=\spnP\bigl(H|_{\ker U_n\times\ker U_n}\ne0\bigr)
 =1-\sum_{k\ge0}u_n(k)q^{-\binom k2}R_{n,k}\longrightarrow\delta_q>0 ,
\end{equation}
with $R_{n,k}$ as in \eqref{spn:eq:local-Rnk}, and
$\abs{\spnTV(A,U_n)-\delta_q}\le\bigl(\tfrac4{3(q-1)}+\tfrac1{q^2-1}\bigr)q^{-n}$. The event that the
kernel is not $H$-isotropic never occurs for $A$ and is an optimal distinguishing event at every
finite $n$.
\end{corollary}

\begin{proof}
Over $\F_q$, which is the residue field of the complete discrete valuation ring $\F_q[[t]]$, the
model \eqref{spn:eq:lagrangian-model} with nonsingular $H$ is the first digit $h=1$ of the model of
\cref{spn:sec:local}, the frame being immaterial by \cref{spn:lem:gram-frames}. At $h=1$ the Smith
type is the nullity, so \eqref{spn:eq:rank-blind}, its asymptotic and the optimal rank event are
\cref{spn:cor:local-smith-tv}; \eqref{spn:eq:full-blind} and the optimality of the isotropy test are
\eqref{spn:eq:local-support-tv} and \eqref{spn:eq:local-field}, the sum being \eqref{spn:eq:local-Delta}
with $f_1(k)=q^{-\binom k2}$; and since $\delta_{q,1}=\delta_q$, the limit with its rate is
\cref{spn:thm:local-limit} at $h=1$. The obstruction is visible directly: for $x,y\in\ker A$,
$x^{\spnT}Hy=x^{\spnT}(A^{\spnT}B-B^{\spnT}A)y=0$.
\end{proof}

The limiting nullity law $p_q$ is the classical one \cite{Fulman2015Stein}. For binary
$20\times20$ matrices the two distances are $2.75409684107\cdot10^{-7}$ and
$0.0688049884541$. The following values come from the exact formulas; the constants $\delta_q$
are certified by rational arithmetic, using $a_M(1-q^{-M}/(q-1))\le a_\infty\le a_M$ and the
tail bound $q^{-(K+1)^2}/\bigl(b_-(1-q^{-(2K+3)})\bigr)$ for the terms $k>K$ of
\eqref{spn:eq:deltaq}, where $b_-$ is the displayed lower bound for $a_\infty$; with $M=140$ and
$K=12$ the rational enclosure of $\delta_2$ has width below $4.94\cdot10^{-44}$.
\begin{center}\small
\begin{tabular}{rll}
\toprule
$q$ & $\delta_q$ & coefficient of $q^{-n}$ in \eqref{spn:eq:rank-blind}\\
\midrule
$2$ & $0.0688050065066945\ldots$ & $0.288788095086602$\\
$3$ & $0.0132121494188950\ldots$ & $0.122527579546740$\\
$4$ & $0.00408561989901872\ldots$ & $0.066303614685662$\\
$5$ & $0.00165069311120747\ldots$ & $0.041580699774208$\\
\bottomrule
\end{tabular}
\end{center}
A kernel test is elementary: if $Z$ is a basis matrix of $\ker M$ and $Z^{\spnT}HZ\ne0$, then
$M$ is not a value of \eqref{spn:eq:lagrangian-model}.

The separation does not need a nonsingular form. Put $P_q=a_\infty$,
$B_q=\prod_{j\ge1}(1-q^{-2j})$ and
\[
 \kappa_q=\frac1{2B_q}\Bigl[\frac1{q^2-1}+\frac1{P_q}\Bigl(\frac{2q+3}{q^2-1}
 +\frac{q}{(q-1)^2}\Bigr)\Bigr].
\]

\begin{theorem}[A growing symplectic part suffices; \stU]\label{spn:thm:growing-form}
For an alternating $H$ of rank $2s$ as above,
\[
 \spnTV\bigl(\dim\ker A,\dim\ker U_n\bigr)\le\kappa_q\,q^{-2s},
\]
and $\kappa_2<12$. Along every sequence with $s\to\infty$ and $2s\le n$, the full-matrix
distance $\spnTV(A,U_n)$ tends to the constant $\delta_q$ of \eqref{spn:eq:deltaq}.
\end{theorem}

\begin{proof}
Condition on the nullity $k$ of $U_n$, so that $K$ is a uniform $k$-subspace, and put
$D_{n,k}=\prod_{i=0}^{k-1}(1-q^{i-n})\ge P_q$, $F_{s,k}=\prod_{i=0}^{k-1}(1-q^{2i-2s})$ for
$k\le s$ and $F_{s,k}=0$ for $k>s$. Counting the $H$-isotropic $K$ with
$K\cap\operatorname{rad}H=0$ by \eqref{spn:eq:kernel-types} with $t=0$ gives
\begin{equation}\label{spn:eq:transverse}
 \spnP\bigl(K\text{ isotropic},\ t=0\bigm|k\bigr)=q^{-\binom k2}\,\frac{F_{s,k}}{D_{n,k}} .
\end{equation}
With $h(K)=\mathbf 1_{\set{K\text{ isotropic}}}q^{\binom k2}\prod_{j=k-t+1}^k(1-q^{-j})$ the
likelihood of nullity $k$ is $\spnE(h(K)\mid k)/C_{n,s}$; its part with $t=0$ is
$F_{s,k}/D_{n,k}$, and the rest lies between $0$ and $q^{\binom k2}\spnP(t>0\mid k)$. Counting
lines of $K$ inside the $d$-dimensional radical,
\begin{equation}\label{spn:eq:radical-lines}
 \spnP(t>0\mid k)\le\frac{(q^d-1)(q^k-1)}{(q-1)(q^n-1)}\le\frac{q}{(q-1)^2}(q^k-1)q^{-2s}.
\end{equation}
The elementary product bound gives
\[
 \Bigl|\frac{F_{s,k}}{D_{n,k}}-1\Bigr|\le P_q^{-1}\Bigl[q^{-2s}\frac{q^{2k}-1}{q^2-1}
 +q^{-n}\frac{q^k-1}{q-1}\Bigr],\qquad 1-C_{n,s}\le\frac{q^{-2s}}{q^2-1},\qquad
 C_{n,s}\ge B_q;
\]
for $k>s$ the first bound still holds, since then $q^{-2s}(q^{2k}-1)/(q^2-1)\ge1$. For the
uniform model, $\spnE(q^k-1)\le1$; counting pairs of kernel vectors by the dimension of their
span gives $\spnE q^{2k}=1+(q+1)(1-q^{-n})+(1-q^{-n})(1-q^{1-n})\le q+3$; and
$u_n(k)\le P_q^{-1}q^{-k^2}$ by \eqref{spn:eq:uniform-nullity}, so that
\[
 \spnE\bigl[q^{\binom k2}(q^k-1)\bigr]\le P_q^{-1}\sum_{k\ge1}
 \bigl(q^{-k(k-1)/2}-q^{-k(k+1)/2}\bigr)=P_q^{-1}.
\]
Write the likelihood minus one as $(F_{s,k}/D_{n,k}-1)/C_{n,s}$ plus the rest divided by
$C_{n,s}$ plus $C_{n,s}^{-1}-1$, average against $u_n(k)$ using $q^{-n}\le q^{-2s}$, and divide
by two; this gives the bound with the constant $\kappa_q$. The bounds $P_2>0.2887$ and
$B_2>0.6885$ give $\kappa_2<12$.

For the full law, again conditional on $k$, the likelihood equals
$q^{\binom k2}/C_{n,s}\ge1$ on the event \eqref{spn:eq:transverse} and vanishes off the
isotropic event. The contributions with $t>0$ have uniform probability tending to zero by
\eqref{spn:eq:radical-lines}, and \eqref{spn:eq:transverse} tends to $q^{-\binom k2}$ for fixed
$k$ as $s\to\infty$. So the conditional negative discrepancy tends to $1-q^{-\binom k2}$, and
the uniform nullity tail bound permits summation over $k$, which gives $\delta_q$.
\end{proof}

This rate is not claimed to be sharp for all degeneracies; the nondegenerate formula
\eqref{spn:eq:rank-blind} is exact. At the isotropic input $s=0$ the conditional law given
rank is uniform and rank detects the whole discrepancy, as in
\cref{spn:thm:symplectic-exact}. The likelihood was compared with exhaustive enumeration of all pairs $(A,B)$
over prime fields for $(q,n)=(2,2),(3,2),(2,3)$, which covers $268{,}961$ pairs and $4{,}371$
comparisons of a matrix with a form, including degenerate forms in characteristic two; the
general laws, normalization, the exact nondegenerate distances and every sign in the rank
comparison were evaluated in exact rational arithmetic for $q\in\set{2,3,4,5,8,9}$,
$1\le n\le12$ and $0\le s\le\floor{n/2}$, where for extension fields only the counting
formulas are
tested.\footnote{\raggedright\artifact{symplectic_rank_blindness_2026-11-26/results/exhaustive_check.json},
\artifact{symplectic_rank_blindness_2026-11-26/results/formula_checks.json},
\artifact{symplectic_rank_blindness_2026-11-26/results/certified_constants.json}}
Classical counts of subspaces in symplectic spaces are in
\cite{Wan1993Anzahl,GuoLiWang2013Anzahl}.

\paragraph{Wide blocks.}
Keep the Lagrangian target, but let the nondegenerate source have more columns than the target has
rows. Fix integers $n\ge c\ge1$ with $b=n+c$ even, and put $\ell=n-c$, which is even with
$0\le\ell\le n-1$, and $Q=q^{-1}$. Let $H$ be a nondegenerate alternating $b\times b$ matrix, choose
$X_0\in\spnMat_{2n,b}(\F_q)$ with $X_0^{\spnT}J_nX_0=H$, which exists since $b\le2n$, and for uniform
$g\in\Sp_{2n}(\F_q)$ write $gX_0=(A;B)$ with $A,B\in\spnMat_{n,b}(\F_q)$; let $U$ be uniform on
$\spnMat_{n,b}(\F_q)$. With $a_j$ as in \eqref{spn:eq:uniform-nullity}, $a_0=1$, put
\[
 C'=\prod_{j=\ell/2+1}^{n}(1-Q^{2j}),\quad D'=\prod_{j=b/2+1}^{n}(1-Q^{2j}),\quad
 E'=\prod_{j=n+1}^{n+c}(1-Q^j),
\]
\[
 u'(k)=Q^{k(k-c)}\frac{a_na_b}{a_ka_{k-c}a_{b-k}}\quad(c\le k\le b),\qquad
 R'_{b,k}=\prod_{i=0}^{k-1}\frac{1-Q^{b-2i}}{1-Q^{b-i}}\quad(k\le b/2),
\]
and $R'_{b,k}=0$ for $k>b/2$. A proposed $M$ has kernel $K$ of dimension $k\ge c$.

\begin{theorem}[Exact law of a wide Lagrangian projection; \stU]\label{spn:thm:wide}
Over every finite field,
\[
 \frac{\spnP(A=M)}{q^{-nb}}=\mathbf 1_{\set{K\text{ is $H$-isotropic}}}\frac{q^{\binom k2}}{C'},\qquad
 \spnTV(A,U)=\spnP\bigl(H|_{\ker U\times\ker U}\ne0\bigr).
\]
Moreover $\spnP(\dim\ker U=k)=u'(k)$ and $\spnP(\dim\ker A=k)=u'(k)R'_{b,k}/C'$; full row rank is the
only rank with more probability under $A$ than under $U$, and
\[
 \delta_{\mathrm{rk}}:=\spnTV(\rank A,\rank U)=\frac{a_n}{a_c}\Bigl(\frac1{D'}-E'\Bigr),\qquad
 \Omega:=1-\spnTV(A,U)=\sum_{k=c}^{b/2}u'(k)Q^{\binom k2}R'_{b,k}.
\]
\end{theorem}

\begin{proof}
This is the case of \cref{spn:thm:complete-law} with a Lagrangian target, and the fibre count of
\cref{spn:thm:lagrangian} gives it directly. By \cref{spn:lem:gram-frames} there are
$\spnEmb_n(b,b/2)=q^{2nb-\binom b2}C'$ frames with Gram matrix $H$. For $M$ of rank $r=b-k$, the
condition $M^{\spnT}B-B^{\spnT}M=H$ vanishes on $K\times K$, so isotropy is necessary; with
$PMT=\operatorname{diag}(I_r,0)$ and blocks of sizes $r,n-r$ for rows and $r,k$ for columns it reads
$B_{11}-B_{11}^{\spnT}=H_{11}$, $B_{12}=H_{12}$, $H_{22}=0$, with $q^{r(r+1)/2}$ choices of $B_{11}$ and
the last $n-r$ rows of $B$ free. That is $q^{nb-\binom b2+\binom k2}$ solutions, each an injective frame
because $\operatorname{rad}H=0$, which gives the likelihood. Every supported likelihood exceeds one, so
all deficit lies off the support. Counting a $k$-subspace of $\F_q^b$ and an injection of its quotient
into $\F_q^n$ gives $u'(k)$, and conditional on $k$ the kernel is uniform; choosing an independent
isotropic frame vector by vector leaves $q^{b-i}-q^i$ choices for the next vector, so a uniform
$k$-subspace is isotropic with probability $Q^{\binom k2}R'_{b,k}$, which gives the rank law and
$\Omega$.

For the signs, cancellation gives $R'_{b,c}/C'=(D'E')^{-1}>1$, and
$R'_{b,k+1}/R'_{b,k}=(1-Q^{b-2k})/(1-Q^{b-k})<1$ for $1\le k<b/2$, so it suffices that
$R'_{b,c+1}<C'$. If $c=n$ this is immediate. Otherwise $\ell\ge2$ and
\[
 \frac{R'_{b,c+1}}{C'}=\frac{1-Q^{\ell}}{\prod_{j=n}^{n+c}(1-Q^j)\prod_{j=b/2+1}^{n}(1-Q^{2j})};
\]
the sum of the quantities subtracted from one in the denominator, divided by $Q^\ell$, is at most
$Q^c(1-Q^{c+1})/(1-Q)+Q^{2c+2}/(1-Q^2)$, which is at most $\tfrac34+\tfrac1{12}$ for $c=1$ and
$\tfrac12+\tfrac1{48}$ for $c\ge2$, using only $Q\le\tfrac12$. By $\prod(1-x_i)\ge1-\sum x_i$ the
denominator exceeds the numerator. The excess at full row rank is
$u'(c)((D'E')^{-1}-1)$ with $u'(c)=a_b/a_c$ and $a_b/E'=a_n$.
\end{proof}

\begin{theorem}[Almost singular matrices with almost equal ranks; \stU]\label{spn:thm:wide-separation}
Fix $q$. Throughout the admissible range, $\abs{(q-1)q^n\delta_{\mathrm{rk}}-1}\le3q^{-c}$, so
$\delta_{\mathrm{rk}}=q^{-n}(1+O_q(q^{-c}))/(q-1)$ uniformly in $n\ge c$. With
\[
 \Omega_0=q^{-\binom c2}\frac{a_n}{a_c}\prod_{j=\ell/2+1}^{b/2}(1-q^{-2j}),\qquad
 \varepsilon_q(c)=\frac{a_\infty^{-2}q^{-(2c+1)}}{1-q^{-(2c+4)}},
\]
one has $\Omega_0\le\Omega\le\Omega_0(1+\varepsilon_q(c))$. Consequently, with
$B_{q,\ell}=\prod_{j=\ell/2+1}^\infty(1-q^{-2j})$, uniformly in $\ell\ge0$,
\[
 1-\spnTV(A,U)=B_{q,\ell}\,q^{-\binom c2}\bigl(1+O_q(q^{-c})\bigr),
\]
so as $c\to\infty$ the rank distance tends to zero while the matrix distance tends to one, with no
condition on $n-c$.
\end{theorem}

\begin{proof}
Put $S_E=\sum_{j=n+1}^{n+c}Q^j=Q^n(1-Q^c)/(q-1)$ and $S_D=\sum_{j=b/2+1}^{n}Q^{2j}\le Q^{n+c}/(q^2-1)\le\tfrac1{12}$.
Bonferroni and the product inequalities give $S_E-\tfrac12S_E^2\le1-E'\le S_E$,
$0\le D'^{-1}-1\le S_D/(1-S_D)$ and $1-Q^c/(q-1)\le a_n/a_c\le1$. Dividing the formula for
$\delta_{\mathrm{rk}}$ by $Q^n/(q-1)$ gives at most $1+\tfrac4{11}Q^c$ and at least
$(1-Q^c/(q-1))(1-Q^c-Q^n/(2(q-1)))\ge1-\tfrac52Q^c$. The term $k=c$ of $\Omega$ is $\Omega_0$. For
$k=c+j$, $R'_{b,c+j}\le R'_{b,c}$ and
$u'(c+j)/u'(c)=Q^{j(c+j)}a_ca_n/(a_{c+j}a_ja_{n-j})\le a_\infty^{-2}Q^{j(c+j)}$, so the $j$th term
is at most $a_\infty^{-2}Q^{2cj+(3j^2-j)/2}$ times $\Omega_0$; successive exponents increase by at
least $2c+4$, and a geometric majorant gives the enclosure. Finally $a_n/a_c=1+O_q(Q^c)$, and the
finite even product in $\Omega_0$ differs from $B_{q,\ell}$ by a relative $O_q(Q^{n+c})$, with
$n+c\ge2c$.
\end{proof}

At the boundary $c=n$ the source frame fills the ambient space, $A$ has full row rank surely,
\[
 \delta_{\mathrm{rk}}=1-\prod_{j=n+1}^{2n}(1-q^{-j}),\qquad\Omega=q^{-\binom n2}\prod_{j=1}^n(1-q^{-2j}),
\]
and the second number is also the probability that the $n$-dimensional kernel of a uniform $n\times2n$
matrix is Lagrangian.

Small total variation for one observation does not by itself give the number of observations needed.
Write $\chi^2(P\Vert V)=\sum_x(P(x)-V(x))^2/V(x)$.

\begin{theorem}[Rank-only information and its sharp sample scale; \stU]\label{spn:thm:wide-sample}
Let $P$ and $V$ be the laws of $\rank A$ and $\rank U$. For fixed $q$ and $c\to\infty$, uniformly in
$n\ge c$,
\[
 \chi^2(P\Vert V)=\frac{q^{-(2n-c)}}{q-1}\bigl(1+O_q(q^{-c})\bigr).
\]
For $t$ independent observations of which only the ranks are kept: if $t=o(q^{2n-c})$ the two joint
rank laws have total variation tending to zero, so every rank-only rule has equal-prior error
tending to at least $\tfrac12$; if $t/q^{2n-c}\to\infty$, a threshold on the number of matrices
without full row rank has error tending to zero. For a fixed target error in $(0,\tfrac12)$ the
required number of rank observations is therefore $\Theta_q(q^{2n-c})$. One complete matrix has
optimal equal-prior error
\[
 \frac\Omega2=\frac{B_{q,\ell}}2\,q^{-\binom c2}\bigl(1+O_q(q^{-c})\bigr).
\]
\end{theorem}

\begin{proof}
Index by the row deficiency $j$, so $k=c+j$, and put $u_j=u'(c+j)$, $r_j=R'_{b,c+j}/C'$; then
$\chi^2(P\Vert V)=\sum_{j=0}^nu_j(r_j-1)^2$. One has
$u_1=Q^c(1-Q^n)\prod_{i=c+2}^{n+c}(1-Q^i)/(q-1)=Q^c(1+O_q(Q^c))/(q-1)$, $r_0=1+O_q(Q^n)$, and
$r_1=r_0(1-Q^\ell)/(1-Q^n)$, so $1-r_1=Q^\ell(1+O_q(Q^c))$ (exactly $r_1=0$ when $\ell=0$), and the term
$j=1$ is $Q^{c+2\ell}(1+O_q(Q^c))/(q-1)$, with $c+2\ell=2n-c$; the term $j=0$ is $O_q(Q^{2n})$. For
$j\ge1$, $0\le1-r_j\le Q^{\ell-2(j-1)}/(1-Q^2)$: if $c+j\le b/2$, then
$r_j\ge R'_{b,c+j}\ge\prod_{i<c+j}(1-Q^{b-2i})$, and otherwise $r_j=0$ and $\ell-2(j-1)\le0$. With
$u_j\le a_\infty^{-1}Q^{j(c+j)}$, the terms $j\ge2$ divided by $Q^{c+2\ell}$ are at most
$Q^{(j-1)c+(j-2)^2}/(a_\infty(1-Q^2)^2)$, which sum to $O(Q^c)$.

For $t$ observations, $1+\chi^2(P^{\otimes t}\Vert V^{\otimes t})=(1+\chi^2(P\Vert V))^t$, and
Cauchy--Schwarz gives
\[
 \spnTV(P^{\otimes t},V^{\otimes t})\le\tfrac12\bigl((1+\chi^2(P\Vert V))^t-1\bigr)^{1/2}.
\]
For the test, let $p_U$ and $p_A$ be the probabilities of rank below $n$; by the sign pattern $p_U-p_A=\delta_{\mathrm{rk}}$,
and $0\le p_A\le p_U\le Q^c/(q-1)$. Comparing the observed fraction of rank-deficient matrices with
$(p_A+p_U)/2$, Chebyshev's inequality bounds each error by
$4p_U/(t\delta_{\mathrm{rk}}^2)=O_q(q^{2n-c}/t)$. The complete-matrix test accepts the symplectic model
exactly when the kernel is $H$-isotropic; it never errs under that model, errs with probability
$\Omega$ under the uniform one, and is optimal by \cref{spn:thm:wide}.
\end{proof}

If $c/n\to\alpha\in(0,1)$, the rank distance is asymptotic to $q^{-n}/(q-1)$ and the rank-only sample
scale is $q^{(2-\alpha)n+o(n)}$, while one complete matrix errs with probability
$\tfrac12q^{-\alpha^2n^2/2+o(n^2)}$. Over $\F_2$, for $(n,c)=(40,20)$ one has
$\delta_{\mathrm{rk}}\approx9.095\cdot10^{-13}$, $\Omega\approx6.372\cdot10^{-58}$ and a rank-only sample
scale $2^{60}$; for $(n,c)=(12,12)$, $\delta_{\mathrm{rk}}\approx2.441\cdot10^{-4}$ and
$\Omega\approx9.331\cdot10^{-21}$. Enumeration of $313{,}170$ prescribed-Gram frames over $\F_2$, $\F_3$
and $\F_4$ checks the likelihood on $11{,}295$ matrices, and exact rational arithmetic checks the rank
law, its sign pattern and the formula for $\delta_{\mathrm{rk}}$ for all $288$ admissible triples with
$q\le5$ and $n\le16$.\footnote{\raggedright\artifact{grh_extensions_2026-11-27/code/verify_symplectic.py},
\artifact{grh_extensions_2026-11-27/results/symplectic_validation.json}}

%======================================================================

\subsection{Lagrangian projections at every precision: cokernels conceal the dependence}\label{spn:sec:local}
\sources{grh_exact_extensions_2026-11-27}

Over a finite field the Lagrangian projection of a nondegenerate frame keeps a positive distance from
independent entries while its rank law becomes indistinguishable (\cref{spn:thm:rank-blind}). Over a
complete discrete valuation ring the whole Smith normal form, including every higher invariant
factor and at every precision, has exactly the distinguishing power of the rank modulo the maximal
ideal, whereas the full matrix gains information with every digit, at the sharp rate $q^{-5h}$.

Let $\OO$ be a complete discrete valuation ring with uniformizer $\varpi$ and residue field $\F_q$,
with no restriction on the characteristic, the ramification or the parity of $q$, and put
$R_h=\OO/\varpi^h\OO$ for $h\ge1$. Let $n=2m\ge2$, let $J$ be the standard alternating matrix of size
$2n$, and let $H\in\spnMat_n(\OO)$ be unimodular and alternating ($x^{\spnT}Hx=0$ for all $x$). Fix
$X_0\in\spnMat_{2n,n}(\OO)$ with $X_0^{\spnT}JX_0=H$; for instance $X_0=(I_n;B_0)$ with $B_0$ strictly upper
triangular and $B_0-B_0^{\spnT}=H$. For Haar $g\in\Sp_{2n}(\OO)$ write $gX_0=(A;B)$, let $U\in\spnMat_n(\OO)$
have independent additive Haar entries, and let a subscript $h$ denote reduction modulo $\varpi^h$;
equivalently $A_h$ comes from uniform $g_h\in\Sp_{2n}(R_h)$. Put
\begin{gather}
 a_r=\prod_{j=1}^r(1-q^{-j}),\quad a_\infty=\prod_{j\ge1}(1-q^{-j}),\quad
 u_r(k)=q^{-k^2}\frac{a_r^2}{a_k^2a_{r-k}}\ (0\le k\le r),\notag\\
 C_n=\prod_{j=m+1}^n(1-q^{-2j}),\notag\\
 R_{n,k}=\prod_{i=0}^{k-1}\frac{1-q^{2i-n}}{1-q^{i-n}}\quad(0\le k\le m),\qquad R_{n,k}=0\quad(k>m),
 \label{spn:eq:local-Rnk}
\end{gather}
with $a_0=1$ and $u_r(k)=0$ outside $[0,r]$; $u_r$ is the nullity law of a uniform $r\times r$ matrix
over $\F_q$ \eqref{spn:eq:uniform-nullity}, and $R_{n,0}=R_{n,1}=1$.

\begin{lemma}[Prescribed nondegenerate frames; \stU]\label{spn:lem:local-frames}
$\Sp_{2n}(R_h)$ acts transitively on $\mathcal X_H(R_h)=\set{X\in\spnMat_{2n,n}(R_h):X^{\spnT}JX=H}$,
every such $X$ is a split injection, and $\abs{\mathcal X_H(R_h)}=q^{hD_n}C_n$ with $D_n=(3n^2+n)/2$.
\end{lemma}
\begin{proof}
$H^{-1}X^{\spnT}J$ is a left inverse of $X$, so the image is a free nondegenerate summand with free
nondegenerate complement of rank $n$. Over a local ring a unimodular alternating form has a
symplectic basis (find a unit pairing, normalize, split off the hyperbolic plane, continue), with no
division by two; applied to the complement this gives transitivity with stabilizer $\Sp_n(R_h)$. The
reduction $\Sp_{2t}(R_{r+1})\to\Sp_{2t}(R_r)$ is onto with kernel of order $q^{t(2t+1)}$, as in the
proof of \cref{spn:thm:padic-tangent}, and
$\abs{\Sp_{2t}(\F_q)}=q^{t(2t+1)}\prod_{j\le t}(1-q^{-2j})$; the orbit size
$\abs{\Sp_{2n}(R_h)}/\abs{\Sp_n(R_h)}$ is the count.
\end{proof}

For $M\in\spnMat_n(R_h)$ choose a Smith decomposition $PMT=\operatorname{diag}(\varpi^{a_1},\dots,\varpi^{a_n})$
with $P,T\in\GL_n(R_h)$ and $0\le a_i\le h$, the value $\varpi^h$ standing for zero, and put
\[
 H'=T^{\spnT}HT,\qquad S_h(a)=\sum_{i<j}\min(a_i,a_j)=\sum_{\ell=1}^h\binom{k_\ell(a)}2,\qquad
 k_\ell(a)=\#\set{i:a_i\ge\ell}.
\]

\begin{theorem}[Density at every finite precision; \stU]\label{spn:thm:local-density}
Call $M$ admissible if $H'_{ij}\in\varpi^{\min(a_i,a_j)}R_h$ for all $i<j$; this does not depend on the
Smith decomposition, and it is equivalent to
\[
 H\bigl(K_\ell(M),K_\ell(M)\bigr)\subseteq\varpi^\ell R_h\quad(1\le\ell\le h),\qquad
 K_\ell(M)=\set{x\in R_h^n:Mx\in\varpi^\ell R_h^n}.
\]
The likelihood relative to independent uniform entries is
\[
 \frac{\spnP(A_h=M)}{\spnP(U_h=M)}=\frac{q^{S_h(a)}}{C_n}\mathbf 1_{\set{M\text{ admissible}}}.
\] Every supported
likelihood exceeds one, so with $\mathcal E_h$ the admissible set
\begin{equation}\label{spn:eq:local-support-tv}
 \Delta_{n,h}:=\spnTV(A_h,U_h)=\spnP(U_h\notin\mathcal E_h),
\end{equation}
and inadmissibility is an optimal distinguishing event: the kernel modules must stay isotropic to
the corresponding precision.
\end{theorem}
\begin{proof}
The fibre over $M$ is the set of $B$ with $M^{\spnT}B-B^{\spnT}M=H$. With $B'=P^{-\spnT}BT$ this becomes
$DB'-B'^{\spnT}D=H'$, one equation $\varpi^{a_i}b'_{ij}-\varpi^{a_j}b'_{ji}=H'_{ij}$ for each $i<j$. The
map $(b'_{ij},b'_{ji})\mapsto\varpi^{a_i}b'_{ij}-\varpi^{a_j}b'_{ji}$ has image $\varpi^{\min(a_i,a_j)}R_h$
and kernel of order $q^{h+\min(a_i,a_j)}$, also when an exponent is $h$, and the diagonal entries are
free. So the fibre is empty unless $M$ is admissible, and otherwise has $q^{hn(n+1)/2+S_h(a)}$
elements; divide by $q^{hD_n}C_n$ and by $q^{-hn^2}$. Admissibility is solvability of the equation,
hence intrinsic. In Smith coordinates the $i$th coordinate of $K_\ell(D)$ is
$\varpi^{\max(\ell-a_i,0)}R_h$, and $\min(a_i,a_j)+\max(\ell-a_i,0)+\max(\ell-a_j,0)\ge\ell$ gives one
direction; conversely, for $a=\min(a_i,a_j)\ge1$ both coordinate vectors lie in $K_a(D)$.
\end{proof}

\begin{lemma}[The support inside a Smith orbit; \stU]\label{spn:lem:local-orientation}
For a fixed Smith type $a$ with $k=\#\set{i:a_i\ge1}$, a uniform matrix of that type is admissible
with probability $q^{-S_h(a)}R_{n,k}$.
\end{lemma}
\begin{proof}
A uniform element of the Smith orbit is $PDT$ with independent uniform $P,T\in\GL_n(R_h)$, and $H'$ is
then uniform among unimodular alternating forms, which form one congruence orbit: its strict upper
entries are uniform subject only to nonsingularity of the reduction, with independent uniform higher
digits. In the first digit the $k$-space of positive exponents must be isotropic, which for a
nondegenerate alternating $n$-space has probability
$\prod_{i<k}(q^{n-i}-q^i)/(q^n-q^i)=q^{-\binom k2}R_{n,k}$ (zero for $k>m$), by transitivity of
$\GL_n(\F_q)$ on $k$-spaces. The remaining divisibility conditions have independent total cost
$q^{-(S_h(a)-\binom k2)}$.
\end{proof}

\begin{theorem}[Exact blindness of the entire Smith form; \stU]\label{spn:thm:local-smith}
For every $h\ge1$ and every capped Smith type $a$, with $k=\#\set{i:a_i\ge1}$,
\[
 \spnP(\mathrm{SN}(A_h)=a)=\frac{R_{n,k}}{C_n}\,\spnP(\mathrm{SN}(U_h)=a).
\]
Conditional on the nullity $k\le m$ modulo $\varpi$, the entire Smith forms of $A_h$ and $U_h$ have
the same law, every higher exponent included, for every $h$.
\end{theorem}
\begin{proof}
On a Smith orbit the supported likelihood is the constant $q^{S_h(a)}/C_n$; average it with
\cref{spn:lem:local-orientation}.
\end{proof}

\begin{corollary}[Exact Smith testing distance; \stU]\label{spn:cor:local-smith-tv}
For every $h\ge1$,
\[
 s_n:=\spnTV(\mathrm{SN}(A_h),\mathrm{SN}(U_h))=(u_n(0)+u_n(1))(C_n^{-1}-1)
 =a_n\Bigl[1+\frac{q(1-q^{-n})}{(q-1)^2}\Bigr](C_n^{-1}-1),
\]
an optimal Smith-only event is nullity at most one modulo $\varpi$, and
$s_n=\frac{a_\infty}{q^2-1}\bigl(1+\frac q{(q-1)^2}\bigr)q^{-n}+O_q(q^{-2n})$. At $h=1$,
\begin{equation}\label{spn:eq:local-field}
 \Delta_{n,1}=\spnP\bigl(H|_{\ker U_1\times\ker U_1}\ne0\bigr).
\end{equation}
\end{corollary}
\begin{proof}
By \cref{spn:thm:local-smith} the Smith distance is the distance between the first nullities, with
likelihood $R_{n,k}/C_n$, which exceeds one for $k=0,1$. For $2\le k\le m$ every factor of $R_{n,k}$ is
at most one, and
$1-R_{n,k}\ge1-R_{n,2}=(q^2-q)q^{-n}/(1-q^{1-n})\ge q(q-1)q^{-n}>q^{-n}/(q^2-1)>\sum_{j=m+1}^nq^{-2j}\ge1-C_n$,
so $R_{n,k}<C_n$; for $k>m$ the likelihood is zero. Summing the positive discrepancies, with
$u_n(0)=a_n$ and $u_n(1)=a_nq(1-q^{-n})/(q-1)^2$, gives the formula, and
$C_n^{-1}-1=q^{-n}/(q^2-1)+O_q(q^{-2n})$, $a_n=a_\infty(1+O_q(q^{-n}))$ give the asymptotic. At $h=1$
admissibility is isotropy of $\ker U_1$.
\end{proof}

\begin{proposition}[Infinite precision; \stU]\label{spn:prop:local-infinity}
$A$ is almost surely nonsingular over the fraction field, and relative to independent Haar entries
its density is
\[
 h_\infty(M)=C_n^{-1}q^{\sum_{i<j}\min(a_i,a_j)}\mathbf 1_{\set{H'_{ij}\in\varpi^{\min(a_i,a_j)}\OO\ \forall i<j}}
\]
at every matrix with finite Smith exponents. Its Smith law satisfies \cref{spn:thm:local-smith}, and
the distance between the full cokernel laws of $A$ and $U$ is exactly $s_n$.
\end{proposition}
\begin{proof}
An iid Haar square matrix is almost surely nonsingular: by induction and Fubini, off the null set
where a chosen $(n-1)$-minor vanishes, the determinant is a nonconstant affine function of the
remaining diagonal entry. Once $h$ exceeds the largest exponent the finite densities stabilize to
$h_\infty$. \Cref{spn:lem:local-orientation} holds with Haar changes of basis in $\GL_n(\OO)$, so the
conditional expectation of $h_\infty$ on a Smith type is $R_{n,k}/C_n$, and
$\int h_\infty=C_n^{-1}\sum_ku_n(k)R_{n,k}=1$ by normalization at $h=1$. The finite densities are
nonnegative with integral one, so Scheff\'e's lemma gives convergence in $L^1$; the limit agrees with
the law of $A$ on cylinder sets, hence on all Borel sets, and there is no singular part.
\end{proof}

\begin{corollary}[Every finite cokernel; \stU]\label{spn:cor:local-cokernel}
Let $G$ be a finite $\OO$-module with $k=\dim_{\F_q}(G/\varpi G)\le n$. Then
\[
 \spnP(\spncoker A\cong G)=\frac{R_{n,k}}{C_n}\,\frac{a_n^2}{a_{n-k}\abs{\Aut_\OO(G)}},
\]
which vanishes for $k>m$. Conditional on any $k\le m$ the cokernel has the independent-entry
conditional law, and the optimal cokernel-only test is whether the cokernel is cyclic, the trivial
module included.
\end{corollary}
\begin{proof}
For iid Haar $U$ a fixed lattice $L\subset\OO^n$ with quotient $G$ is the image with probability
$\abs G^{-n}a_n$, and there are $\abs G^n\abs{\Aut_\OO(G)}^{-1}\prod_{i<k}(1-q^{i-n})$ such lattices, by
reduction modulo $\varpi$ and Nakayama's lemma, so
$\spnP(\spncoker U\cong G)=a_n^2/(a_{n-k}\abs{\Aut_\OO(G)})$ \cite{WangStanley2017Smith}. Multiply by
the Smith likelihood.
\end{proof}

As even $n$ grows, this cokernel law tends in total variation to the Cohen--Lenstra law
$a_\infty/\abs{\Aut_\OO(G)}$: the iid first-nullity tail is at most $2q^{-K}$ beyond $K$, and in the
chain of \cref{spn:lem:local-markov} below each step stays positive with probability at most
$1-a_\infty$, so the largest exponent is at least $h$ with probability at most $(1-a_\infty)^h$,
uniformly in $n$; this tightness, the pointwise limit of the iid formula and $s_n\to0$ give the
claim.

\begin{lemma}[Successive iid nullities; \stU]\label{spn:lem:local-markov}
For an iid Haar square matrix with Smith exponents $a_i$, put $k_\ell=\#\set{i:a_i\ge\ell}$. Then $k_1$
has law $u_n$, and $(k_\ell)$ is a Markov chain with $\spnP(k_{\ell+1}=j\mid k_\ell=k)=u_k(j)$; the same
holds through level $h$ for a uniform matrix over $R_h$.
\end{lemma}
\begin{proof}
Given that the reduction has rank $n-k$, lift row and column operations making an upper
$(n-k)$-block invertible and eliminate. Conditional on the other blocks, the Schur complement is a
translate of a uniform lift of its zero residue, so it is uniform in $\varpi\spnMat_k(\OO)$; dividing by
$\varpi$ leaves a fresh iid $k\times k$ Haar matrix, and the Smith form consists of $n-k$ units and
$\varpi$ times the factors of that matrix.
\end{proof}

Put $f_0(k)=1$ and $f_h(k)=q^{-\binom k2}\sum_{j=0}^ku_k(j)f_{h-1}(j)$ for $h\ge1$, so that
$f_h(0)=f_h(1)=1$, $f_1(k)=q^{-\binom k2}$, and $f_h(k)$ is the expectation of
$q^{-\sum_{\ell\le h}\binom{k_\ell}2}$ along the chain started at $k_1=k$.

\begin{theorem}[Exact matrix distance at every precision; \stU]\label{spn:thm:local-precision}
For every $h\ge1$,
\begin{equation}\label{spn:eq:local-Delta}
 \Delta_{n,h}=1-\sum_{k=0}^mu_n(k)R_{n,k}f_h(k).
\end{equation}
With $f_\infty(0)=f_\infty(1)=1$ and, for $k\ge2$ in increasing order,
\[
 f_\infty(k)=\frac{q^{-\binom k2}\sum_{j<k}u_k(j)f_\infty(j)}{1-q^{-(3k^2-k)/2}},
\]
$f_h(k)$ decreases to $f_\infty(k)$, the full $\OO$-valued matrix distance is
$\Delta_{n,\infty}=1-\sum_{k\le m}u_n(k)R_{n,k}f_\infty(k)$, and $f_\infty(2)=(q^4-1)/(q^5-1)$.
\end{theorem}
\begin{proof}
By \eqref{spn:eq:local-support-tv} and \cref{spn:lem:local-orientation},
$1-\Delta_{n,h}=\spnE[q^{-S_h(a)}R_{n,k_1(a)}]$ under iid entries, and the formula
$S_h(a)=\sum_\ell\binom{k_\ell}2$ with \cref{spn:lem:local-markov} gives \eqref{spn:eq:local-Delta}. The
path interpretation shows that $f_h$ decreases; the limit of the recurrence uses $u_k(k)=q^{-k^2}$, and
\cref{spn:prop:local-infinity} passes from cylinder distances to the full distance. For $k=2$ the
states below have value one, so $f_\infty(2)=q^{-1}(1-q^{-4})/(1-q^{-5})$.
\end{proof}

\begin{theorem}[The sharp fifth-power precision law; \stU]\label{spn:thm:local-five}
For every even $n\ge2$ and $h\ge1$, $0\le\Delta_{n,\infty}-\Delta_{n,h}\le q^{-5h}$. For even $n\ge4$
also
\[
 \Delta_{n,\infty}-\Delta_{n,h}\ge u_n(2)R_{n,2}\Bigl(1-\frac{q^4-1}{q^5-1}\Bigr)q^{-5h},
\]
so the exponent five is sharp and the full distance increases strictly with every digit. For $n=2$
the distance is $q^{-4}$ at every precision.
\end{theorem}
\begin{proof}
On the states $k\ge2$ let $Q(k,j)=q^{-\binom k2}u_k(j)$, $2\le j\le k$. The error
$e_h=f_h-f_\infty$ vanishes at $k=0,1$, satisfies $e_h=Qe_{h-1}$ and $0\le e_0\le1$. Every row sum of
$Q$ is at most $q^{-5}$: at $k=2$ it is $q^{-1}q^{-4}$; at $k=3$ counting independent ordered pairs in
the kernel gives $\spnP(\text{nullity}\ge2)\le1/((q^2-1)(q^2-q))$, so the row sum is at most
$1/(q^4(q-1)^2(q+1))\le q^{-5}$; for $k\ge4$ the discount alone is at most $q^{-6}$. So
$0\le e_h\le q^{-5h}$, and averaging in \eqref{spn:eq:local-Delta} with $0\le R_{n,k}\le1$ gives the upper
bound. At state two the recursion is exact, $f_h(2)-f_\infty(2)=(1-f_\infty(2))q^{-5h}$, which gives the
lower bound, positive for $n\ge4$, and the strict increase. For $n=2$ only $k=0,1$ carry weight, and
the forbidden first-digit matrix is the zero matrix.
\end{proof}

After the first nullity has reached two, staying at two for one more digit costs $q^{-4}$ and keeping
admissibility at that digit costs a further $q^{-1}$; higher nullities have smaller diagonal factors
$q^{-(3k^2-k)/2}$.

\begin{theorem}[Simultaneous dimension and precision limits; \stU]\label{spn:thm:local-limit}
Let $p_q(k)=a_\infty q^{-k^2}/a_k^2$ be the limiting iid nullity law \cite{Fulman2015Stein}, and for
$h\in\set{1,2,\dots,\infty}$ put $\delta_{q,h}=\sum_{k\ge0}p_q(k)(1-f_h(k))$. Uniformly in $h$,
\[
 \abs{\Delta_{n,h}-\delta_{q,h}}\le q^{-n}\Bigl(\frac1{(q-1)-q^{-n}}+\frac1{q^2-1}\Bigr)\le
 \Bigl(\frac4{3(q-1)}+\frac1{q^2-1}\Bigr)q^{-n},
\]
for finite $h$ moreover $\abs{\Delta_{n,h}-\delta_{q,\infty}}\le(\frac4{3(q-1)}+\frac1{q^2-1})q^{-n}+q^{-5h}$,
and
\[
 \delta_{q,\infty}-\delta_{q,1}\ge p_q(2)\frac{q-1}{q(q^5-1)}>0 .
\]
\end{theorem}
\begin{proof}
With $B_n=a_n/a_\infty$, $u_n(k)/p_q(k)=B_na_n/a_{n-k}\le B_n$ for $k\le n$, and trivially for $k>n$,
so $\spnTV(u_n,p_q)\le B_n-1\le q^{-n}/((q-1)-q^{-n})$ by $\prod_{j>n}(1-q^{-j})\ge1-q^{-n}/(q-1)$.
Normalization of the first-digit law gives $\sum_ku_n(k)R_{n,k}=C_n$ with $1-C_n\le q^{-n}/(q^2-1)$.
Since $0\le f_h\le1$,
$\abs{\Delta_{n,h}-\delta_{q,h}}=\abs{\sum_k(p_q(k)-u_n(k))f_h(k)+\sum_ku_n(k)(1-R_{n,k})f_h(k)}\le\spnTV(u_n,p_q)+1-C_n$,
and $q^{-n}/(q-1)\le\tfrac14$. The second bound adds $0\le f_h-f_\infty\le q^{-5h}$ from the proof of
\cref{spn:thm:local-five}, and the gap is the term $k=2$:
$p_q(2)(q^{-1}-(q^4-1)/(q^5-1))$.
\end{proof}

Hence every coupled growth of even $n$ and of $h$ has the full-matrix limit $\delta_{q,\infty}$, while
the Smith distance tends to zero uniformly in the precision; at $h=1$, $\delta_{q,1}$ is the constant
$\delta_q$ of \eqref{spn:eq:deltaq}. All quantities are rational for given $q,n,h$, and the limiting
constants are enclosed by the method described after \cref{spn:thm:rank-blind}. For $q=2$ and $n=4$,
\[
 s_4=\frac{7337}{417792},\quad \Delta_{4,1}=\frac{2213}{32768},\quad\Delta_{4,2}=\frac{72391}{1048576},\quad
 \Delta_{4,h}=\frac{35089}{507904}-\frac{1575}{31744}\,2^{-5h}\ (h\ge1),
\]
only the state $k=2$ contributing to the dependence on $h$. Further values, rounded:
\begin{center}\small
\begin{tabular}{rrlll}
\toprule
$q$ & $n$ & Smith, every precision & full matrix, one digit & full matrix, all digits\\
\midrule
$2$ & $4$ & $0.017561370251$ & $0.067535400391$ & $0.069085890247$\\
$2$ & $8$ & $0.0011263070836$ & $0.068730737837$ & $0.070799468377$\\
$2$ & $20$ & $2.7540968411\cdot10^{-7}$ & $0.068804988454$ & $0.070908663442$\\
$3$ & $4$ & $0.0014973991965$ & $0.013120650932$ & $0.013168571188$\\
$3$ & $20$ & $3.5140566616\cdot10^{-11}$ & $0.013212149417$ & $0.01326643453$\\
\midrule
$q$ & & $\delta_{q,1}$ & $\delta_{q,\infty}$ &\\
\midrule
$2$ & & $0.068805006506694520573$ & $0.070908690034008657231$ &\\
$3$ & & $0.013212149418894995143$ & $0.013266434532721205987$ &\\
$4$ & & $0.0040856198990187158379$ & $0.0040896087111604532303$ &\\
$5$ & & $0.0016506931112074708975$ & $0.0016512212975639941216$ &\\
\bottomrule
\end{tabular}
\end{center}
For binary $20\times20$ matrices every test based on the complete cokernel has probability gap at most
$2.7540968411\cdot10^{-7}$, while inspecting the kernel pairings of the matrix gives
$0.070908663442\ldots$; higher precision improves the full-matrix distinction and leaves the
Smith-only one unchanged. Exact enumeration checks the binary fibres of all $65{,}536$ matrices in
$\spnMat_4(\F_2)$, all pairs of $2\times2$ matrices over $\Z/4$ and $\F_3$, the $15$ capped Smith types
over $\Z/4$ in dimension four, and normalization, Smith distances and monotonicity for $70$ pairs
$(q,n)$.\footnote{\raggedright\artifact{grh_exact_extensions_2026-11-27/code/verify_exact.py},
\artifact{grh_exact_extensions_2026-11-27/results/exact_verification.json},
\artifact{grh_exact_extensions_2026-11-27/results/tables_and_certificates.json}}
Classical rank and Smith distributions \cite{Fulman2015Stein,WangStanley2017Smith} are the background;
the symplectic likelihood, its cancellation on every Smith orbit and the precision estimates are
specific to the compression. The model differs from that of \cref{spn:sec:padic}, whose target is a
symplectic summand and whose full distance does not change with the precision.

%======================================================================

\subsection{A hidden source form shared by many wide blocks}\label{spn:sec:hidden}
\sources{grh_hidden_forms_and_filters_2026-11-27}

When the source form of a wide Lagrangian projection is unknown and uniform, one observed block carries no
information about it, but many blocks sharing the same unknown form constrain it jointly. For
blocks with two-dimensional kernels the problem is exactly that of recovering a nondegenerate
alternating form from random isotropic planes, and it has a sharp additive threshold at the number
of free parameters of the form.

Fix a prime power $q$, an even integer $n\ge4$ and $V=\F_q^n$. An alternating form is a bilinear
form $H$ with $H(v,v)=0$ for every $v$, also in characteristic two. Put
\[
 D_n=\binom n2,\qquad d=D_n-1,\qquad \gamma_q=\prod_{j\ge1}(1-q^{-j}),
\]
\[
 c_n=\prod_{j=2}^{n/2}\bigl(1-q^{-(2j-1)}\bigr),\qquad c_q=\prod_{j\ge2}\bigl(1-q^{-(2j-1)}\bigr),
\]
with empty products equal to one. Let $\mathcal H_n$ be the set of nondegenerate alternating forms
on $V$ modulo nonzero scalars and $\mathcal G_n$ the set of two-dimensional subspaces of $V$; a
plane $K$ is $H$-isotropic if $H|_{K\times K}=0$. Then
\begin{equation}\label{spn:eq:hidden-counts}
 \abs{\mathcal H_n}=M_n=q^dc_n,\qquad
 p_n=\frac{q^{n-2}-1}{q^{n-1}-1}<q^{-1},
\end{equation}
where $p_n$ is the fraction of planes that are $H$-isotropic, the same for every nondegenerate $H$.
Under $Q_{n,t}$ an ordered $t$-tuple of planes is independent and uniform on $\mathcal G_n$. Under
$P_{n,t}$ one first chooses $[H]\in\mathcal H_n$ uniformly, once for all the observations, and then
$t$ independent uniform $H$-isotropic planes. For observed planes $K_1,\dots,K_t$ let
\[
 \mathcal N(K_1,\dots,K_t)=\set{G\in\spnAlt(V):G|_{K_i\times K_i}=0\text{ for every }i},
\]
the solution space of $t$ homogeneous linear equations in the $D_n$ entries of $G$ above the
diagonal, and let $Z$ be the number of projective nondegenerate forms in $\mathcal N$; under
$P_{n,t}$ it contains $\langle H\rangle$.

\begin{theorem}[A sharp additive sample threshold; \stU]\label{spn:thm:hidden-threshold}
Fix $q$, let $n\to\infty$ through even integers, and let $w_n\to\infty$ be positive.
\begin{enumerate}
\item If $0\le t_n\le d-w_n$, then $\spnTV(P_{n,t_n},Q_{n,t_n})\to0$, and the optimal probability of
recovering the hidden projective form from the observations tends to zero.
\item If $t_n\ge d+w_n$, then $\spnTV(P_{n,t_n},Q_{n,t_n})\to1$, and, uniformly in the nondegenerate
hidden form $H$, with conditional probability tending to one,
$\mathcal N(K_1,\dots,K_{t_n})=\langle H\rangle$; so Gaussian elimination recovers $[H]$.
\end{enumerate}
\end{theorem}

\begin{theorem}[The bounded window; \stU]\label{spn:thm:hidden-window}
For each fixed integer $a$,
\[
 \chi^2(P_{n,d+a}\Vert Q_{n,d+a})\longrightarrow c_q^{-1}q^a .
\]
Consequently, for each fixed integer $w\ge0$,
\begin{gather*}
 \limsup_{n}\spnTV(P_{n,d-w},Q_{n,d-w})\le\tfrac12c_q^{-1/2}q^{-w/2},\qquad
 \liminf_{n}\spnTV(P_{n,d+w},Q_{n,d+w})\ge1-c_qq^{-w},\\
 \limsup_{n}\spnP(\text{optimal recovery at }d-w)\le c_q^{-1}q^{-w},\\
 \limsup_{n}\sup_H\spnP_H\bigl(\mathcal N\ne\langle H\rangle\text{ at }d+w\bigr)\le\frac{q^{-w}}{q-1}.
\end{gather*}
\end{theorem}

The window needed for any fixed accuracy is therefore bounded independently of $n$. The second
moment limit does not identify the limiting total variation inside the window, nor the law of $Z$.
Both theorems follow from explicit finite bounds. Put
\[
 \rho_q=\frac{2q-1}{q^2}<1,\qquad K_q=4q(q-1),\qquad A_q=\frac q{\gamma_qc_q}.
\]

\begin{proposition}[Finite bounds; \stU]\label{spn:prop:hidden-finite}
The likelihood is
\begin{equation}\label{spn:eq:hidden-likelihood}
 \frac{dP_{n,t}}{dQ_{n,t}}=\frac{Z}{M_np_n^t},
\end{equation}
so in particular $P_{n,1}=Q_{n,1}$. For every $t$ and every integer $0\le R<n$,
\begin{gather}
 \chi^2(P_{n,t}\Vert Q_{n,t})\le e^{K_qtq^{-R}}-1+\frac1{M_np_n^t}
 +A_qq^{-\binom{n-R}2}\Bigl(\frac{\rho_q}{p_n}\Bigr)^t,\label{spn:eq:hidden-chi}\\
 Q_{n,t}(Z>0)\le M_np_n^t,\qquad
 \spnP(\text{optimal recovery})=\frac{Q_{n,t}(Z>0)}{M_np_n^t}\le\frac1{M_np_n^t},
 \label{spn:eq:hidden-bayes}\\
 \sup_H\spnP_H\bigl(\mathcal N\ne\langle H\rangle\bigr)\le\frac{q^d}{q-1}\bigl[p_ne^{K_qq^{-R}}\bigr]^t
 +\frac{q^{D_n-\binom{n-R}2}}{(q-1)\gamma_q}\rho_q^t.\label{spn:eq:hidden-recovery}
\end{gather}
\end{proposition}

\begin{proof}[Counts and the likelihood]
$\GL_n(\F_q)$ acts transitively on the nondegenerate alternating forms with stabilizer the
symplectic group, and $\abs{\GL_n(\F_q)}/\abs{\Sp_n(\F_q)}=q^{D_n}\prod_{j=1}^{n/2}(1-q^{-(2j-1)})$;
dividing by $q-1$ gives $M_n$. For an isotropic plane choose an ordered independent pair $(u,v)$: for
$u\ne0$ the functional $H(u,\cdot)$ is nonzero, its kernel has $q^{n-1}$ elements, and removing
$\langle u\rangle$ leaves $q^{n-1}-q$ choices of $v$ out of $q^n-q$, which gives $p_n$. Each compatible
$[H]$ has conditional likelihood $p_n^{-t}$ relative to $Q_{n,t}$, and averaging gives
\eqref{spn:eq:hidden-likelihood}; hence $\spnE_QZ=M_np_n^t$, which bounds $Q_{n,t}(Z>0)$. The
posterior under $P_{n,t}$ is uniform among the $Z$ compatible forms, so the optimal Bayes success is
$\spnE_P(1/Z)=\sum_{Z>0}Q(\boldsymbol K)\frac{Z}{M_np_n^t}\frac1Z=Q(Z>0)/(M_np_n^t)$. Transitivity of
$\GL(V)$ on planes makes the number of compatible forms constant on $\mathcal G_n$, so one plane is
uniform after averaging the form.
\end{proof}

\begin{lemma}[Common isotropic planes; \stU]\label{spn:lem:hidden-pair}
Fix a nondegenerate alternating $H$ and an alternating $G$, and put
\[
 e(H,G)=\sum_{a\in\F_q}\bigl(q^{\dim\ker(G-aH)}-1\bigr),\qquad r(H,G)=\frac{e(H,G)}{q^n-1},\qquad
 x=q^{n-1}.
\]
The probability $p_2(H,G)$ that a uniform plane is isotropic for both forms satisfies
\begin{gather}
 p_2(H,G)=\frac{e(H,G)(q^{n-1}-q)+(q^n-1-e(H,G))(q^{n-2}-q)}{(q^n-1)(q^n-q)},\notag\\
 \frac{p_2(H,G)}{p_n^2}=1+\frac{(q-1)x\bigl[(x-1)r(H,G)-(q-1)\bigr]}{(x-q)^2}.\label{spn:eq:hidden-explicit}
\end{gather}
If $G$ is not proportional to $H$, then $p_2(H,G)/p_n\le\rho_q$, and with
$s=\min_a\rank(G-aH)$, $p_2(H,G)/p_n^2\le1+K_qq^{-s}$.
\end{lemma}

\begin{proof}
For nonzero $u$ the functionals $H(u,\cdot)$ and $G(u,\cdot)$ are dependent exactly when
$(G-aH)u=0$ for some scalar $a$, which is unique because $H$ is nondegenerate; so there are
$e(H,G)$ such $u$. Their common kernel then has dimension $n-1$, and otherwise $n-2$, and both
kernels contain $u$ because both forms are alternating. Counting the admissible $v$ gives $p_2$, and
substitution gives \eqref{spn:eq:hidden-explicit}. Since $(q^n-1)(q^n-q)p_n=(q^n-1)(q^{n-1}-q)$,
\[
 \frac{p_2(H,G)}{p_n}=r+(1-r)\frac{q^{n-2}-q}{q^{n-1}-q},\qquad r=r(H,G).
\]
If $G$ is not proportional to $H$, each $G-aH$ is a nonzero alternating form, of rank at least two,
so $r\le q(q^{n-2}-1)/(q^n-1)<q^{-1}$, and the display is at most
$q^{-1}+(1-q^{-1})q^{-1}=\rho_q$. For the last bound, $n\ge4$ gives $x\ge q^3$ and
$x(x-1)/(x-q)^2\le2$, while $r\le q\,q^{n-s}/(q^n-1)\le2q\,q^{-s}$; dropping the negative term of
\eqref{spn:eq:hidden-explicit} gives $1+4q(q-1)q^{-s}$.
\end{proof}

\begin{lemma}[Few forms are close to a scalar multiple; \stU]\label{spn:lem:hidden-rank}
The number $B_{n,R}$ of alternating forms on $V$ of rank at most $R$ satisfies
$B_{n,R}\le\gamma_q^{-1}q^{D_n-\binom{n-R}2}$. For a fixed nondegenerate $H$ and a uniform
nondegenerate $G$, $\spnP\set{\min_a\rank(G-aH)\le R}\le A_qq^{-\binom{n-R}2}$.
\end{lemma}

\begin{proof}
A form of rank at most $R$ has a kernel containing some $(n-R)$-subspace $W$; there are at most
$\gamma_q^{-1}q^{R(n-R)}$ such $W$, by the product formula for the Gaussian coefficient, and the
forms vanishing on $W$ descend to the $R$-dimensional quotient, which carries $q^{\binom R2}$
alternating forms; $R(n-R)+\binom R2=D_n-\binom{n-R}2$. For nondegenerate $G$ and $R<n$ the scalar
$a=0$ does not contribute, the union over the other $q-1$ scalars has at most $(q-1)B_{n,R}$
elements, and there are $(q-1)M_n$ nondegenerate $G$; divide, using $c_n\ge c_q$.
\end{proof}

\begin{proof}[Proof of \cref{spn:prop:hidden-finite}]
Squaring \eqref{spn:eq:hidden-likelihood} and averaging under $Q$ gives
\begin{equation}\label{spn:eq:hidden-second}
 1+\chi^2(P_{n,t}\Vert Q_{n,t})=\spnE_{[H],[G]\in\mathcal H_n}\Bigl(\frac{p_2(H,G)}{p_n^2}\Bigr)^t,
\end{equation}
with independent uniform classes. Equal classes have probability $1/M_n$ and contribute exactly
$1/(M_np_n^t)$. Pairs with $\min_a\rank(G-aH)>R$ contribute at most $e^{K_qtq^{-R}}$ by
\cref{spn:lem:hidden-pair} and $(1+z)^t\le e^{tz}$, and the remaining nonproportional pairs at most
$A_qq^{-\binom{n-R}2}(\rho_q/p_n)^t$ by \cref{spn:lem:hidden-pair,spn:lem:hidden-rank}; this is
\eqref{spn:eq:hidden-chi}, and \eqref{spn:eq:hidden-bayes} was proved above.

Fix $H$. A competitor $G\notin\langle H\rangle$ defines the same event on $H$-isotropic planes as
$aG+bH$ for every $a\ne0$ and every $b$, so the distinct events are indexed by the $(q^d-1)/(q-1)$
projective points of $\spnAlt(V)/\langle H\rangle$, each represented by $q(q-1)$ forms. The quantity
$\min_b\rank(G-bH)$ does not depend on the representative; at most $qB_{n,R}$ representatives have
it at most $R$, so at most $B_{n,R}/(q-1)$ points are near, and each survives all $t$ observations
with probability at most $\rho_q^t$. A far point survives with probability at most
$[p_ne^{K_qq^{-R}}]^t$, by \cref{spn:lem:hidden-pair}. The union bound and \cref{spn:lem:hidden-rank}
give \eqref{spn:eq:hidden-recovery}, and when no competitor survives the solution space is exactly
$\langle H\rangle$.
\end{proof}

\begin{proof}[Proof of \cref{spn:thm:hidden-threshold,spn:thm:hidden-window}]
Take $R=\floor{n/5}$, so that $tq^{-R}\to0$ whenever $t=O(n^2)$, and put
$\lambda_q=\log_q(q\rho_q)=\log_q(2-q^{-1})$; then $\lambda_2=\log_2(3/2)<16/25$ and
$\lambda_q<\log_32<16/25$ for $q\ge3$. Since $\binom{n-R}2=(16/25+o(1))D_n$,
$p_n=q^{-1}(1+O_q(q^{2-n}))$ and $M_n=q^d(c_q+o(1))$, the near terms of \eqref{spn:eq:hidden-chi}
and \eqref{spn:eq:hidden-recovery} are at most $q^{-(16/25-\lambda_q+o(1))D_n}=o(1)$, uniformly for
$0\le t\le d+O(1)$.

At $t=d+a$ the proportional term tends to $c_q^{-1}q^a$, which gives the upper bound for the
chi-square limit. For the lower bound, \eqref{spn:eq:hidden-explicit} with $r\ge0$ gives
$p_2/p_n^2\ge1-\delta_n$ with $\delta_n=(q-1)^2q^{n-1}/(q^{n-1}-q)^2=O_q(q^{1-n})$, so the unequal
pairs contribute at least $(1-M_n^{-1})(1-\delta_n)^t\to1$ to \eqref{spn:eq:hidden-second}; their
contribution is also at most $1+o(1)$ by the near and far estimates. This proves the chi-square
limit. Cauchy--Schwarz, $\spnTV(P,Q)\le\frac12\chi^2(P\Vert Q)^{1/2}$, and \eqref{spn:eq:hidden-bayes}
give the two bounds below the threshold. The event $\set{Z>0}$ has planted probability one, so
$\spnTV(P,Q)\ge1-Q(Z>0)\ge1-M_np_n^t$, which tends to $1-c_qq^{-w}$ at $t=d+w$; and the far term of
\eqref{spn:eq:hidden-recovery} tends to $q^{-w}/(q-1)$ there. This proves
\cref{spn:thm:hidden-window}.

For diverging gaps, additional observations cannot decrease the total variation or the optimal
recovery probability, nor increase the probability that $\dim\mathcal N>1$. For instance, for each
fixed $w$ the experiment at $t_n\le d-w_n$ is eventually a projection of the one at $d-w$; let
$n\to\infty$ and then $w\to\infty$. The other assertions are obtained in the same way.
\end{proof}

\paragraph{The connection with wide blocks.}
In \cref{spn:thm:wide} take the ambient dimension $2(n-2)$, a target with $n-2$ rows and a
nondegenerate source of dimension $n$, so that $c=2$ there.

\begin{proposition}[Conditioned wide blocks with a shared unknown form; \stU]\label{spn:prop:hidden-wide}
Conditional on full row rank, the kernel of $A$ is uniform among the $H$-isotropic planes and, given
its kernel, $A$ is uniform among the full-row-rank matrices with that kernel; for the
independent-entry block $U$ conditioned on full row rank, the kernel is uniform on $\mathcal G_n$
and, given it, $U$ is uniform in the same set. The full-row-rank probability of $A$ does not depend
on $H$. Consequently $t$ blocks sharing one uniform unknown $[H]$, each conditioned on full row
rank, against $t$ independent conditioned uniform blocks, is the experiment of
\cref{spn:thm:hidden-threshold}, and passing to the kernels loses no information for testing or for
recovery.
\end{proposition}

\begin{proof}
By \cref{spn:thm:wide}, the likelihood of a full-row-rank $M$ is $q/C'$ when $\ker M$ is
$H$-isotropic and zero otherwise, a constant on the supported full-row-rank matrices. For each plane
$K$ the full-row-rank matrices with kernel $K$ are in bijection with the isomorphisms
$V/K\to\F_q^{n-2}$, so there are $\abs{\GL_{n-2}(\F_q)}$ of them for every $K$. This gives both
conditional descriptions, and since the number of $H$-isotropic planes is $p_n\abs{\mathcal G_n}$ for
every nondegenerate $H$, the full-row-rank probability does not depend on $H$; conditioning each
block therefore leaves the uniform prior on $[H]$ unchanged.
\end{proof}

This settles, for blocks with two-dimensional kernels and full row rank, the problem of
observations with an unknown shared source form. No conclusion is drawn for unconditioned blocks.

\begin{remark}[Recovery by elimination]\label{spn:rem:hidden-recovery}
Represent each observed plane by an ordered basis $(u_i,v_i)$ and form the row
$z_i=(u_{i,r}v_{i,s}-u_{i,s}v_{i,r})_{r<s}\in\F_q^{D_n}$; a change of basis multiplies the row by a
nonzero scalar, so the null space of the stacked $t\times D_n$ matrix is intrinsic and consists of
the coefficient vectors of $\mathcal N$. If the null space is one-dimensional, put a generator above
the diagonal of an alternating matrix and check that it is nonsingular. This takes
$O(tD_n^2+D_n^3)$ field operations, hence $O(n^6)$ at the threshold. Under the planted law a
one-dimensional null space automatically has a nonsingular generator; under the null law an output
implies $Z>0$, whose probability is bounded by \eqref{spn:eq:hidden-bayes}. The same procedure is
therefore a polynomial-time test above the threshold.
\end{remark}

For $n=4$ independent enumeration of reduced row spaces gives $35$ planes, $28$ projective
nondegenerate forms and $15$ isotropic planes per form over $\F_2$, and $130$, $234$ and $40$ over
$\F_3$; the largest conditional overlap $p_2/p_n$ between nonproportional forms is $3/5$ and $2/5$,
against the bound $\rho_q=3/4$ and $5/9$. For $(q,n)=(2,4)$, where $d=5$, exact transition counts
over all $2356$ reachable spans of exterior coordinates give the following distances
$\spnTV(P_{4,t},Q_{4,t})$ and probabilities of recovery by elimination, rounded to five places:
\begin{center}\small
\begin{tabular}{rcccccccc}
\toprule
$t$ & $1$ & $2$ & $4$ & $5$ & $6$ & $8$ & $10$ & $12$\\
\midrule
distance & $0$ & $0.11429$ & $0.37280$ & $0.66308$ & $0.84036$ & $0.96869$ & $0.99417$ & $0.99893$\\
elimination & $0$ & $0$ & $0$ & $0.23893$ & $0.49557$ & $0.81123$ & $0.93361$ & $0.97679$\\
\bottomrule
\end{tabular}
\end{center}
These finite values are not limits; statistical distinction, and uniqueness among nonsingular
forms, can occur before the linear equations determine a unique form. For $q=2$ and $n=200$, where
$d=19899$, the finite bounds certify by exact rational comparison that at $t=19889$ the distance is
below $0.0171$ and the optimal recovery probability below $0.001165$, and that at $t=19909$
elimination fails with probability below $0.000977$ while its false-positive probability under the
null law is below
$0.00082$.\footnote{\raggedright\artifact{grh_hidden_forms_and_filters_2026-11-27/data/independent_hidden_forms.json},
\artifact{grh_hidden_forms_and_filters_2026-11-27/data/finite_bound_checks.json},
\artifact{grh_hidden_forms_and_filters_2026-11-27/data/recovery_demo.json}}
The pair-incidence counts are classical finite geometry, close to the eigenvector counts for
symplectic Grassmann codes \cite{CardinaliGiuzzi2016Minimum}, and the Pl\"ucker description of the
symplectic Grassmannian is treated in \cite{Blok2003Generating}.

%======================================================================

\subsection{Nondegenerate frames: residual forms and rank blindness}\label{spn:sec:nondegenerate}
\sources{exact_symplectic_compressions_2026-11-26; symplectic_compression_2026-11-26_2; nondegenerate_symplectic_blocks_2026-11-27; uniform_symplectic_blocks; rectangular_symplectic_blocks; grh_uniform_extensions_2026-11-27; grh-two-open-questions}

When the target frame is nondegenerate, the observed matrix fixes the component of the source
frame along the target, and what remains is one alternating Gram matrix in the orthogonal
complement. This reduction gives exact likelihoods and limiting distances in regimes where the
target frame, or both frames, grow with the ambient dimension: for a fixed source, and for two
nondegenerate frames at every offset from complementary dimension, uniformly in the aspect
ratio. For complementary frames the rank, the residual defect and their joint law have sharp
first-order expansions, uniform in the shape, while the full matrix stays a fixed distance from
uniform. The notation is that of
\cref{spn:sec:complete-law}: $V$ has dimension $2n$, $Y$ is an $a$-frame with Gram matrix $H_1$,
$X_0$ is a $b$-frame with Gram matrix $H_0$, $M(g)=Y^{\spnT}JgX_0$, and $Z$ is uniform on
$\spnMat_{a,b}(\F_q)$. Throughout $H_1$ is nonsingular, so $a$ is even, and $2c=2n-a$ is the
dimension of $(\spnim Y)^\perp$. By \cref{spn:lem:gram-frames} and the invariance of uniform $g$,
the law of $M(g)$ depends on the frames only through $H_1$ and $H_0$, and replacing $H_1$ by a
congruent form multiplies $M(g)$ by a fixed invertible matrix, which changes neither ranks nor
distances from $Z$. Write $J_v$ for the standard alternating matrix of size $2v$,
$\spnAlt_b(\F_q)$ for the additive group of alternating $b\times b$ matrices, of order
$q^{D_b}$ with $D_b=\binom b2$, $U_{\spnAlt}$ for a uniform element of it, and put
\[
 \gamma_r=\prod_{i=1}^r(1-q^{-i}),\qquad
 \eta_r=\prod_{i=1}^r(1-q^{-2i}),\qquad
 \theta_r=\prod_{i=1}^r(1-q^{-(2i-1)}),
\]
\[
 \Pi_{n,b}=\prod_{j=n-b/2+1}^{n}(1-q^{-2j}),
\]
with empty products equal to one and infinite products $\gamma_\infty=\eta_\infty\theta_\infty$;
thus $\gamma_\infty=P_q$ and $\eta_\infty=B_q$ in the notation of \cref{spn:sec:lagrangian}.
The elementary inequality $\prod_i(1-x_i)\ge1-\sum_ix_i$ gives
$\eta_\infty\ge(q^2-2)/(q^2-1)$.

\paragraph{Random alternating Gram matrices.}
For $v\ge0$ and $b\ge1$ let $T_v$ be uniform on $\spnMat_{2v,b}(\F_q)$ and
$G_v=T_v^{\spnT}J_vT_v\in\spnAlt_b(\F_q)$; $G_0=0$.

\begin{lemma}[Random alternating Gram matrices; \stU]\label{spn:lem:alt-gram}
Fix a nontrivial additive character $\psi$ of $\F_q$, and index the characters of
$\spnAlt_b(\F_q)$ by $W\in\spnAlt_b(\F_q)$ through $G\mapsto\psi(\sum_{i<j}W_{ij}G_{ij})$.
\begin{enumerate}
\item The character indexed by $W$ has expectation $q^{-v\rank W}$ under $G_v$.
\item If $b\ge2$, then $\abs{q^{D_b}\spnP(G_v=G)-1}\le(q^{D_b}-1)q^{-2v}$ for every $G$.
\item If $2v\ge b$, then
\[
 \spnTV(G_v,U_{\spnAlt})\le\frac{1}{2\sqrt{\eta_\infty}}
 \Bigl(\sum_{t\ge1}q^{-2(2v-b)t-2t^2-t}\Bigr)^{1/2}
 \le\frac{q^{-(2v-b)-3/2}}{2\sqrt{\eta_\infty\bigl(1-q^{-2(2v-b)-7}\bigr)}} .
\]
\end{enumerate}
If $Q$ is a nonsingular alternating matrix of size $2v$ and $T$ is uniform on
$\spnMat_{2v,b}(\F_q)$, then $T^{\spnT}QT$ has the law of $G_v$. All statements hold in
characteristic two.
\end{lemma}

\begin{proof}
Group the rows of $T_v$ into $v$ independent uniform pairs $(x_i,y_i)$, so that
$G_v=\sum_i(x_iy_i^{\spnT}-y_ix_i^{\spnT})$; pairing one summand with $W$ gives
$\psi(x_i^{\spnT}Wy_i)$, with no division by two. For fixed $x_i$ the average over $y_i$ is one
if $Wx_i=0$ and zero otherwise, so the average over the pair is $q^{-\rank W}$, and
independence gives the first claim. Every nonzero alternating matrix has rank at least two, so
Fourier inversion gives the second. For the third, the number of alternating $b\times b$
matrices of rank $2t$ is
\begin{equation}\label{spn:eq:alt-rank-count}
 A_b(t)=\spnqbin{b}{2t}{q}\frac{\abs{\GL_{2t}(\F_q)}}{\abs{\Sp_{2t}(\F_q)}}
 =q^{t(2b-2t-1)}\frac{\gamma_b}{\gamma_{b-2t}\eta_t}\le\eta_\infty^{-1}q^{t(2b-2t-1)},
\end{equation}
by choosing the radical and then a nondegenerate form on the quotient. Plancherel and
Cauchy--Schwarz give
$\spnTV(G_v,U_{\spnAlt})\le\frac12(\sum_{W\ne0}q^{-2v\rank W})^{1/2}$, and inserting
\eqref{spn:eq:alt-rank-count} gives the first bound. With $d=2v-b\ge0$ the first summand is
$q^{-2d-3}$ and every successive ratio is at most $q^{-2d-7}$, which gives the second. Finally
$Q=P^{\spnT}J_vP$ for an invertible $P$, and $PT$ is uniform.
\end{proof}

For $0\le r\le\floor{b/2}$ the number of $T\in\spnMat_{2c,b}(\F_q)$ with $T^{\spnT}J_cT$ equal to
a prescribed form of rank $2r$ is
\[
 Z_c(b,r)=\sum_{k=0}^{b-2r}\spnqbin{b-2r}{k}{q}\spnEmb_c(b-k,r),
\]
by choosing the kernel inside the radical and embedding the quotient
(\cref{spn:lem:gram-frames}). The \emph{complement discrepancy} is therefore the finite sum
\begin{equation}\label{spn:eq:complement-discrepancy}
 \Delta_{q,b,c}=\spnTV(G_c,U_{\spnAlt})
 =\frac12\sum_{r=0}^{\floor{b/2}}A_b(r)\,
 \Bigl|\frac{Z_c(b,r)}{q^{2cb}}-q^{-D_b}\Bigr| .
\end{equation}

\begin{proposition}[Residual Gram likelihood; \stU]\label{spn:prop:residual}
Let $H_0$ also be nonsingular, so that $b$ is even, and put
\begin{equation}\label{spn:eq:residual-form}
 S(M)=H_0+M^{\spnT}H_1^{-1}M\in\spnAlt_b(\F_q).
\end{equation}
For every $M\in\spnMat_{a,b}(\F_q)$,
\[
 q^{ab}\spnP\bigl(M(g)=M\bigr)=\frac{q^{D_b}\,\spnP\bigl(G_c=S(M)\bigr)}{\Pi_{n,b}} ,
\]
and $M$ occurs with positive probability exactly when $\rank S(M)\le2c$.
\end{proposition}

\begin{proof}
Decompose $V=\spnim Y\oplus(\spnim Y)^\perp$ and choose a basis of the second summand with Gram
matrix $J_c$. A $b$-frame $X$ satisfies $Y^{\spnT}JX=M$ exactly when its component in
$\spnim Y$ is $YH_1^{-1}M$, whose Gram matrix is $-M^{\spnT}H_1^{-1}M$. Writing the other
component as $T\in\spnMat_{2c,b}(\F_q)$, the condition $X^{\spnT}JX=H_0$ becomes
$T^{\spnT}J_cT=S(M)$, and $X$ is automatically injective because $H_0$ is nonsingular. By
\cref{spn:lem:gram-frames}, $gX_0$ is uniform among the
$\spnEmb_n(b,b/2)=q^{2nb-D_b}\Pi_{n,b}$ frames with Gram matrix $H_0$, while the admissible
$T$ number $q^{2cb}\spnP(G_c=S(M))$; since $a+2c=2n$, the ratio is the claim. A form of rank
$2r$ is a Gram matrix in a symplectic $2c$-space exactly when $r\le c$, which gives the support.
\end{proof}

The residual form is the Schur complement recorded in the first special case after
\cref{spn:thm:complete-law}: $\rank\Gamma_M=a+\rank S(M)$.

\paragraph{A fixed source and a growing nondegenerate target.}

\begin{theorem}[Complement-Gram boundary law; \stU]\label{spn:thm:complement-boundary}
Fix $q$, $b\ge2$, $c\ge0$ and an alternating $H_0$ of size $b$ and rank $2h_0$, and let
$a=2s$ with $s\to\infty$, so that $n=s+c$ and frames exist once $n\ge b-h_0$. Then, uniformly
in $H_0$,
\[
 \spnTV\bigl(M(g),Z\bigr)=\Delta_{q,b,c}+O_{q,b,c}(q^{-2s}),\qquad
 \spnTV\bigl(\rank M(g),\rank Z\bigr)=O_{q,b,c}(q^{-2s}).
\]
Explicitly, put $\eps_v=(q^{D_b}-1)q^{-2v}$ and $L_b=(q^b-1)/(q-1)$; if $\eps_n<1$, then
\[
 \Bigl|\spnTV\bigl(M(g),Z\bigr)-\Delta_{q,b,c}\Bigr|
 \le\frac{q^{D_b}L_bq^{-2n}}{1-\eps_n}+\tfrac12\eps_n+\eps_s\Delta_{q,b,c},
\]
\[
 \spnTV\bigl(\rank M(g),\rank Z\bigr)\le\Bigl(1+\frac{q^{D_b}}{\eta_\infty}\Bigr)L_bq^{-2s},
\]
and the first term may be omitted when $H_0$ is nonsingular. The limit does not depend on
$H_0$, not even on its rank, and
\[
 \tfrac12q^{-2c}\le\Delta_{q,b,c}\le\min\set{1,\tfrac12\sqrt{q^{D_b}-1}\,q^{-2c}} ,
\]
so $\Delta_{q,b,c}>0$ for every fixed $c$ and $\Delta_{q,b,c}=\Theta_{q,b}(q^{-2c})$ as
$c\to\infty$.
\end{theorem}

A hidden symplectic complement of fixed dimension therefore leaves a persistent constraint on
the whole observed matrix, carried by its alternating Gram matrix, although the rank laws
converge.

\begin{proof}
Take $H_1=J_s$ and $Y$ a symplectic basis of a nondegenerate $2s$-space $U$, so that
$M(g)=J_sA$ where $A\in\spnMat_{2s,b}(\F_q)$ is the block of $U$-coordinates of $gX_0$; it
suffices to treat $A$. Let $X=(A;T)$ be an unrestricted uniform $2n\times b$ matrix, with
$T$ its block of coordinates in $U^\perp$. Its Gram matrix is $A^{\spnT}J_sA+T^{\spnT}J_cT$.
Let $E$ be the event that it equals $H_0$, and put $F_c(G)=q^{D_b}\spnP(G_c=G)$ and
$R_n=q^{D_b}\spnP(G_n=H_0)$. Conditional on $E$, the law $Q_s$ of $A$ has likelihood
$F_c(H_0-A^{\spnT}J_sA)/R_n$ relative to the uniform law, and $\spnE F_c(H_0-G_s)=R_n$.

By \cref{spn:lem:alt-gram}, $\abs{R_n-1}\le\eps_n$. Put $f(G)=\abs{F_c(H_0-G)-1}$. Translation
and negation preserve uniform measure on $\spnAlt_b$, so $\spnE f(U_{\spnAlt})=2\Delta_{q,b,c}$,
and the pointwise bound of \cref{spn:lem:alt-gram} together with $f\ge0$ gives
$\abs{\spnE f(G_s)-2\Delta_{q,b,c}}\le2\eps_s\Delta_{q,b,c}$. Replacing $F_c/R_n$ by $F_c$
inside the total variation costs at most
$\frac12\abs{R_n^{-1}-1}\spnE F_c(H_0-G_s)=\frac12\abs{R_n-1}$. Hence
$\abs{\spnTV(Q_s,Z)-\Delta_{q,b,c}}\le\frac12\eps_n+\eps_s\Delta_{q,b,c}$.

An unrestricted matrix kills a given line with probability $q^{-2n}$, and
$\spnP(E)\ge q^{-D_b}(1-\eps_n)$, so $X$ fails to be injective given $E$ with probability at
most $q^{D_b}L_bq^{-2n}/(1-\eps_n)$. Given $E$ and injectivity, $X$ is uniform among the frames
with Gram matrix $H_0$, which is the law of $gX_0$. Conditioning changes the joint law by at
most the discarded probability and projection cannot increase total variation, which proves the
explicit full-matrix bound. When $H_0$ is nonsingular the Gram condition forces injectivity.

For the ranks, $\spnEmb_n(b,h_0)=q^{2nb-D_b}C$ with
$C=\prod_{j=n-b+h_0+1}^{n}(1-q^{-2j})\ge\eta_\infty$. A uniform $2s\times b$ matrix fails to
have full column rank with probability at most $L_bq^{-2s}$, so dividing by the probability
$q^{-D_b}C$ that an unrestricted matrix is a frame with Gram matrix $H_0$ gives
$\spnP(\rank A<b)\le q^{D_b}L_bq^{-2s}/\eta_\infty$. Comparing both rank laws with the point
mass at $b$ gives the rank bound.

Finally, a character indexed by a rank-two alternating matrix has expectation $q^{-2c}$ under
$G_c$ and zero under $U_{\spnAlt}$; expectations of a function of modulus at most one differ by
at most twice the total variation, which gives the lower bound. Parseval and Cauchy--Schwarz give
$\Delta_{q,b,c}\le\frac12(\sum_{W\ne0}q^{-2c\rank W})^{1/2}\le\frac12\sqrt{q^{D_b}-1}\,q^{-2c}$.
\end{proof}

\begin{corollary}[Sources of dimension two and three; \stU]\label{spn:cor:small-source}
For $b\in\set{2,3}$, $\Delta_{q,b,c}=(1-q^{-D_b})q^{-2c}$.
\end{corollary}

\begin{proof}
Every nonzero $W\in\spnAlt_b$ has rank two, so by \cref{spn:lem:alt-gram} the law of $G_c$ is
$q^{-2c}\delta_0+(1-q^{-2c})U_{\spnAlt}$, whose distance from $U_{\spnAlt}$ is the displayed
quantity.
\end{proof}

At $c=0$ the hidden Gram matrix is zero and $\Delta_{q,b,0}=1-q^{-D_b}$. Over $\F_2$, exact
rational evaluation of \eqref{spn:eq:complement-discrepancy} gives:
\begin{center}\small
\begin{tabular}{ccccc}
\toprule
$b$ & $c=0$ & $c=1$ & $c=2$ & $c=3$\\
\midrule
$2$ & $1/2$ & $1/8$ & $1/32$ & $1/128$\\
$3$ & $7/8$ & $7/32$ & $7/128$ & $7/512$\\
$4$ & $63/64$ & $7/16$ & $133/1024$ & $553/16384$\\
\bottomrule
\end{tabular}
\end{center}
For a symplectic plane as source the exact finite-$s$ laws, including the optimal rank-only
test, are in \cref{spn:prop:plane-rank}.

\paragraph{Nondegenerate frames at every offset.}
Now let both $H_1$ and $H_0$ be nonsingular, so that $a$ and $b$ are even, put $h=b/2$, and call
\[
 e=c-h=n-\tfrac{a+b}{2}
\]
the \emph{offset}: $2e$ is the ambient dimension left over after both frames are counted, and
$e=0$ exactly when $a+b=2n$. Frames exist exactly when $c\ge0$, that is $e\ge-h$. Put
$F_v(G)=q^{D_b}\spnP(G_v=G)$ for $G\in\spnAlt_b(\F_q)$. Since $T_vP$ is uniform for invertible
$P$, $F_v$ is constant on congruence classes, hence a function of $\rank G$, and by
\cref{spn:prop:residual} the likelihood of $M$ is $F_c(S(M))/\Pi_{n,b}$.

\begin{lemma}[Monotone Gram densities; \stU]\label{spn:lem:gram-monotone}
For all $v\ge0$ and $b\ge1$, $F_v(G)$ is a nonincreasing function of $\rank G$.
\end{lemma}

\begin{proof}
Passing from $G_v$ to $G_{v+1}$ adds $u\wedge w=uw^{\spnT}-wu^{\spnT}$ for independent uniform
$u,w\in\F_q^b$. This increment is symmetric in law under negation, so
$F_{v+1}(G)=\spnE F_v(G+u\wedge w)$: the update of densities is the transition operator of the
walk $G\mapsto G+u\wedge w$. By congruence invariance the half-rank of the walk is a Markov chain,
and one step changes the rank by at most two. Let $G$ have rank $2i$, column space $L$, and put
$f=b-2i$. The rank increases exactly when $u,w$ are independent modulo $L$: then a complement of
$L$ containing $u,w$ exhibits the two extra ranks in a separate block, and otherwise the new
column space lies in a space of dimension at most $2i+1$. Thus the upward probability is
$p_i^+=(1-q^{-f})(1-q^{1-f})$, which vanishes for $f\le1$. A decrease from rank $2i+2$ requires
both vectors in the column space $L'$ of the current matrix: for a nonzero increment, equality
in $\rank G'\le\rank(G'+u\wedge w)+\rank(u\wedge w)$ forces the column space of $u\wedge w$
into $L'$. On $L'$ the block $B'$ of $G'$ is invertible, and the determinant of the updated block
divided by $\det B'$ is $(1+w^{\spnT}B'^{-1}u)^2$, also in characteristic two; given $u\ne0$,
singularity prescribes one value of a nonzero linear functional of $w$. Hence, for $f\ge2$,
\[
 p_{i+1}^-=q^{-2(f-2)-1}\bigl(1-q^{-2(i+1)}\bigr)\le q^{1-f}\le1-p_i^+ .
\]
So for adjacent starting states the probability of crossing their common boundary is ordered,
at every other threshold the ordering is automatic, and the transition operator preserves
nonincreasing functions of the half-rank. The initial density $F_0$, supported at rank zero,
is nonincreasing; iterate.
\end{proof}

\begin{theorem}[The exact test at every offset; \stU]\label{spn:thm:offset-test}
Let $H_0,H_1$ be nonsingular with offset $e$.
\begin{enumerate}
\item If $e\ge0$, the likelihood $q^{ab}\spnP(M(g)=M)$ is less than one exactly when $S(M)$ is
nonsingular, where it equals
\[
 \alpha_e=\frac{\eta_{b/2+e}\,\eta_{a/2+e}}{\eta_e\,\eta_{(a+b)/2+e}} .
\]
Hence $\set{S(M)\text{ singular}}$ is an exactly optimal distinguishing event and
$\spnTV(M(g),Z)=(1-\alpha_e)\,\spnP\bigl(S(Z)\text{ nonsingular}\bigr)$.
\item If $e=-r<0$, then $M$ has positive probability exactly when $\dim\ker S(M)\ge2r$, its
likelihood is then greater than one, and
$\spnTV(M(g),Z)=\spnP\bigl(\dim\ker S(Z)<2r\bigr)$.
\end{enumerate}
\end{theorem}

\begin{proof}
Splitting off symplectic pairs, a nonsingular $b$-form has
$\prod_{i=0}^{h-1}(q^{2(v-i)}-1)q^{2(v-i)-1}$ realizations as $T^{\spnT}J_vT$ when $v\ge h$, so
$F_v(G)=\eta_v/\eta_{v-h}$ at nonsingular $G$; in particular $\Pi_{n,b}=\eta_n/\eta_{n-h}$.

Let $e\ge0$, so $c=h+e$ and $n-h=a/2+e$. At a nonsingular residual the likelihood is
$(\eta_c/\eta_{c-h})(\eta_{n-h}/\eta_n)=\alpha_e$, which is less than one because each of the $h$
factors of $\eta_{h+e}/\eta_e$ has smaller index than the corresponding factor of
$\eta_{h+a/2+e}/\eta_{a/2+e}$. At nullity two, after $h-1$ symplectic pairs have been embedded,
the remaining two vectors lie in a symplectic space of dimension $f=2e+2$ and must pair to zero;
there are $q^f+(q^f-1)q^{f-1}$ such pairs against $(q^f-1)q^{f-1}$ with a prescribed nonzero
pairing, so the density there is
\[
 \frac{\eta_{h+e}}{\eta_e}\Bigl(1+\frac{q}{q^{2e+2}-1}\Bigr).
\]
With $u=q^{-2e}\le1$ the product in front is at least $1-u/(q^2-1)$, and
$(1-u/(q^2-1))(1+qu/(q^2-u))-1$ has the sign of $q^3-q^2-q-(q-1)u\ge q^3-q^2-2q+1>0$. Division
by $\Pi_{n,b}\le1$ only increases the likelihood, and \cref{spn:lem:gram-monotone} extends the
strict inequality to every larger nullity. So all deficient mass lies on the nonsingular
residuals, where the likelihood is the constant $\alpha_e$.

Let $e=-r<0$, so $c=h-r$ and residual ranks are at most $2c$ (\cref{spn:prop:residual}). At rank
exactly $2c$ a realizing $T$ must have kernel equal to the radical and induce a symplectic
isomorphism on the quotient, so $F_c(G)=q^{D_b-2cb}\abs{\Sp_{2c}(\F_q)}=q^{r(2r-1)}\eta_c$, and
$q^{r(2r-1)}\eta_c\ge q\bigl(1-1/(q^2-1)\bigr)>1$. By \cref{spn:lem:gram-monotone} the likelihood
exceeds one on the whole support, and every other matrix has likelihood zero. The distance is the
uniform mass of the complement of the support.
\end{proof}

For complementary frames ($e=0$) this is the sign statement of \cref{spn:thm:balanced} below,
which in addition gives the likelihood at every corank. To evaluate either distance exactly, note
that the independent row pairs of $Z$ add independent uniform wedges to $H_0$, so the half-rank of
$S(Z)$ is the chain of the proof of \cref{spn:lem:gram-monotone} started at $h$, with
$p_i^+=(1-q^{-f})(1-q^{1-f})$ and $p_i^-=q^{-2f-1}(1-q^{-2i})$ at state $i$, $f=b-2i$, run for $a/2$
steps; this needs $O(ab)$ rational operations.

Let $\lambda_b(j)$ be the probability that a uniform element of $\spnAlt_b(\F_q)$ has nullity $j$.
Choosing the radical and a nonsingular form on the quotient gives
\begin{equation}\label{spn:eq:finite-alt}
 \lambda_b(2k)=q^{-k(2k-1)}\frac{\gamma_b}{\gamma_{2k}\,\eta_{h-k}}\qquad(0\le k\le h),
\end{equation}
and $\lambda_b$ vanishes at odd arguments; in particular $\lambda_b(0)=\theta_h$.

\begin{lemma}[Corank mixing with its first term; \stU]\label{spn:lem:corank-mix}
Let $a,b$ be positive and even and $H_0,H_1$ nonsingular. Then
\[
 \spnTV\bigl(\dim\ker S(Z),\dim\ker U_{\spnAlt}\bigr)\le\tfrac12q^{-a}\sqrt{\theta_h^{-1}-1}.
\]
More precisely, with
\[
 H_b(j)=\frac{(q+1)q^b-q^{b+j}-q}{q(q^2-1)(q^{b-1}-1)},
\]
one has
$\sum_j\bigl|\spnP(\dim\ker S(Z)=j)-\lambda_b(j)\bigl(1+q^{-a}H_b(j)\bigr)\bigr|
\le q^{-2a}\sqrt{\theta_h^{-1}-1}$.
\end{lemma}

\begin{proof}
Let $\nu$ be the uniform law on the nonsingular orbit, of uniform mass $\theta_h$, with density
$f=\mathbf 1_{\mathrm{nonsing}}/\theta_h$. By \cref{spn:lem:alt-gram}, $Z^{\spnT}H_1^{-1}Z$ has the
law of $G_{a/2}$, which is congruence invariant, so replacing $H_0$ by an independent $B_0\sim\nu$
does not change the corank law of the translate, and the law of $B_0+G_{a/2}$ is congruence
invariant. Its density is the convolution of $f$ with the law of $G_{a/2}$, whose Fourier transform
at $W$ is $q^{-(a/2)\rank W}$ (\cref{spn:lem:alt-gram}). So convolution multiplies every
nonconstant coefficient by at most $q^{-a}$ and every coefficient of rank at least four by at most
$q^{-2a}$. By Parseval the nonconstant coefficients of $f$ have squared sum $\theta_h^{-1}-1$, and
Cauchy--Schwarz bounds the mean absolute error by its root mean square. This proves the first bound,
and the second up to evaluation of the rank-two terms.

Every rank-two $W$ is $u\wedge w$ for independent $u,w$, in exactly $\abs{\SL_2(\F_q)}=q(q^2-1)$
ways, and its character at $B$ is $\psi(u^{\spnT}Bw)$. At $B$ of nullity $j$, the sum of
$\psi(u^{\spnT}Bw)$ over all ordered pairs is $q^{b+j}$, and the $(q+1)q^b-q$ dependent pairs
each contribute one, so the sum of the rank-two characters at $B$ is
$K(B)=(q^{b+j}-(q+1)q^b+q)/(q(q^2-1))$. By transitivity the rank-two coefficients of $f$ are all
equal, to $K$ at a nonsingular matrix divided by $K(0)$, that is $-1/(q^{b-1}-1)$; $K$ is real.
Multiplication by $q^{-a}$ gives exactly $q^{-a}H_b(j)$, and summing the absolute error over the
corank classes proves the second bound.
\end{proof}

\begin{theorem}[Uniform limiting profile; \stU]\label{spn:thm:offset-profile}
Fix $q$. For all positive even $a,b$, all admissible offsets $e$ and all nonsingular $H_0,H_1$,
\[
 \bigl|\spnTV\bigl(M(g),Z\bigr)-\mathcal D_q(e)\bigr|\le C_q\bigl(q^{-a}+q^{-b}\bigr),\qquad
 \mathcal D_q(e)=
 \begin{cases}
 \theta_\infty\bigl(1-\eta_\infty/\eta_e\bigr),&e\ge0,\\[2pt]
 \displaystyle\sum_{j=0}^{-e-1}\frac{\theta_\infty q^{-j(2j-1)}}{\gamma_{2j}},&e<0,
 \end{cases}
\]
with $C_q$ depending only on $q$. In particular, when $a,b\to\infty$ the distance tends to
$\mathcal D_q(e)$ at every fixed offset, to $\Delta_q=\theta_\infty(1-\eta_\infty)$ at $e=0$ in
every aspect ratio, and $\mathcal D_q(e)\to0$ as $e\to+\infty$ and $\mathcal D_q(e)\to1$ as
$e\to-\infty$.
\end{theorem}

\begin{proof}
For $e\ge0$, \cref{spn:lem:corank-mix} gives
$\spnP(S(Z)\text{ nonsingular})=\theta_h+O_q(q^{-a})$, and the product tails give
$\theta_h=\theta_\infty+O_q(q^{-b})$ and, using $\abs{\eta_j-\eta_\infty}\le q^{-2j}/(q^2-1)$ and
the positive lower bound on all denominators,
$\alpha_e=\eta_\infty/\eta_e+O_q(q^{-a}+q^{-b})$ uniformly in $e\ge0$; insert these in
\cref{spn:thm:offset-test}. For $e<0$, let $\lambda_q(2j)=\theta_\infty q^{-j(2j-1)}/\gamma_{2j}$.
The finite probabilities \eqref{spn:eq:finite-alt} are at most $C_qq^{-j(2j-1)}$, so
$\sum\lambda_q=1$ by dominated convergence, and the product tails with
$\abs{\gamma_b-\gamma_\infty}\le q^{-b}/(q-1)$ give
$\abs{\lambda_b(2j)-\lambda_q(2j)}\le C_qq^{-b}(1+q^{2j})\lambda_q(2j)$, also for $j>h$; the right
side has total $O_q(q^{-b})$. Apply \cref{spn:lem:corank-mix} and part~(2) of
\cref{spn:thm:offset-test} to the event $\dim\ker S(Z)<2r$.
\end{proof}

\begin{corollary}[First correction off the boundary; \stU]\label{spn:cor:offset-first}
Let $e\ge0$, $A_e=\eta_\infty/\eta_e$, $u=q^{-2e}$ and $w=q^{-a}+q^{-b}$. Uniformly in all such
dimensions,
\[
 \spnTV\bigl(M(g),Z\bigr)=\mathcal D_q(e)+\frac{\theta_\infty}{q^2-1}\bigl[q(1-A_e)-uA_e\bigr]w
 +O_q(uw^2).
\]
If $e,a,b\to\infty$ at any relative rates, then
$\spnTV(M(g),Z)\sim\theta_\infty q^{-2e}/(q^2-1)$.
\end{corollary}

\begin{proof}
$H_b(0)=(q^b-1)/((q^2-1)(q^{b-1}-1))=q/(q^2-1)+O_q(q^{-b})$, so \cref{spn:lem:corank-mix} and the
first-order expansion of $\theta_h$ give
$\spnP(S(Z)\text{ nonsingular})=\theta_\infty[1+qw/(q^2-1)]+O_q(w^2)$. Likewise
$\alpha_e=A_e[1+uw/(q^2-1)]+O_q(uw^2)$: the three product tails in $\alpha_e$ have first-order
sizes $uq^{-b}/(q^2-1)$, $uq^{-a}/(q^2-1)$ and $O_q(uq^{-a-b})$, and
$\prod(1-t_i)=1-\sum t_i+O((\sum t_i)^2)$ with uniformly positive denominators. Since
$1-A_e=O_q(u)$, substitution in \cref{spn:thm:offset-test} proves the expansion, and
$1-A_e=u/(q^2-1)+O_q(u^2)$ gives the last assertion.
\end{proof}

Over $\F_2$ the limiting distances at $e=-1,0,1$ are about $0.41942$, $0.13063$ and $0.03437$, so a
change of two in the ambient dimension survives the limit in which both observed dimensions grow.
Rounded to six places, \cref{spn:thm:offset-profile} gives:
\begin{center}\small
\begin{tabular}{lrrrrrr}
\toprule
 & $e=-2$ & $e=-1$ & $e=0$ & $e=1$ & $e=2$ & $e=3$\\
\midrule
$q=2$ & $0.978652$ & $0.419422$ & $0.130634$ & $0.034372$ & $0.008702$ & $0.002182$\\
$q=3$ & $0.998445$ & $0.639005$ & $0.078878$ & $0.008863$ & $0.000986$ & $0.000110$\\
$q=5$ & $0.999933$ & $0.793335$ & $0.033003$ & $0.001322$ & $0.000053$ & $0.000002$\\
\bottomrule
\end{tabular}
\end{center}
At $e=-1$ the limit is $\theta_\infty$. Twenty-seven parameter choices over $\F_2$ and $\F_3$ were
checked by independently counting row pairs with a prescribed alternating Gram
matrix.\footnote{\artifact{grh-two-open-questions/symplectic/verification.json}}

\paragraph{Complementary nondegenerate frames.}
Now let $e=0$, that is $a+b=2n$, so that $(\spnim Y)^\perp$ has the dimension $b=2c$ of the source.
Neither frame is Lagrangian. Write $h=b/2$ and let $2\ell$ be the corank of $S(M)$.

\begin{theorem}[The residual singularity test is optimal; \stU]\label{spn:thm:balanced}
For every $M\in\spnMat_{a,b}(\F_q)$,
\[
 q^{ab}\spnP\bigl(M(g)=M\bigr)=\frac{\eta_h}{\Pi_{n,b}}
 \sum_{d=\ell}^{2\ell}\spnqbin{2\ell}{d}{q}\frac{q^{-d(d+1)/2}}{\eta_{d-\ell}} .
\]
This likelihood is less than one when $S(M)$ is nonsingular and greater than one when $S(M)$ is
singular. Hence $\set{S(M)\text{ singular}}$ is an exactly optimal distinguishing event, and
\[
 \spnTV\bigl(M(g),Z\bigr)=\Bigl(1-\frac{\eta_h}{\Pi_{n,b}}\Bigr)
 \spnP\bigl(S(Z)\text{ is nonsingular}\bigr).
\]
\end{theorem}

\begin{proof}
By \cref{spn:prop:residual} with $c=h$ we need the number of $T\in\spnMat_b(\F_q)$ with
$T^{\spnT}J_hT=S(M)$. The kernel of such a $T$ is a $d$-subspace of the $2\ell$-dimensional
radical of $S(M)$, and the quotient, of dimension $b-d$ and rank $b-2\ell$, embeds in a
symplectic $b$-space exactly when $\ell\le d\le2\ell$. By \eqref{spn:eq:E-count},
$\spnEmb_h(b-d,h-\ell)=q^{b(b+1)/2-d(d+1)/2}\eta_h/\eta_{d-\ell}$, and summation over the
$\spnqbin{2\ell}dq$ kernels and division by $q^{b^2}$ gives the formula.

If $\ell=0$ the sum is one, and $\eta_h/\Pi_{n,b}<1$ because each of the $h$ factors of
$\Pi_{n,b}$ has a larger index than the corresponding factor of $\eta_h$, as $a>0$. For $\ell\ge1$
use $\eta_h/\Pi_{n,b}\ge\eta_\infty\ge(q^2-2)/(q^2-1)$. For $\ell=1$ the sum is
$(q+1)q^{-1}+q^{-3}/(1-q^{-2})=1+q/(q^2-1)$, and
\[
 \frac{q^2-2}{q^2-1}\Bigl(1+\frac{q}{q^2-1}\Bigr)-1=\frac{q^3-q^2-2q+1}{(q^2-1)^2}>0 .
\]
For $\ell\ge2$ the term $d=\ell$ alone gives
$\eta_\infty\spnqbin{2\ell}{\ell}{q}q^{-\ell(\ell+1)/2}\ge\eta_\infty q^{\ell(\ell-1)/2}\ge
q\eta_\infty>1$, since $\spnqbin{2\ell}\ell q\ge q^{\ell^2}$. So all excess mass lies on the
singular residuals and all deficient mass on the nonsingular ones, where the likelihood is the
constant $\eta_h/\Pi_{n,b}$; this gives the distance.
\end{proof}

For $q=3$ and $a=b=2$ the likelihoods are $99/80$ and $9/10$, as in \cref{spn:thm:two-plane},
while the distance between the two rank laws is only $11/270$ against $19/270$ for the matrices:
matrices of determinants one and two have full rank but lie on opposite sides of the threshold.

\begin{corollary}[An explicit lower bound for every complementary pair; \stU]\label{spn:cor:balanced-lower}
For all positive even $a,b$ with $a+b=2n$,
\[
 \spnTV\bigl(M(g),Z\bigr)\ge\Bigl(1-\frac{\eta_h}{\Pi_{n,b}}\Bigr)
 \Bigl(1-\frac{q(1-q^{-a})(1-q^{-b})}{q^2-1}\Bigr).
\]
\end{corollary}

\begin{proof}
For a nonzero $x\in\F_q^b$ change coordinates so that $x$ is the first basis vector, and let
$z_1,\dots,z_b$ be the independent uniform columns of $Z$. If $z_1=0$, then $S(Z)x$ is the
nonzero first column of $H_0$. If $z_1\ne0$, the $i$th entry of $S(Z)x$ for $i\ge2$ is a
prescribed constant plus the nonzero linear functional $z_i\mapsto z_i^{\spnT}H_1^{-1}z_1$, and
the first entry vanishes identically. Hence $\spnP(x\in\operatorname{rad}S(Z))=(1-q^{-a})q^{-(b-1)}$,
and summing over $x$ gives
$\spnE\bigl(q^{\dim\operatorname{rad}S(Z)}-1\bigr)=q(1-q^{-a})(1-q^{-b})$. A singular
alternating matrix of even size has nullity at least two, so $S(Z)$ is singular with probability
at most this expectation divided by $q^2-1$. Apply \cref{spn:thm:balanced}.
\end{proof}

\paragraph{Complementary frames, uniformly in the aspect ratio.}
For complementary frames the rank law, the residual-defect law and the full-matrix law have
first-order expansions whose errors are uniform in the shape. Let $a,b\ge2$ be even, write
$J_a$ and $J_b$ for nonsingular alternating matrices of sizes $a$ and $b$, let $g$ be uniform in the
symplectic group of $V=(\F_q^a,J_a)\perp(\F_q^b,J_b)$, and write the image under $g$ of the
coordinate frame of the second summand as
\begin{equation}\label{spn:eq:block-frame}
 \begin{pmatrix}A\\B\end{pmatrix},\qquad A\in\spnMat_{a,b}(\F_q),\quad B\in\spnMat_b(\F_q),\qquad
 A^{\spnT}J_aA+B^{\spnT}J_bB=J_b .
\end{equation}
With $Y$ and $X_0$ the coordinate frames of the two summands, $H_1=J_a$, $H_0=J_b$, $M(g)=J_aA$ and
$S(J_aA)=\Lambda(A)$, where
\[
 \Lambda(M)=J_b-M^{\spnT}J_aM ;
\]
so this is the case $e=0$ of the preceding paragraphs, $\spnTV(M(g),Z)=\spnTV(A,U)$, and
$\eta_h/\Pi_{n,b}=\alpha_{a,b}=\eta_{a/2}\eta_{b/2}/\eta_{(a+b)/2}$. For $a=b$, with all position
coordinates first, the first summand would instead be Lagrangian, as in \cref{spn:sec:lagrangian}.
Let $U$ be uniform on $\spnMat_{a,b}(\F_q)$ and $T$ uniform on $\spnAlt_b(\F_q)$, and put
\[
 K_A=\dim\ker A,\quad K_U=\dim\ker U,\quad R_U=\dim\ker\Lambda(U),\quad R_T=\dim\ker T,
\]
and $d=a-b$, $\tau=q^{-b}$, $\sigma=q^{-a}$.
Rank and nullity have the same distances. The limiting nullity laws of independent-entry
rectangular and alternating matrices, and the signed sequences describing the corrections, are
\begin{gather}
 \pi_d(k)=\frac{\gamma_\infty q^{-k(k+d)}}{\gamma_k\gamma_{k+d}},\qquad
 \lambda_q(2k)=\frac{\theta_\infty q^{-k(2k-1)}}{\gamma_{2k}},\qquad \lambda_q(2k+1)=0,
 \label{spn:eq:limit-laws}\\
 \xi_d(k)=\frac{1+q^{-d}-q^k}{q-1}\,\pi_d(k),\qquad
 \zeta_q(k)=\frac{q+1-q^k}{q^2-1}\,\lambda_q(k).\label{spn:eq:signed-laws}
\end{gather}
Both signed sequences have total mass zero and their only positive mass is at zero, so
\begin{equation}\label{spn:eq:signed-norms}
 \tfrac12\norm{\xi_d}_1=\xi_d(0)=\frac{q^{-d}\gamma_\infty}{(q-1)\gamma_d},\qquad
 \tfrac12\norm{\zeta_q}_1=\zeta_q(0)=\frac{q\theta_\infty}{q^2-1};
\end{equation}
see \eqref{spn:eq:signed-moments} below. Finally put
\[
 c_q=\frac{\theta_\infty}{q^2-1}\bigl(q-(q+1)\eta_\infty\bigr)<0 .
\]

\begin{theorem}[Uniform rank and defect laws; \stU]\label{spn:thm:uniform-marginals}
There is $C_q<\infty$, depending only on $q$, such that for all even $a\ge b\ge2$
\begin{gather*}
 \norm{\mathcal L(K_A)-\mathcal L(K_U)-\tau\xi_d}_1\le C_q\,q^{-a-b},\\
 \norm{\mathcal L(R_U)-\mathcal L(R_T)-\sigma\zeta_q}_1\le C_q\,q^{-a-b}.
\end{gather*}
Consequently
\begin{gather*}
 \spnTV(\rank A,\rank U)=\frac{\gamma_\infty}{(q-1)\gamma_{a-b}}\,q^{-a}+O_q(q^{-a-b}),\\
 \spnTV(R_U,R_T)=\frac{q\theta_\infty}{q^2-1}\,q^{-a}+O_q(q^{-a-b}),
\end{gather*}
and as $b\to\infty$ the test for full column rank is asymptotically optimal among all tests that
use only the rank.
\end{theorem}

The coefficient of $q^{-a}$ in the rank distance lies between $\gamma_\infty/(q-1)$ and $1/(q-1)$,
so the relative error is $O_q(q^{-b})$ even when $a-b$ grows arbitrarily fast; a remainder
$O_q(q^{-2b})$ would not have this property. For $a=b$ the correction $\tau\xi_0$ has
$\xi_0(k)=(2-q^k)\pi_0(k)/(q-1)$. Thus the rank of the block is exponentially close in law to that
of a uniform matrix, at the sharp rate $q^{-\max(a,b)}$, and the nullity of the quadratic
expression $\Lambda(U)$, whose entries are dependent, follows the alternating law to first order.

\begin{remark}[An explicit rank bound for arbitrary frames]\label{spn:rem:rank-union}
In the setting of \cref{spn:thm:balanced}, with arbitrary complementary nondegenerate frames
$Y$ and $X_0$, at every size
\[
 \spnTV\bigl(\rank M(g),\rank Z\bigr)\le\frac{q^b-1}{q-1}\Bigl(q^{-a}+\frac{q^b-1}{q^{a+b}-1}\Bigr)
 =O_q(q^{b-a}),
\]
and for the coordinate frames this reads $\spnTV(\rank A,\rank U)$. Indeed, the domain has
$(q^b-1)/(q-1)$ lines. A fixed nonzero vector lies in the kernel of $Z$ with probability $q^{-a}$.
Its image under $gX_0$ is uniform among the nonzero vectors of $V$, and it is killed by
$Y^{\spnT}J$ exactly when it lies in $(\spnim Y)^\perp$, of dimension $b$, which has probability
$(q^b-1)/(q^{a+b}-1)$. Union bounds over lines and comparison of both rank laws with the point
mass at full column rank give the bound. In particular the rank of $M(g)$ carries no information
in the limit $a-b\to\infty$, for every choice of frames.
\end{remark}

\begin{theorem}[Uniform full-matrix law; \stU]\label{spn:thm:uniform-full}
For all positive even $a,b$:
\begin{enumerate}
\item $\spnTV(A,U)=(1-\alpha_{a,b})\spnP\bigl(\Lambda(U)\text{ nonsingular}\bigr)
=(\alpha_{a,b}^{-1}-1)\spnP(B\text{ invertible})$;
\item if $a\ge b$, then $\spnTV(A,U)=(1-\eta_{b/2})\theta_{b/2}+c_q\,q^{-a}+O_q(q^{-a-b})$;
\item $\spnTV(A,U)=\Delta_q+c_q(q^{-a}+q^{-b})+O_q\bigl((q^{-a}+q^{-b})^2\bigr)$, without
ordering $a$ and $b$.
\end{enumerate}
In particular $\spnTV(A,U)\to\Delta_q>0$ whenever $\min(a,b)\to\infty$, in every aspect ratio,
while every rank statistic loses its power. For instance
$\Delta_2=0.1306343467085051764310563857\ldots$.
\end{theorem}

The finite-product baseline in (2) matters for very tall blocks: its error is $O_q(q^{-a-b})$,
whereas replacing it by $\Delta_q$ introduces a term of order $q^{-2b}$. With equal prior
probabilities the best probability of identifying the model from the whole matrix is
$\frac12+\frac12\spnTV(A,U)$, attained by declaring the symplectic model when
$\det(J_b-M^{\spnT}J_aM)=0$; it tends to $0.5653171734\ldots$ over $\F_2$. The best success
probability of any procedure that uses only the rank, randomized or not, is
$\frac12+\gamma_\infty q^{-a}/(2(q-1)\gamma_{a-b})+O_q(q^{-a-b})$ for $a\ge b$.

\begin{remark}[A fixed source]\label{spn:rem:fixed-source}
If $b$ is fixed and $a\to\infty$, (1) gives $\spnTV(A,U)\to\theta_{b/2}(1-\eta_{b/2})$, since
$\alpha_{a,b}\to\eta_{b/2}$ and $\Lambda(U)$ becomes uniform among alternating matrices
(\cref{spn:lem:alt-gram}). Both dimensions must therefore grow for the universal constant
$\Delta_q$. Comparing with \cref{spn:thm:complement-boundary} at $c=b/2$ identifies the complement
discrepancy $\Delta_{q,b,b/2}=\theta_{b/2}(1-\eta_{b/2})$, in agreement with the values $1/8$ and
$133/1024$ of the table there. By \cite[Theorem~4.9]{DeBoeckVandeVoorde2024Anzahl}, which counts
nondegenerate subspaces disjoint from a given one, the probability in (1) also has the positive-term
form
\[
 \spnP\bigl(\Lambda(U)\text{ nonsingular}\bigr)
 =\sum_{m=0}^{h}\theta_{h-m}\,\frac{\eta_h}{\eta_m\eta_{h-m}}\,q^{-(a+1)m},
\]
since $\Lambda(A)=B^{\spnT}J_bB$ is nonsingular exactly when the image of the source frame is
disjoint from the first summand. Its terms $m=0,1$ reproduce the first-order expansion of
\eqref{spn:eq:defect-zero} below.
\end{remark}

\begin{theorem}[Joint law of ordinary and residual nullities; \stU]\label{spn:thm:joint-defect}
Let $P_K$ and $P_R$ be the finite laws of $K_U$ and $R_U$. Uniformly for even $a\ge b\ge2$,
\[
 \norm{\mathcal L(K_U,R_U)-P_K\otimes P_R+\tau(q^2-1)\,\xi_d\otimes\zeta_q}_1\le C_q\,q^{-a-b},
\]
and consequently
\[
 \spnTV\bigl(\mathcal L(K_U,R_U),P_K\otimes P_R\bigr)
 =\frac{2q\theta_\infty\gamma_\infty}{(q-1)\gamma_{a-b}}\,q^{-a}+O_q(q^{-a-b}).
\]
\end{theorem}

Centering at the true finite marginals measures the dependence itself rather than a mixture of
dependence and marginal convergence. The negative sign of the leading tensor at $(0,0)$ says that
$U$ and $\Lambda(U)$ are simultaneously of full rank slightly less often than independence
predicts.

\begin{corollary}[Two bits carry the leading dependence; \stU]\label{spn:cor:xor}
Let $E=\set{K_U=0}$, $F=\set{R_U=0}$ and $X=E\mathbin{\triangle}F$. Then
\begin{gather*}
 \operatorname{Cov}(\mathbf 1_E,\mathbf 1_F)
 =-\frac{q\theta_\infty\gamma_\infty}{(q-1)\gamma_{a-b}}\,q^{-a}+O_q(q^{-a-b}),\\
 \mathcal L(K_U,R_U)(X)-(P_K\otimes P_R)(X)
 =\spnTV\bigl(\mathcal L(K_U,R_U),P_K\otimes P_R\bigr)+O_q(q^{-a-b}),
\end{gather*}
so the test that exactly one of $U$ and $\Lambda(U)$ has full rank is asymptotically optimal for
detecting their dependence.
\end{corollary}

\begin{remark}[Finite enclosures]\label{spn:rem:enclosure}
By \cref{spn:lem:corank-mix} with $\eps_{a,b}=\frac12q^{-a}\sqrt{\theta_h^{-1}-1}$ and
\cref{spn:thm:uniform-full}(1),
\[
 (1-\alpha_{a,b})\max(0,\theta_h-\eps_{a,b})\le
 \spnTV(A,U)\le(1-\alpha_{a,b})\min(1,\theta_h+\eps_{a,b})
\]
for all positive even $a,b$. For the constants, after $L$ factors the tails are
$\sum_{j>L}q^{-2j}=q^{-2L}/(q^2-1)$ and $\sum_{j>L}q^{-(2j-1)}=q^{1-2L}/(q^2-1)$; a finite product
times one minus its tail is a lower bound and the finite product an upper bound. With $L=48$ and
outward rounding this gives:
\begin{center}\small
\begin{tabular}{rlrl}
\toprule
$q$ & $\Delta_q$ & $q$ & $\Delta_q$\\
\midrule
$2$ & $0.13063434670850517643$ & $5$ & $0.03300267418688550148$\\
$3$ & $0.07887849870952883542$ & $8$ & $0.01385794748300541419$\\
$4$ & $0.04897471703346141005$ & $9$ & $0.01109398786238867742$\\
\bottomrule
\end{tabular}
\end{center}
\end{remark}

The proofs occupy the rest of this subsection.

\emph{Kernel moments.} For $0\le j\le b$ and $0\le2s\le j$ the number of $j$-subspaces of
$(\F_q^b,J_b)$ on which $J_b$ restricts to rank $2s$ is
\begin{equation}\label{spn:eq:type-count}
 N_b(j,s)=\frac{\spnEmb_{b/2}(j,s)}{\abs{\Sp_{2s}(\F_q)}\,\abs{\GL_{j-2s}(\F_q)}\,q^{2s(j-2s)}},
\end{equation}
since the isometries of such a subspace are the block matrices
$\left(\begin{smallmatrix}P&0\\R&S\end{smallmatrix}\right)$ with $P\in\Sp_{2s}$, $S\in\GL_{j-2s}$
and $R$ arbitrary. Write $\spnqbin{x}{j}{q}$ for the number of $j$-subspaces of an
$x$-dimensional space, so that $\spnE\spnqbin{K}{j}{q}$ is the expected number of $j$-subspaces of
a kernel of dimension $K$.

\begin{lemma}[Kernel moments; \stU]\label{spn:lem:block-moments}
For $0\le j\le b$,
\begin{gather*}
 \mu_j=\spnE\spnqbin{K_A}{j}{q}=\sum_sN_b(j,s)\frac{\spnEmb_{b/2}(j,s)}{\spnEmb_{(a+b)/2}(j,s)},\qquad
 \mu_j^0=\spnE\spnqbin{K_U}{j}{q}=\spnqbin{b}{j}{q}q^{-aj},\\
 \nu_j=\spnE\spnqbin{R_U}{j}{q}=q^{-(a+b)j+j^2}\sum_sN_b(j,s)\spnEmb_{a/2}(j,s),\qquad
 \nu_j^0=\spnE\spnqbin{R_T}{j}{q}=\spnqbin{b}{j}{q}q^{-bj+j(j+1)/2}.
\end{gather*}
\end{lemma}

\begin{proof}
The frame \eqref{spn:eq:block-frame} is injective, since its Gram matrix $J_b$ is nonsingular, and
uniform among the frames with that Gram matrix (\cref{spn:lem:gram-frames}). A source $j$-subspace
of restricted rank $2s$ is therefore carried to a uniform one of the $\spnEmb_{(a+b)/2}(j,s)$
frames with its Gram matrix, and $A$ kills it exactly when that image lies in the second summand,
which allows $\spnEmb_{b/2}(j,s)$ frames. For $U$, killing a $j$-subspace imposes $aj$ independent
conditions. For the defect, fix a $j$-subspace $K$ of restricted rank $2s$. If
$K\subseteq\ker\Lambda(U)$ and $Ux=0$ with $x\in K$, then
$x^{\spnT}J_by=(Ux)^{\spnT}J_a(Uy)=0$ for all $y$, so $x=0$; thus $U|_K$ is an isometric
injection, with $\spnEmb_{a/2}(j,s)$ choices, and then each of $b-j$ complementary basis vectors
has $j$ independent prescribed pairings and $q^{a-j}$ images. Dividing by $q^{ab}$ gives $\nu_j$.
For $T$, killing a $j$-subspace imposes $bj-j(j+1)/2$ independent conditions.
\end{proof}

Put $m_j^{(d)}=q^{-j(j+d)}/\gamma_j$, $m'_j=q^{-j(j-1)/2}/\gamma_j$ and $Q_j=(q^j-1)/(q-1)$. For even
$D$ let $\Phi_D(j,s)=\prod_{i=D/2-j+s+1}^{D/2}(1-q^{-2i})$ on the feasible range $j-s\le D/2$ and
$\Phi_D(j,s)=0$ otherwise, so that $\spnEmb_{D/2}(j,s)=q^{Dj-\binom j2}\Phi_D(j,s)$ by
\eqref{spn:eq:E-count}. Let $H$ be uniform on $\spnAlt_j(\F_q)$ with rank $2s$ and put
$W(H)=(q^{2(j-s)}-1)/(q^2-1)$. A rank type of $H$ has probability
$\abs{\GL_j(\F_q)}/\bigl(q^{\binom j2}\abs{\Sp_{2s}(\F_q)}\abs{\GL_{j-2s}(\F_q)}q^{2s(j-2s)}\bigr)$,
so $N_b(j,s)=q^{j(b-j)}\gamma_j^{-1}\spnP(\rank H=2s)\,\Phi_b(j,s)$, and
\cref{spn:lem:block-moments} becomes
\begin{equation}\label{spn:eq:unified-moments}
 \mu_j=m_j^{(d)}\spnE_H\frac{\Phi_b^2}{\Phi_{a+b}},\quad \mu_j^0=m_j^{(d)}\spnE_H\Phi_b,\quad
 \nu_j=m'_j\spnE_H\Phi_a\Phi_b,\quad \nu_j^0=m'_j\spnE_H\Phi_b,
\end{equation}
with $\Phi_D=\Phi_D(j,s)$; for $j\le b\le a$ every type is feasible in dimension $a+b$ and
$\Phi_{a+b}\ge\eta_\infty$. On the feasible range the omitted factors of $\Phi_D$ sum to exactly
$q^{-D}W(H)$, so $0\le\prod(1-x_i)-1+\sum x_i\le\frac12(\sum x_i)^2$ gives
$0\le1-\Phi_D\le q^{-D}W(H)$ and $\Phi_D=1-q^{-D}W(H)+O_q(q^{-2D}q^{4j})$; on the infeasible range
$\Phi_D=0$ and $q^{-D}W(H)\ge1$, and the same bounds hold. A fixed nonzero vector is killed by $H$
with probability $q^{-(j-1)}$, so $\spnE_Hq^{\dim\ker H}=1+(q^j-1)q^{1-j}$ and, as
$2(j-s)=j+\dim\ker H$, $\spnE_HW(H)=Q_j$. Writing
$\Phi_b^2/\Phi_{a+b}-\Phi_b=(\Phi_b-1)+(\Phi_b-1)^2+\Phi_b^2(\Phi_{a+b}^{-1}-1)$ with
$\abs{\Phi_{a+b}^{-1}-1}\le\eta_\infty^{-1}\sigma\tau W$, and
$\Phi_b(\Phi_a-1)=(\Phi_a-1)+(\Phi_b-1)(\Phi_a-1)$ with $\sigma\le\tau$, gives for $1\le j\le b$
\begin{equation}\label{spn:eq:first-moment-expansion}
 \mu_j-\mu_j^0=-\tau m_j^{(d)}Q_j+O_q\Bigl(\sigma\tau\frac{q^{-j^2+4j}}{\gamma_j}\Bigr),\qquad
 \nu_j-\nu_j^0=-\sigma m'_jQ_j+O_q(\sigma\tau m'_jq^{4j}),
\end{equation}
while both differences vanish at $j=0$. The first remainder is where uniformity in $d$ enters: the
available bound $\tau^2m_j^{(d)}q^{4j}$ is at most $\sigma\tau q^{-j^2+4j}/\gamma_j$ because
$q^{-jd}\le q^{-d}$ for $j\ge1$ and $\tau^2q^{-d}=\sigma\tau$.

Counting matrices by their kernel, and alternating forms by their radical, the nullity of a uniform
$(m+d)\times m$ matrix has the law
\begin{equation}\label{spn:eq:finite-rect}
 \pi_{m,d}(k)=q^{-k(k+d)}\frac{\gamma_{m+d}\gamma_m}{\gamma_k\gamma_{k+d}\gamma_{m-k}}\qquad
 (0\le k\le m),
\end{equation}
and \eqref{spn:eq:finite-alt} holds; their tails are at most $C_qq^{-k(k+d)}$ and
$C_qq^{-k(2k-1)}$, so the limits \eqref{spn:eq:limit-laws} are probability laws whose Gaussian
moments $\sum_k\spnqbin kjq\pi_d(k)=m_j^{(d)}$ and $\sum_k\spnqbin kjq\lambda_q(k)=m'_j$ are the
limits of the finite kernel counts $\spnqbin mjqq^{-(m+d)j}$ and $\spnqbin{2h}jqq^{-2hj+j(j+1)/2}$.
The identity $(q^k-q^j)\spnqbin kjq=q^j(q^{j+1}-1)\spnqbin k{j+1}q$ and the relations
$(q^{j+1}-1)m_{j+1}^{(d)}=q^{-j-d}m_j^{(d)}$, $(q^{j+1}-1)m'_{j+1}=qm'_j$ show that the
$q^k$-tilted moments of $\pi_d$ and $\lambda_q$ are $(q^j+q^{-d})m_j^{(d)}$ and $(q+1)q^jm'_j$, so
\begin{equation}\label{spn:eq:signed-moments}
 \sum_k\spnqbin kjq\xi_d(k)=-m_j^{(d)}Q_j,\qquad
 \sum_k\spnqbin kjq\zeta_q(k)=-m'_jQ_j,
\end{equation}
matching the first-order terms of \eqref{spn:eq:first-moment-expansion}; the case $j=0$ gives
total mass zero, and $1+q^{-d}-q^k\le0$ for $k\ge1$ and $q+1-q^k<0$ for $k\ge2$ give
\eqref{spn:eq:signed-norms}.

\begin{lemma}[Interpolation; \stU]\label{spn:lem:interpolation}
For $0\le k\le N$ let $L_{k,N}(x)=\prod_{0\le j\le N,\,j\ne k}(x-q^j)/(q^k-q^j)$, which equals the
indicator of $\set{z=k}$ at $x=q^z$, $0\le z\le N$, and let $L^{\mathrm{ev}}_{k,h}$ be the analogous
polynomial on the grid $1,q^2,\dots,q^{2h}$, indicating nullity $2k$. Then for every $i$
\[
 \abs{[x^i]L_{k,N}}\le\gamma_\infty^{-3}q^{-k(k+1)/2}q^{-i(i-1)/2},\qquad
 \abs{[x^i]L^{\mathrm{ev}}_{k,h}}\le\eta_\infty^{-3}q^{-k(k+1)}q^{-i(i-1)},
\]
and for $z>N$, respectively even $z>2h$,
\[
 \abs{L_{k,N}(q^z)}\le\gamma_\infty^{-2}q^{Nz-N(N+1)/2-k(k-1)/2},\qquad
 \abs{L^{\mathrm{ev}}_{k,h}(q^z)}\le\eta_\infty^{-2}q^{hz-h(h+1)-k(k-1)}.
\]
\end{lemma}

\begin{proof}
Writing the factors as multiples of $1-xq^{-j}$, the constant prefactor of $L_{k,N}$ has
absolute value $q^{-k(k+1)/2}/(\gamma_k\gamma_{N-k})$. The sum of the products of $i$ distinct
members of $1,q^{-1},q^{-2},\dots$ is $q^{-i(i-1)/2}/\gamma_i$, by removing the staircase
$0,1,\dots,i-1$ and summing the nonnegative gaps; this bounds the coefficients after one index is
omitted. For the values, bound each numerator factor by $q^z$ and use
$\prod_{j\ne k}\abs{q^k-q^j}=q^{N(N+1)/2+k(k-1)/2}\gamma_k\gamma_{N-k}$. Replacing $q$ by $q^2$ gives
the even case.
\end{proof}

Counting maps $\F_q^i\to\F_q^z$ by their images gives
\begin{equation}\label{spn:eq:power-conversion}
 q^{iz}=\sum_{j=0}^i\Bigl(\prod_{r=0}^{j-1}(q^i-q^r)\Bigr)\spnqbin zjq .
\end{equation}

\begin{proof}[Proof of \cref{spn:thm:uniform-marginals}]
Let $\rho=\mathcal L(K_A)-\mathcal L(K_U)-\tau\xi_d$. By \eqref{spn:eq:first-moment-expansion} and
\eqref{spn:eq:signed-moments} its $j$th Gaussian moment vanishes at $j=0$ and is at most
$C_q\sigma\tau q^{-j^2+4j}$ for $1\le j\le b$, so by \eqref{spn:eq:power-conversion} its $i$th power
moment, $i\le b$, is at most $C_q\sigma\tau\sum_{j\le i}q^{-j^2+(i+4)j}\le C'_q\sigma\tau q^{(i+4)^2/4}$
in absolute value. Pairing with the coefficients of $L_{k,b}$ leaves the convergent exponent
$-i^2/4+5i/2+4$, so $\abs{\sum_zL_{k,b}(q^z)\rho(z)}\le C_q\sigma\tau q^{-k(k+1)/2}$. Beyond the grid
only $-\tau\xi_d$ survives, and $\abs{\xi_d(z)}\le C_qq^{-d}q^{-z^2+z}$ for $z\ge1$ uniformly in $d$;
with $z=b+1+t$ the tail bound of \cref{spn:lem:interpolation} leaves the exponent $-(b+1)t-t^2$, so
$\sum_{k\le b}\abs{\sum_{z>b}L_{k,b}(q^z)\xi_d(z)}\le C_qq^{-d}q^{-b(b+1)/2}$, and after
multiplication by $\tau$ this is $O_q(\sigma\tau)$; so is the mass $\tau\sum_{z>b}\abs{\xi_d(z)}$.
Summing over $k$ proves the rank expansion.

For the defect let $\rho'=\mathcal L(R_U)-\mathcal L(R_T)-\sigma\zeta_q$, supported on even
nullities. For $i\le h$ its power-moment remainder is at most
$C_q\sigma\tau\sum_{j\le i}q^{-j(j-1)/2+(i+4)j}\le C'_q\sigma\tau q^{i(i+9)/2}$, the terms increasing
by a factor at least $q^5$ up to $j=i$, and the even-grid coefficients leave the convergent exponent
$-i^2/2+11i/2$, so $\abs{\sum_zL^{\mathrm{ev}}_{k,h}(q^z)\rho'(z)}\le C_q\sigma\tau q^{-k(k+1)}$. For
even $z\ge2$, $\abs{\zeta_q(z)}\le C_qq^{-z(z-3)/2}$; with $z=2h+2+2t$ the exponent of the
interpolated tail is $1-h^2$ at $t=0$ and decreases quadratically, so after multiplication by
$\sigma$ the tail is at most $C_q\sigma q^{1-h^2}\le C'_q\sigma q^{-2h}=C'_q\sigma\tau$, because
$1-h^2+2h\le2$; the direct tail mass is smaller. This proves the defect expansion.

Finally $\abs{\frac12\norm{x+y}_1-\frac12\norm x_1}\le\frac12\norm y_1$ and
\eqref{spn:eq:signed-norms} give the two distances, and since the only positive mass of $\xi_d$ is
at zero, the event of full column rank attains the leading term.
\end{proof}

\begin{proof}[Proof of \cref{spn:thm:uniform-full}]
Part~(1) is \cref{spn:thm:balanced}. The Gram condition \eqref{spn:eq:block-frame} says
$\Lambda(A)=B^{\spnT}J_bB$, which is nonsingular exactly when $B$ is invertible, and on that event the
likelihood is the constant $\alpha_{a,b}$, so
$\spnP(B\text{ invertible})=\alpha_{a,b}\spnP(\Lambda(U)\text{ nonsingular})$.

For (2) let $a\ge b$. Since $\spnP(R_T=0)=\theta_{b/2}$, \cref{spn:thm:uniform-marginals} at
nullity zero gives
\begin{equation}\label{spn:eq:defect-zero}
 \spnP\bigl(\Lambda(U)\text{ nonsingular}\bigr)=\theta_{b/2}+\sigma\frac{q\theta_\infty}{q^2-1}
 +O_q(\sigma\tau).
\end{equation}
The ratio $\eta_{a/2}/\eta_{(a+b)/2}$ is the reciprocal of a product whose omitted factors sum to
$\sigma(1-\tau)/(q^2-1)$, so $\alpha_{a,b}=\eta_{b/2}\bigl(1+\sigma/(q^2-1)+O_q(\sigma\tau)\bigr)$.
Substituting in (1),
\[
 \spnTV(A,U)=(1-\eta_{b/2})\theta_{b/2}
 +\frac{\sigma}{q^2-1}\bigl((1-\eta_{b/2})q\theta_\infty-\eta_{b/2}\theta_{b/2}\bigr)+O_q(\sigma\tau),
\]
and replacing the finite products in the coefficient by their limits costs $O_q(\sigma\tau)$.
The product tails give $\eta_{b/2}=\eta_\infty(1+\tau/(q^2-1)+O_q(\tau^2))$ and
$\theta_{b/2}=\theta_\infty(1+q\tau/(q^2-1)+O_q(\tau^2))$, hence
$(1-\eta_{b/2})\theta_{b/2}=\Delta_q+c_q\tau+O_q(\tau^2)$, which is (3) for $a\ge b$; (3) without
ordering is also the case $e=0$ of \cref{spn:cor:offset-first}, and the transpose of
$Y^{\spnT}JgX_0$ is $-X_0^{\spnT}Jg^{-1}Y$ with $g^{-1}$ again uniform. Finally
$c_q<0$ because $\eta_\infty>q/(q+1)$: for $q>2$ this follows from
$\eta_\infty\ge(q^2-2)/(q^2-1)$, and for $q=2$ the product inequality is strict, so
$\eta_\infty>2/3$.
\end{proof}

For the joint law, the conditioning that defines the defect moments leaves an independent uniform
matrix of smaller size.

\begin{lemma}[Finite rectangular nullity laws; \stU]\label{spn:lem:finite-rect}
For all integers $m,d\ge0$,
$\norm{\pi_{m,d}-\pi_d-q^{-m}\xi_d}_1\le C_q\,q^{-2m-d}$.
\end{lemma}

\begin{proof}
Put $t=q^{-m}$ and $\epsilon=q^{-d}$. For $0\le k\le m$,
$\pi_{m,d}(k)/\pi_d(k)=\prod_{i=m-k+1}^m(1-q^{-i})/\prod_{i>m+d}(1-q^{-i})$, and the product
inequality gives
$\pi_{m,d}(k)/\pi_d(k)=1+t(1+\epsilon-q^k)/(q-1)+O_q\bigl(t^2((q^k-1)^2+\epsilon(q^k-1)+\epsilon^2)\bigr)$
uniformly, every reciprocal product being at most $\gamma_\infty^{-1}$. At $k=0$ the error is
$O_q(t^2\epsilon^2)$, and for $k\ge1$, $\pi_d(k)\le C_q\epsilon q^{-k^2}$, so the weighted errors
sum to $O_q(t^2\epsilon)$. For $k>m$ the finite law vanishes, and
$\pi_d(k)\le C_q\epsilon q^{-k^2}$, $\abs{\xi_d(k)}\le C_q\epsilon q^{-k^2+k}$ bound the tails of
$\pi_d+t\xi_d$ by $C_qt^2\epsilon$.
\end{proof}

\begin{lemma}[Conditioning on a residual radical; \stU]\label{spn:lem:conditional}
Let $K\le\F_q^b$ be a $j$-subspace on which $J_b$ has rank $2s$. Whenever the event has positive
probability, the conditional law of $\dim\ker U$ given $K\subseteq\ker\Lambda(U)$ is
$\pi_{b-2j+2s,d}$.
\end{lemma}

\begin{proof}
Write $K=S\oplus R$ with $S$ nondegenerate of dimension $2s$ and $R$ the radical, of dimension
$r=j-2s$; choose dual partners $R^*$ and complete to
$\F_q^b=S\perp(R\oplus R^*)\perp W_0$ with $\dim W_0=b-2j+2s$. The condition says
$\langle Ux,Uy\rangle_a=\langle x,y\rangle_b$ for $x\in K$ and all $y$, so $U|_K$ is an isometric
injection, as in \cref{spn:lem:block-moments}. Fix it, and make the corresponding choices in the
target, with residual space $W_1$ of dimension $a-2j+2s$. In the decompositions ordered as
$S,R,R^*,W_0$ and $U(S),U(R),U(R)^*,W_1$ the admissible extensions are exactly
\[
 U=\begin{pmatrix}I_{2s}&0&0&0\\0&I_r&C&E\\0&0&I_r&0\\0&0&G&F\end{pmatrix}
\]
with $C,E,G,F$ arbitrary: pairings with $S$ fix the first block row and pairings with $R$ fix the
third, and nothing constrains pairings between two vectors outside $K$. So the free entries are
independent and uniform, and $t\mapsto(0,-Et,0,t)$ identifies $\ker F$ with $\ker U$, where $F$ is a
uniform $(a-2j+2s)\times(b-2j+2s)$ matrix. This law does not depend on the chosen injection.
\end{proof}

\begin{proof}[Proof of \cref{spn:thm:joint-defect}]
Let $\mathcal C=\mathcal L(K_U,R_U)-P_K\otimes P_R$, and regard its Gaussian moment in the second
coordinate as an element $v_j$ of $\ell^1(\N_0)$ in the first. By \cref{spn:lem:conditional} and
the count behind $\nu_j$, for $j\le h$, where every type is feasible,
\[
 v_j=m'_j\spnE_H\bigl[\Phi_a\Phi_b\,(\pi_{b-r,d}-\pi_{b,d})\bigr],\qquad r=2(j-s),
\]
since $P_K=\pi_{b,d}$; also $v_0=0$. By \cref{spn:lem:finite-rect},
$\pi_{b-r,d}-\pi_{b,d}=\tau(q^r-1)\xi_d+O_{\ell^1}(\sigma\tau q^{2r})$ with constants depending only on
$q$. Replacing $\Phi_a\Phi_b$ by one in the leading term costs
$O_{\ell^1}(\sigma\tau m'_jq^{4j})$, since $\norm{\xi_d}_1=O_q(q^{-d})$,
$1-\Phi_a\Phi_b\le2\tau W$ and $q^r-1=(q^2-1)W$; and $\spnE_HW=Q_j$. Hence
\[
 v_j=\tau(q^2-1)Q_jm'_j\,\xi_d+O_{\ell^1}(\sigma\tau m'_jq^{4j}),
\]
and by \eqref{spn:eq:signed-moments} the leading term is the Gaussian moment of
$-\tau(q^2-1)\xi_d\otimes\zeta_q$. The interpolation estimates used for the defect law hold
verbatim for sequences with values in a normed space, with norms in place of absolute values. So
the even-grid argument applies to $\mathcal C+\tau(q^2-1)\xi_d\otimes\zeta_q$ with values in
$\ell^1(\N_0)$: its Gaussian moment errors have the bounds used for $\rho'$, and its tail beyond the
grid is multiplied by $\tau\norm{\xi_d}_1=O_q(\sigma)$, contributing
$O_q(\sigma q^{1-h^2})=O_q(\sigma\tau)$. The first coordinate is already summed in the norm, so no
second inversion is needed. This proves the expansion. Half the $\ell^1$ norm of the leading
tensor is $\frac12\tau(q^2-1)\norm{\xi_d}_1\norm{\zeta_q}_1$, which by \eqref{spn:eq:signed-norms}
is the stated leading term.
\end{proof}

A second proof computes the mixed Gaussian moments directly: by \cref{spn:lem:conditional},
$\spnE\bigl[\spnqbin{K_U}iq\spnqbin{R_U}jq\bigr]=m_i^{(d)}m'_j\spnE_H\bigl[\Phi_a\Phi_b
\Psi_{b-2j+2s}(i)\bigr]$ with $\Psi_D(i)=\prod_{r<i}(1-q^{r-D})$, and two interpolations, on the
ordinary and on the even grid, recover the joint law with the same remainder.

\begin{proof}[Proof of \cref{spn:cor:xor}]
The $(0,0)$ entry of the expansion gives
$\operatorname{Cov}(\mathbf 1_E,\mathbf 1_F)=\mathcal C(E\cap F)=-\tau(q^2-1)\xi_d(0)\zeta_q(0)
+O_q(\sigma\tau)$. Since $\mathcal C$ has zero marginals, $\mathcal C(E\cap F^c)=\mathcal C(E^c\cap
F)=-\mathcal C(E\cap F)$, so $\mathcal C(X)=-2\mathcal C(E\cap F)$, which is the leading term of
\cref{spn:thm:joint-defect}.
\end{proof}

For $a=b$ the Gram matrix $U^{\spnT}J_aU$ itself is not asymptotically uniform among alternating
matrices: its density relative to $T$ at corank $2k$ is $\eta_{b/2}F(k)$, with $F(k)$ the sum in
\cref{spn:thm:balanced}, which depends on $k$ for all $b$; for $a-b\to\infty$ it is
(\cref{spn:lem:alt-gram}). Tilting the joint law by the likelihood recovers the law of $K_A$ by a
second route.

\paragraph{Exact checks.} Exact rational evaluation of the laws by the subspace counts and
triangular moment inversion, exact joint laws by mixed moments, and element-level enumeration of
all pairs $(A,B)$ in small cases, grouping $B$ by $B^{\spnT}J_bB$ and not using the embedding or
moment formulas, agree with the theorems in every case tested; the cases $q=4$ use genuine $\F_4$
arithmetic. Enumerated cases give:
\begin{center}\small
\begin{tabular}{rrrccc}
\toprule
$q$&$a$&$b$&$\spnTV(A,U)$&$\spnTV(\rank A,\rank U)$&$\spnTV(\mathcal L(K_U,R_U),P_K\otimes P_R)$\\
\midrule
$2$&$2$&$2$&$1/8$&$1/8$&$15/32$\\
$2$&$4$&$2$&$85/672$&$115/2688$&$345/2048$\\
$2$&$4$&$4$&$6767/52224$&$1087/52224$&$1959705/33554432$\\
$2$&$6$&$2$&$273/2176$&$399/34816$&$5985/131072$\\
$3$&$2$&$2$&$19/270$&$11/270$&$176/729$\\
$3$&$4$&$2$&$1630/22113$&$1070/199017$&$17120/531441$\\
$4$&$4$&$2$&$13073/279552$&$5423/4472832$&\\
\bottomrule
\end{tabular}
\end{center}
and for $(q,a,b)$ equal to $(4,2,2)$ and $(5,2,2)$ the full and rank distances are $49/1088$,
$19/1088$ and $101/3250$, $29/3250$. Over $\F_2$ along the gap-two sequence the exact distances
are:
\begin{center}\small
\begin{tabular}{rrll}
\toprule
$a$&$b$&$\spnTV(\rank A,\rank U)$&$\spnTV(A,U)$\\
\midrule
$6$&$4$&$1.1816321888\cdot10^{-2}$&$0.129825251831$\\
$10$&$8$&$7.5131145533\cdot10^{-4}$&$0.130589162483$\\
$18$&$16$&$2.9376932110\cdot10^{-6}$&$0.130634171739$\\
$26$&$24$&$1.1475407711\cdot10^{-8}$&$0.130634346025$\\
$50$&$48$&$6.8398761046\cdot10^{-16}$&$0.130634346709$\\
\bottomrule
\end{tabular}
\end{center}
For the binary $26\times24$ block the leading rank term of \cref{spn:thm:uniform-marginals} differs
from the exact rank distance by a relative $1.49\cdot10^{-8}$, the defect distance is
$4.1666\cdot10^{-9}$ and the joint-versus-product distance $1.9252\cdot10^{-8}$. For binary
$48\times24$ and $64\times32$ blocks the rank distances are $3.5527\cdot10^{-15}$ and
$5.4210\cdot10^{-20}$, while the full distances are $0.1306343462$ and $0.1306343467$; for $q=2$
and $a=b=24$ the rank distance is $1.72131\cdot10^{-8}$. Enumeration of $329{,}957$ frame pairs
over $\F_2$, $\F_3$ and $\F_4$ checks the likelihood signs and the optimal-test identity, and
enumeration over $\F_2$ and $\F_3$ checks \cref{spn:thm:complement-boundary} on $85$ Gram-matrix
comparisons and $5$ exact boundary
distances.\footnote{\raggedright\artifact{uniform_symplectic_blocks/results/exhaustive_results.json},
\artifact{uniform_symplectic_blocks/results/exact_distances.csv},
\artifact{rectangular_symplectic_blocks/results/exact_laws.json},
\artifact{rectangular_symplectic_blocks/results/joint_exact.json},
\artifact{rectangular_symplectic_blocks/results/independent_enumeration.json},
\artifact{rectangular_symplectic_blocks/results/constant_enclosures.json},
\artifact{grh_uniform_extensions_2026-11-27/data/symplectic_checks.csv},
\artifact{grh_uniform_extensions_2026-11-27/data/symplectic_verification.json},
\artifact{nondegenerate_symplectic_blocks_2026-11-27/results/exhaustive.json},
\artifact{nondegenerate_symplectic_blocks_2026-11-27/results/exact_distances.csv},
\artifact{nondegenerate_symplectic_blocks_2026-11-27/results/limiting_constants.json},
\artifact{exact_symplectic_compressions_2026-11-26/results/exhaustive_checks.json},
\artifact{exact_symplectic_compressions_2026-11-26/results/certified_constants.json},
\artifact{symplectic_compression_2026-11-26_2/results.json}}
Sharp convergence of classical rank laws is due to \cite{Fulman2015Stein}; exact
rank-intersection counts for pairs of alternating matrices are in \cite{PottSchmidtZhou2016Pairs},
direct sums of nondegenerate spaces are studied in \cite{GlasbyNiemeyerPraeger2022Random}, and
\cite{DeBoeckVandeVoorde2024Anzahl} contains the counts behind the scalar probability
$\spnP(\Lambda(U)\text{ nonsingular})$ and the single probability $\spnP(B\text{ invertible})$.

%======================================================================

\subsection{Many frames: total dimension, hidden dependence and cycle laws}\label{spn:sec:many}
\sources{symplectic_joint_observations_2026-11-27; symplectic_dependence_research_2026-11-27}

\Cref{spn:thm:joint-law} gives the exact law of all pairings among $k$ independent frames. Compared
with independent uniform cross pairings, this law has a transition at total frame dimension equal to
the ambient dimension, located by one exact kernel moment; inside the critical window the limiting
distance is universal when no frame dominates; and observations indexed by a forest are exactly
independent, so the dependence can be invisible in every proper collection of frames. For the
simplest frames, single nonzero vectors observed along the edges of a graph, every cycle law is
explicit and the detection threshold for a growing cactus is sharp.

The notation is that of \cref{spn:thm:joint-law}: $V$ is symplectic of dimension $2n$ over $\F_q$,
the frames $X_i\colon(E_i,H_i)\to V$, $1\le i\le k$, are independent and uniform, $d_i-s_i\le n$,
$D=\sum_id_i$, and $\Gamma$ is the assembled alternating form. Let $\mu$ be its law and $\nu$ the law
that keeps the diagonal blocks $H_i$ and makes every entry of every cross block independent and
uniform, so that $\nu(\Gamma)=q^{-B}$. Put $\eta_j=\prod_{i=1}^j(1-q^{-2i})$,
$\gamma_j=\prod_{i=1}^j(1-q^{-i})$ and $\theta_\infty=\gamma_\infty/\eta_\infty$, with empty products
equal to one. All results hold in every characteristic.

\paragraph{The transition at total dimension $2n$.}
Define
\begin{equation}\label{spn:eq:joint-A}
 A=1+\Bigl(q^D-1-\sum_i(q^{d_i}-1)\Bigr)q^{1-D}+\sum_i(q^{d_i-2s_i}-1)q^{d_i-D}.
\end{equation}

\begin{theorem}[Uniform joint comparison; \stU]\label{spn:thm:joint-threshold}
\begin{enumerate}
\item If $c=2n-D\ge0$, then $\spnTV(\mu,\nu)\le\min\set{1,Aq^{-c}/(q^2-1)}$.
\item If $w=D-2n>0$, then $\spnTV(\mu,\nu)\ge1-(A-1)/(q^w-1)$: the test $\rank\Gamma\le2n$ accepts
with probability one under $\mu$ and with probability at most $(A-1)/(q^w-1)$ under $\nu$.
\item $A=\spnE_\nu q^{\dim\ker\Gamma}$ and $A\le q+1+\sum_iq^{2(d_i-s_i)-D}$; in particular
$A\le q+1+k$ when $\max_i(d_i-s_i)\le D/2$, which always holds when $D>2n$. Hence, for $c\ge0$,
\begin{equation}\label{spn:eq:joint-shape}
 \spnTV(\mu,\nu)\le\frac{(q+1)q^{D-2n}+\sum_iq^{-2(n-d_i+s_i)}}{q^2-1}.
\end{equation}
\end{enumerate}
\end{theorem}

\begin{corollary}[Total dimension determines the transition; \stU]\label{spn:cor:joint-transition}
Fix $q$ and $k$. Along any sequence of admissible frames with $\max_i(d_i-s_i)\le D/2$, the distance
$\spnTV(\mu,\nu)$ tends to zero if $2n-D\to+\infty$ and to one if $D-2n\to+\infty$, with errors
$O_{q,k}(q^{-\abs{D-2n}})$.
\end{corollary}

For two frames, \eqref{spn:eq:joint-shape} improves the exponent $(2n+1-a-b)/2$ of
\cref{spn:thm:symplectic} to $2n-a-b$ when the frames are balanced. For unbalanced frames, for
instance a Lagrangian one, a term $q^{-2(n-d_i+s_i)}$ is not small, and \cref{spn:thm:symplectic}
remains the bound that is uniform in the frames. The bounds do not give a limit for fixed $D-2n$;
that is \cref{spn:thm:joint-profile} below.

\begin{lemma}[The evaluation map; \stU]\label{spn:lem:joint-eval}
Let $\mathcal W$ be the space of alternating forms on $E$ with zero diagonal blocks. For nonzero
$v=(v_1,\dots,v_k)\in E$ the linear map $\mathcal W\to E^*$, $B\mapsto Bv$, has rank $D-d_i$ if
$v_i$ is the only nonzero component of $v$, and otherwise rank $D-1$, with image the annihilator of
$v$.
\end{lemma}

\begin{proof}
If only $v_i$ is nonzero, the $i$th output block vanishes and every other output block is free,
obtained by varying the cross block incident to $i$. If at least two components are nonzero, a
vector $w$ annihilates every $Bv$ exactly when $w_iv_j^{\spnT}-v_iw_j^{\spnT}=0$ for all $i<j$. For two
nonzero components $v_i,v_j$, equality of these rank-one products forces $w_i=\lambda v_i$ and
$w_j=\lambda v_j$ with one scalar $\lambda$, also when one of $w_i,w_j$ vanishes, and pairing a
nonzero component with every other one gives $w_h=\lambda v_h$ for all $h$. So the cokernel is the
line of $v$; by alternation the image lies in the annihilator of $v$, which has the same dimension.
No division by two occurs.
\end{proof}

\begin{lemma}[Exact first kernel moment; \stU]\label{spn:lem:joint-moment}
$\spnE_\nu q^{\dim\ker\Gamma}=A$. If every $H_i$ is zero,
$A=q+1-q^{1-D}+q^{-D}\sum_i(q^{d_i}-1)(q^{d_i}-q)$.
\end{lemma}

\begin{proof}
$\spnE_\nu\abs{\ker\Gamma}=\sum_{v\in E}\spnP_\nu(\Gamma v=0)$. The zero vector contributes one. A
nonzero $v$ supported in $E_i$ can be in the kernel only if $v_i\in\operatorname{rad}H_i$; there are
$q^{d_i-2s_i}-1$ such $v$, each with probability $q^{-(D-d_i)}$ by \cref{spn:lem:joint-eval}. For
each of the $q^D-1-\sum_i(q^{d_i}-1)$ vectors supported in at least two blocks, the random part of
$\Gamma v$ is uniform on the annihilator of $v$, and the fixed diagonal part lies in that annihilator
because $\sum_iH_i(v_i,v_i)=0$, so the probability is $q^{-(D-1)}$. Summing gives
\eqref{spn:eq:joint-A}, and $s_i=0$ gives the isotropic form. The vectors supported in at least two
blocks contribute at most $q$, and those in $E_i$ at most $q^{2(d_i-s_i)-D}$, which is the bound in
\cref{spn:thm:joint-threshold}(3).
\end{proof}

\begin{lemma}[Injective realizations; \stU]\label{spn:lem:joint-injective}
Let $c=2n-D\ge0$, and let $\mu_0\le\mu$ be the contribution of the frame tuples whose combined map
$E\to V$ is injective. At every $\Gamma$ of nullity $r$,
\[
 \frac{\mu_0(\Gamma)}{\nu(\Gamma)}=\mathbf 1_{\set{r\le c}}
 \frac{\prod_i\eta_{n-d_i+s_i}}{\eta_n^{k-1}\,\eta_{(c-r)/2}} ,
\]
and $r\equiv c\pmod2$ when $r\le c$.
\end{lemma}

\begin{proof}
This is the term $\ell=0$ of \eqref{spn:eq:joint-law}, where $T_0=1$, multiplied by $q^B$. By
\eqref{spn:eq:E-count}, $\spnEmb_n(a,s)=q^{2na-\binom a2}\eta_n/\eta_{n-a+s}$ when $a-s\le n$, the
powers of $q$ cancel because $B=\binom D2-\sum_i\binom{d_i}2$, and $n-D+s=(c-r)/2$. The embedding
count vanishes exactly when $r>c$, and an alternating form has even rank, so $r\equiv D\equiv c$.
\end{proof}

\begin{proof}[Proof of \cref{spn:thm:joint-threshold}]
(1) The sequence $\eta_j$ decreases and $0\le n-d_i+s_i\le n$, so the constant
$\prod_i\eta_{n-d_i+s_i}/\eta_n^{k-1}$ is at least $\eta_n\ge\eta_\infty$. If $r\le c$ put
$m=(c-r)/2$; then $f_0=\mu_0/\nu\ge\eta_\infty/\eta_m=\prod_{j>m}(1-q^{-2j})\ge1-q^{-2m}/(q^2-1)
=1-q^{r-c}/(q^2-1)$. If $r>c$, parity gives $r-c\ge2$ and $q^{r-c}/(q^2-1)>1$. So
$(1-f_0)_+\le q^{r-c}/(q^2-1)$ everywhere, and since $f=\mu/\nu\ge f_0$,
\[
 \spnTV(\mu,\nu)=\spnE_\nu(1-f)_+\le\spnE_\nu(1-f_0)_+\le\frac{q^{-c}}{q^2-1}\spnE_\nu q^r
 =\frac{A}{q^2-1}q^{-c},
\]
by \cref{spn:lem:joint-moment}. It is the missing likelihood that is bounded, not the
contributions of the dependent tuples one by one.

(2) Every actual observation is $F^{\spnT}JF$, of rank at most $2n$ and nullity at least $w$. Under
$\nu$, Markov's inequality for $q^r-1\ge0$ gives $\spnP_\nu(r\ge w)\le(A-1)/(q^w-1)$. Part (3) is
\cref{spn:lem:joint-moment}, and \eqref{spn:eq:joint-shape} is (1) with that bound for $A$.
\end{proof}

Below the threshold the injective tuples alone supply almost all of the uniform likelihood, and
above it one rank obstruction rejects almost all independent data; neither step needs the law of the
invariants $(s,t_1,\dots,t_k)$.

\paragraph{The critical window when no frame dominates.}
For an integer $w$ of parity $\varepsilon\in\set{0,1}$ put
\begin{equation}\label{spn:eq:joint-profile}
 \delta_q(w)=
 \begin{cases}
 \dfrac{\theta_\infty}{\gamma_\varepsilon}\Bigl(1-\dfrac{\eta_\infty}{\eta_{(\varepsilon-w)/2}}\Bigr),
 &w\le1,\\[10pt]
 \displaystyle\sum_{0\le j<w,\ j\equiv w\ (\mathrm{mod}\ 2)}\frac{\theta_\infty q^{-j(j-1)/2}}{\gamma_j},
 &w\ge2.
 \end{cases}
\end{equation}

\begin{theorem}[Universal critical profile; \stU]\label{spn:thm:joint-profile}
Fix $q$ and an integer $w$, let $D=2n+w$, and let $d_{\max}=\max_id_i=o(D)$. Then, uniformly over all
partitions of $D$ and all prescribed forms $H_i$, degenerate ones included,
\[
 \spnTV(\mu,\nu)\longrightarrow\delta_q(w).
\]
For $w\le1$ the test $\dim\ker\Gamma>\varepsilon$, and for $w\ge2$ the test $\rank\Gamma\le2n$,
attains the limiting distance. Quantitatively, if $D\ge8$ and $d_{\max}\le n$, then
\begin{equation}\label{spn:eq:joint-profile-bound}
 \bigl|\spnTV(\mu,\nu)-\Delta_{q,D,n}\bigr|
 \le\frac{q^{2d_{\max}-\floor{D/4}}}{1-q^{-1}}+\beta^D+e^{C_1S}-1+C_2S,
\end{equation}
where $\Delta_{q,D,n}$ is the complement discrepancy \eqref{spn:eq:complement-discrepancy} with
$b=D$ and $c=n$, $\beta=(151959/160000)^{1/4}<1$, $S=\sum_iq^{2d_i-2n}\le Dq^{2d_{\max}-2n}$,
$C_1=q^{-3}/(\eta_\infty(1-q^{-3}))$ and $C_2=\eta_\infty^{-1}(C_1+(q^2-1)^{-1})$; and
$\Delta_{q,D,n}=\delta_q(w)+O_{q,w}(q^{-2n})$.
\end{theorem}

At even $w$ the profile is $\delta_q(w)=\mathcal D_q(-w/2)$, the two-frame profile of
\cref{spn:thm:offset-profile} at offset $e=-w/2$. The two theorems never apply to the same
configuration, since two frames have $d_{\max}\ge D/2$; together they show that the even profile is
the same for two large nondegenerate frames and for arbitrarily many small frames with arbitrary
forms. The odd values, rounded to six places, are:
\begin{center}\small
\begin{tabular}{lrrr}
\toprule
 & $w=-1$ & $w=1$ & $w=3$\\
\midrule
$q=2$ & $0.068743$ & $0.261269$ & $0.838845$\\
$q=3$ & $0.013294$ & $0.118318$ & $0.958507$\\
$q=4$ & $0.004096$ & $0.065300$ & $0.983350$\\
$q=5$ & $0.001653$ & $0.041253$ & $0.991669$\\
\bottomrule
\end{tabular}
\end{center}
Over $\F_2$ the limit is about $0.130634$ at $D=2n$ and $0.261269$ at $D=2n+1$. As $w\to-\infty$
along a parity class, $\delta_q(w)\sim\theta_\infty q^{w-\varepsilon}/(\gamma_\varepsilon(q^2-1))$, and as
$w\to+\infty$, $1-\delta_q(w)\sim\eta_\infty^{-1}q^{-w(w-1)/2}$; these concern the limiting profile,
not offsets growing with $n$.

\begin{corollary}[Dependence invisible to bounded collections; \stU]\label{spn:cor:joint-local}
Under the hypotheses of \cref{spn:thm:joint-profile}, fix $\ell\ge1$. The joint cross pairings of
any collection of at most $\ell$ of the frames converge in total variation to independent uniform
entries, uniformly over the collection and the forms, while the distance for the complete
collection tends to $\delta_q(w)>0$.
\end{corollary}

The proof compares the nullity of a block-constrained uniform alternating matrix with an
unrestricted one, and then transfers the unrestricted Gram law to the conditioned frame experiment.

\begin{lemma}[Block-nullity universality; \stU]\label{spn:lem:block-nullity}
Let $A_D$ be uniform among the alternating $D\times D$ matrices with arbitrary prescribed diagonal
blocks of sizes at most $d_{\max}$, and let $T_D$ be uniform on $\spnAlt_D(\F_q)$. If $D\ge8$, then
\[
 \spnTV(\dim\ker A_D,\dim\ker T_D)\le\frac{q^{2d_{\max}-\floor{D/4}}}{1-q^{-1}}+\beta^D .
\]
\end{lemma}

\begin{proof}
Expose the rows and columns in a fixed order. At size $m$ let $S$ be the set of previously exposed
coordinates in the same prescribed block as the next one, so $\abs S\le d_{\max}$. Conditional on
the exposed matrix $A_m$, the new column is fixed on $S$ and independent uniform off $S$. If
$\ker A_m\cap\F_q^S=0$, its restriction to $(\ker A_m)^*$ is uniform, since the annihilator in that
dual of the image of the random coordinate subspace is exactly this intersection. Bordering an
alternating matrix of nullity $k$ by a column whose restriction to the kernel is uniform raises the
nullity to $k+1$ with probability $q^{-k}$ and lowers it to $k-1$ otherwise, with no division by
two. For each nonzero $v$ supported on $S$ the equations $A_mv=0$ outside its block are
$m-\abs S$ independent uniform equations, so $\spnP(\ker A_m\cap\F_q^S\ne0)\le q^{2d_{\max}-m}$, and
coupling with the standard alternating nullity chain from size $m_0$ onwards costs at most
$q^{2d_{\max}-m_0}/(1-q^{-1})$.

The standard chain forgets its initial nullity: couple the chains from all initial states at most
$m_0$ of the right parity by one sequence of uniform variables. The upward probability $q^{-k}$
decreases in $k$, so the coupling preserves order, states of equal parity never cross, and all
chains coalesce by the first time $\tau$ the chain started at $m_0$ reaches one (or zero). Before
$\tau$ the state is at least two, so with $r=6/5$,
$\spnE[r^{Z_{t+1}}\mid Z_t=k]=r^k(q^{-k}r+(1-q^{-k})r^{-1})\le\frac{37}{40}r^k$, and
$\spnP(\tau>L)\le r^{m_0-2}(37/40)^L$. The nullity at size $m_0$ of an unrestricted uniform
alternating matrix is at most $m_0$, of the same parity. Take $m_0=\floor{D/4}$ and $L=D-m_0$; since
$\frac65(\frac{37}{40})^3=\frac{151959}{160000}$, the last bound is at most $\beta^D$.
\end{proof}

Let $T$ be uniform on $\spnMat_{2n,D}(\F_q)$ and $G=T^{\spnT}J_nT$; this is $G_n$ of
\cref{spn:lem:alt-gram} with $b=D$, and its density relative to $T_D$ depends only on the nullity
$k$. With $w=D-2n$ and $k\equiv w\pmod 2$, choosing the kernel of $T$ inside the radical of the
Gram matrix and embedding the quotient by \cref{spn:lem:gram-frames} gives
\begin{equation}\label{spn:eq:unrestricted-gram}
 f_{n,D}(k)=\eta_n\sum_{t=\max(0,(w+k)/2)}^{k}\spnqbin ktq\frac{q^{wt-t(t+1)/2}}{\eta_{t-(w+k)/2}},
\end{equation}
an empty sum meaning zero.

\begin{lemma}[Unrestricted Gram laws; \stU]\label{spn:lem:unrestricted-gram}
Every feasible nullity $k\ge2$, that is $k\ge w$ of the parity of $w$, has $f_{n,D}(k)>1$. Write
$D=2d+\varepsilon$ with $\varepsilon\in\set{0,1}$. If $d\le n$, then
$\spnTV(G,T_D)=\spnP(\dim\ker T_D=\varepsilon)(1-\eta_n/\eta_{n-d})$; if $D\ge2n+2$, then
$\spnTV(G,T_D)=\spnP(\rank T_D>2n)$. Moreover
\[
 \spnP(\dim\ker T_D=k)=q^{-k(k-1)/2}\frac{\gamma_D}{\gamma_k\,\eta_{(D-k)/2}}\qquad(k\equiv D\ (\mathrm{mod}\ 2)),
\]
which tends to $\theta_\infty q^{-k(k-1)/2}/\gamma_k$, and for fixed $w$,
$\spnTV(G,T_D)=\delta_q(w)+O_{q,w}(q^{-2n})$.
\end{lemma}

\begin{proof}
Put $F=f_{n,D}/\eta_n$. If $(w+k)/2\ge2$, the first summand of \eqref{spn:eq:unrestricted-gram} is at
least $q^{t(t-1)/2}\ge q$, using $\spnqbin ktq\ge q^{t(k-t)}$; if $(w+k)/2=1$ it is
$\spnqbin k1qq^{1-k}\ge1+q^{-1}$. Both exceed $\eta_\infty^{-1}$, because $q\eta_\infty>1$ and
$(1+q^{-1})\eta_\infty>1$, by $\eta_\infty>1-(q^2-1)^{-1}$, strictly at $q=2$. Otherwise
$h=-(w+k)/2\ge0$; with $Q=q^{-2h-2}$ the terms $t=0,1$ give $F\ge(1+aQ)/\eta_{h+1}$ with
$a=(1-q^{1-k})/(q-1)\ge q^{-1}$, and $\eta_\infty/\eta_{h+1}\ge1-Q/(q^2-1)$, so
$\eta_\infty F\ge(1+aQ)(1-Q/(q^2-1))>1$, since $a(q^2-1-Q)>1$ for $a\ge q^{-1}$, $Q\le q^{-2}$ and
$q\ge2$. In each case $f_{n,D}(k)\ge\eta_\infty F>1$.

If $d\le n$, the only possible deficit is at the least nullity $\varepsilon$, where
$f_{n,D}(\varepsilon)=\eta_n/\eta_{n-d}$; for $\varepsilon=1$ this follows from the two terms of
\eqref{spn:eq:unrestricted-gram} or by placing the radical vector after the nondegenerate part. If
$D\ge2n+2$ the likelihood vanishes off rank at most $2n$ and exceeds one on it. The rank law follows
by choosing the radical and then a nonsingular form on the quotient; it tends to the stated limit,
whose tails are bounded by a constant times $q^{-k(k-1)/2}$, so the limits form a probability law.
For fixed $w$ the two distance formulas converge to \eqref{spn:eq:joint-profile} with the stated
error.
\end{proof}

\begin{proof}[Proof of \cref{spn:thm:joint-profile}]
Note that $\spnTV(G,T_D)=\Delta_{q,D,n}$ by definition. For a uniform $2n\times d$ matrix with
$d\le n$ let $F_{n,d}(H)$ be the density of its Gram matrix at $H$ relative to the uniform law on
$\spnAlt_d(\F_q)$. By \cref{spn:lem:alt-gram} and Fourier inversion, a character indexed by a form of
rank $2t$ contributes $q^{-2nt}$, and there are at most $\eta_\infty^{-1}q^{t(2d-2t-1)}$ such forms
\eqref{spn:eq:alt-rank-count}; summing the geometric majorant,
\begin{equation}\label{spn:eq:small-block}
 \abs{F_{n,d}(H)-1}\le C_1q^{2d-2n}.
\end{equation}
If $\rank H=2s$, the injective maps alone contribute $\Phi_{n,d,s}=\eta_n/\eta_{n-d+s}$, with
$1-\Phi_{n,d,s}\le q^{2d-2n}/(q^2-1)$; since $F_{n,d}(H)\ge\Phi_{n,d,s}\ge\eta_\infty$, a map with Gram
matrix $H$ fails to be injective with conditional probability at most $C_2q^{2d-2n}$.

Now condition a uniform $2n\times D$ matrix only on its diagonal-block Gram matrices $H_i$, and let
$P'$ be the resulting law of the cross blocks. Its density relative to $\nu$ is
$f_{n,D}(\dim\ker\Gamma)/Z$ with $Z=\prod_iF_{n,d_i}(H_i)$, by independence of the column blocks, so
$\abs{Z-1}\le e^{C_1S}-1$. Conditioning each block further on injectivity produces the frame law,
so $\spnTV(\mu,P')\le C_2S$. Let $g(k)=(1-f_{n,D}(k))_+\in[0,1]$. Since $\spnE_\nu f_{n,D}=Z$, the
Lipschitz property of $x\mapsto(1-x)_+$ gives
$\abs{\spnTV(P',\nu)-\spnE_\nu g(\dim\ker\Gamma)}\le\abs{Z-1}$. Under $\nu$, $\Gamma$ is uniform with
prescribed diagonal blocks, so replacing its nullity law by that of $T_D$ costs at most the bound of
\cref{spn:lem:block-nullity}, and $\spnE g(\dim\ker T_D)=\spnTV(G,T_D)$. This proves
\eqref{spn:eq:joint-profile-bound}, and with \cref{spn:lem:unrestricted-gram} the limit, since all
error terms tend to zero when $d_{\max}=o(D)$.

For $w\le1$ the likelihood is constant at nullity $\varepsilon$, so
$P'(\dim\ker\Gamma>\varepsilon)-\nu(\dim\ker\Gamma>\varepsilon)=\nu(\dim\ker\Gamma=\varepsilon)
(1-f_{n,D}(\varepsilon)/Z)\to\delta_q(w)$, and replacing $P'$ by $\mu$ costs $o(1)$. For $w\ge2$ every
actual Gram matrix has rank at most $2n$, and the probability under $\nu$ that this fails tends to
$\delta_q(w)$.
\end{proof}

\begin{proof}[Proof of \cref{spn:cor:joint-local}]
A collection of at most $\ell$ frames has total dimension $D'\le\ell d_{\max}=o(n)$, eventually at
most $n$. Apply \eqref{spn:eq:small-block} to its whole $D'$-column Gram matrix and to its diagonal
blocks: the unnormalized density differs uniformly from one by $O_q(q^{2D'-2n})$, the product
normalizer by $O_{q,\ell}(q^{2d_{\max}-2n})$, and conditioning on injectivity costs
$O_{q,\ell}(q^{2d_{\max}-2n})$.
\end{proof}

Enumeration of $1{,}607{,}096$ vector arrays in $16$ cases over $\F_2,\F_3,\F_4,\F_5$ checks
\eqref{spn:eq:unrestricted-gram} at all $55{,}631$ Gram fibres, with both parities and all three
regimes; $473$ block-constrained prefixes and all their $16{,}274$ bordering columns confirm the
upward probability $q^{-k}$; and $262{,}144$ arrays and $53{,}550$ frame tuples check the normalizer
$Z$, the density $f/Z$ and the injectivity conditioning, which separates the Gram-conditioned law
from the frame law for degenerate
forms.\footnote{\raggedright\artifact{symplectic_dependence_research_2026-11-27/companion/results/verification.json},
\artifact{symplectic_dependence_research_2026-11-27/companion/results/summary.json}}
Rank universality for matrices with some deterministic entries goes back to
\cite{KoenigNguyen2022Rank}, whose alternating theorem assumes odd $q$; the block argument above
works in every characteristic. Classical alternating rank laws are in \cite{Fulman2015Stein}.

\paragraph{Forests and hidden dependence.}

\begin{theorem}[Forest observations are independent; \stU]\label{spn:thm:forest}
Let $\mathcal F$ be a forest on the frame labels $\set{1,\dots,k}$. The cross blocks $M_{ij}$,
$ij\in\mathcal F$, are mutually independent, for arbitrary prescribed forms and at every finite size.
\end{theorem}

\begin{proof}
Choose a leaf $i$ with neighbour $j$ and condition on all frames except $X_i$. The law of $M_{ij}$
does not depend on the value of $X_j$: by \cref{spn:lem:gram-frames} any two values of $X_j$ are
conjugate under $\Sp(V)$, and that transformation preserves the law of the independent uniform
$X_i$. So the leaf block has its unconditional law, independently of all other frames and blocks.
Remove the leaf and iterate.
\end{proof}

For three frames every pair of the three blocks is a forest, so the three blocks are pairwise
independent. Cycles are treated below.

\begin{theorem}[Dependence invisible in every proper frame marginal; \stU]\label{spn:thm:hierarchy}
Fix $q$, $k\ge3$ and $2/k<\alpha<2/(k-1)$, and let all $k$ frames be isotropic of dimension
$d=\floor{\alpha n}$. For large $n$ the frames are admissible, there is $\varepsilon>0$ such that the
complete cross pairings of every proper subset of the frames are within $O_{q,k,\alpha}(q^{-\varepsilon n})$
of independent uniform entries, and the full observation is at distance
$1-O_{q,k,\alpha}(q^{-\varepsilon n})$ from that law and from the product of its own block marginals.
\end{theorem}

\begin{proof}
Here $\alpha<1$, so $d\le n$. A proper subset of $j\ge2$ frames has total dimension at most
$(k-1)d$, is balanced, and has gap at least $[2-(k-1)\alpha]n$, so
\cref{spn:thm:joint-threshold} bounds its distance by $\frac{q+1+j}{q^2-1}q^{-[2-(k-1)\alpha]n}$;
a single frame has no cross blocks. The full observation has $kd-2n\ge(k\alpha-2)n-k$, and the
supercritical bound applies. Any $\varepsilon$ below $2-(k-1)\alpha$ and $k\alpha-2$ works. Each block
marginal is exponentially close to uniform, the product of the actual marginals is within the sum of
the $\binom k2$ marginal distances of the uniform product, and the triangle inequality preserves the
separation.
\end{proof}

Proper frame marginals discard all blocks incident to a removed frame; deleting a single block from
the full observation need not remove its distinguishing power.

\begin{example}[Three almost uniform binary blocks; \stU]\label{spn:ex:three-blocks}
Let $q=2$, $n=50$, and let $X_1,X_2,X_3$ be independent uniform isotropic $40$-frames in $\F_2^{100}$.
The three blocks $M_{12},M_{13},M_{23}$ are exactly pairwise independent, each is within
$\frac53\,2^{-20}<1.59\cdot10^{-6}$ of a uniform $40\times40$ binary matrix, and the triple is at
distance greater than $0.99999809$ from three independent uniform matrices and greater than
$0.99999332$ from the product of its own marginals. The separating event is
$\rank\Gamma\le100$ for the assembled $120\times120$ matrix, which always holds for the blocks and has
probability below $0.00000191$ for independent uniform blocks; with equal priors the rank test
succeeds with probability above $0.99999904$.
\end{example}

\begin{proof}
For two frames the isotropic form of \cref{spn:lem:joint-moment} gives
$A_2=5-6\cdot2^{-40}+2\cdot2^{-80}<5$ with gap $20$, which proves the marginal bound by
\cref{spn:thm:joint-threshold}(1). For three frames $A_3=3+3\cdot2^{-40}-9\cdot2^{-80}+4\cdot2^{-120}$ and
$3d-2n=20$, so the distance exceeds
$1-(2+3\cdot2^{-40}-9\cdot2^{-80}+4\cdot2^{-120})/(2^{20}-1)>0.99999809$. Pairwise independence is
\cref{spn:thm:forest}, and if $P$ is the product of the actual block laws,
$\spnTV(\mu,P)\ge\spnTV(\mu,\nu)-3A_2/(3\cdot2^{20})>0.99999332$.
\end{proof}

In the smallest case, three independent nonzero vectors of $\F_2^2$ have pairings $000$ with
probability $1/9$ and $011,101,110,111$ with probability $2/9$ each: each bit is Bernoulli with
parameter $2/3$, any two are independent, and the distance from three fair bits is $7/18$. The
enumeration recorded after \cref{spn:thm:joint-support}, with exact distances
$25/32$, $2/9$ and $179/243$ for four vectors over $\F_2$ in dimension two, three vectors over $\F_3$
in dimension two, and a nondegenerate plane with two vectors over $\F_3$ in dimension two, confirms the moment and both
bounds, and all $83$ nontrivial forest marginals checked are
independent.\footnote{\raggedright\artifact{symplectic_joint_observations_2026-11-27/results/exhaustive_checks.json},
\artifact{symplectic_joint_observations_2026-11-27/results/independent_moment_results.json},
\artifact{symplectic_joint_observations_2026-11-27/results/headline_example.json},
\artifact{symplectic_joint_observations_2026-11-27/results/certified_envelopes.csv}}
Finite-field rank laws and total variation are studied in \cite{Fulman2015Stein}, a symmetric-Gram
comparison over odd fields is \cite[Proposition~5.2]{Li2026Average}, finite symplectic matrices occur
in random Clifford circuits \cite{Farshi2022Mixing}, and a different random walk on symplectic forms
has a cutoff \cite{He2020Random}.

\paragraph{Cycles of vector observations.}
The simplest frames are single vectors: an isometric injection of a line with the zero form is a
nonzero vector. Let $\mathcal G$ be a finite simple graph, choose independent uniform nonzero vectors
$X_v\in V$ at its vertices, and observe $Y_{uv}=\langle X_u,X_v\rangle$ along oriented edges. Put
$N=q^{2n}$. One edge has the law
\[
 p_N(0)=\frac{N-q}{q(N-1)},\qquad p_N(a)=\frac N{q(N-1)}\quad(a\ne0),
\]
positive since $N\ge q^2$; reversing an edge negates its label, which preserves $p_N$ and every
distance below. The comparison law is now the product $Q_{\mathcal G}=p_N^{\otimes E(\mathcal G)}$ of
the true one-edge marginals, which removes the bias already visible on single edges. For a cyclically
oriented cycle of length $\ell\ge3$ let $\mu_{\ell,N}$ be its edge law and
$\pi_{\ell,N}=p_N^{\otimes\ell}$, and put $\kappa(0)=q-1$, $\kappa(a)=-1$ for $a\ne0$. Fix a nontrivial
additive character $\psi$ of $\F_q$ and write $\widehat\mu(t)=\spnE\psi(\sum_it_iY_i)$.

\begin{theorem}[Complete cycle Fourier law; \stU]\label{spn:thm:cycle-fourier}
If some $t_i$ vanishes, then
$\widehat\mu_{\ell,N}(t)=(-1/(N-1))^{\#\set{i:t_i\ne0}}=\widehat\pi_{\ell,N}(t)$. If all $t_i\ne0$, then
\[
 \widehat\mu_{2k+1,N}(t)=\frac{N^{k+1}-1}{(N-1)^{2k+1}},\qquad
 \widehat\mu_{2k,N}(t)=\frac{1+N^k\bigl((N-1)\mathbf 1_{\set{t_1t_3\cdots t_{2k-1}=t_2t_4\cdots t_{2k}}}-1\bigr)}{(N-1)^{2k}} .
\]
\end{theorem}

\begin{proof}
Zero frequencies remove edges and leave a forest, whose labels are independent by
\cref{spn:thm:forest}. For the full cycle define matrices indexed by $V\setminus\set0$,
$F_t(x,y)=\psi(t\langle x,y\rangle)$, $P_c(x,z)=\mathbf 1_{\set{z=cx}}$ and $\mathbf J(x,z)=1$. For
$t\ne0$, additive orthogonality gives $F_t\mathbf 1=-\mathbf 1$ and $F_sF_t=NP_{s/t}-\mathbf J$. The
constants and their orthogonal complement $\mathcal K$ are invariant. On $\mathcal K$, pairing the
factors in the cycle trace, an even cycle gives $N^kP_c$ with
$c=t_1t_3\cdots t_{2k-1}/(t_2t_4\cdots t_{2k})$, whose trace on $\mathcal K$ is
$(N-1)\mathbf 1_{\set{c=1}}-1$; an odd cycle gives $N^kP_cF_t$, every diagonal entry of $P_cF_t$ is one
because the form is alternating, so its trace is $N-1$, its constant eigenvalue is $-1$, and its
trace on $\mathcal K$ is $N$. Adding the constant eigenvalue $(-1)^\ell$ and dividing by
$(N-1)^\ell$ gives the formulas, with no division by two.
\end{proof}

\begin{theorem}[Exact pointwise cycle law; \stU]\label{spn:thm:cycle-law}
For odd $\ell=2k+1$,
\begin{equation}\label{spn:eq:cycle-odd}
 \mu_{\ell,N}(a)-\pi_{\ell,N}(a)=\frac{N^{k+1}}{q^\ell(N-1)^\ell}\prod_{i=1}^\ell\kappa(a_i).
\end{equation}
For even $\ell=2k$,
\begin{equation}\label{spn:eq:cycle-even}
 \mu_{\ell,N}(a)-\pi_{\ell,N}(a)=\frac{N^k}{q^\ell(N-1)^\ell}\Bigl((N-1)K_\ell(a)-\prod_i\kappa(a_i)\Bigr),
\end{equation}
where $K_\ell(a)=\sum\psi(-\sum_it_ia_i)$ over $t\in(\F_q^\times)^\ell$ with
$t_1t_3\cdots t_{\ell-1}=t_2t_4\cdots t_\ell$ has the evaluation
\[
 K_\ell(a)=\begin{cases}
 \dfrac{\prod_i\kappa(a_i)}{q-1},&\text{some }a_i=0,\\[8pt]
 \dfrac{1+q^k\bigl((q-1)\mathbf 1_{\mathcal R(a)}-1\bigr)}{q-1},&\text{all }a_i\ne0,
 \end{cases}
\]
where $\mathcal R(a)$ is the multiplicative resonance $a_1a_3\cdots a_{2k-1}=(-1)^ka_2a_4\cdots a_{2k}$.
\end{theorem}

\begin{proof}
By Fourier inversion and \cref{spn:thm:cycle-fourier}, only frequencies with all $t_i\ne0$
contribute to $\mu-\pi$. For odd $\ell$ the Fourier difference is then constant and
$\sum_{t\ne0}\psi(-ta)=\kappa(a)$ gives \eqref{spn:eq:cycle-odd}; for even $\ell$, separating the
balanced-frequency indicator from the constant gives \eqref{spn:eq:cycle-even}. For $K_\ell$, insert
multiplicative-character orthogonality for the constraint. The trivial character contributes
$\prod_i\kappa(a_i)/(q-1)$, and if some $a_i=0$ every other term vanishes. If all $a_i\ne0$, the
nontrivial characters contribute $k$ copies of each of the Gauss sums
$G(\chi)=\sum_{x\ne0}\chi(x)\psi(x)$ and $G(\chi^{-1})$, and $G(\chi)G(\chi^{-1})=\chi(-1)q$, which
follows by putting $x=uy$ and summing over $y$; so the contribution is
$\frac{q^k}{q-1}\sum_{\chi\ne1}\chi\bigl((-1)^ka_2a_4\cdots a_{2k}/(a_1a_3\cdots a_{2k-1})\bigr)$, and
the character sum is $q-2$ on $\mathcal R(a)$ and $-1$ otherwise.
\end{proof}

\begin{theorem}[Exact one-cycle distances; \stU]\label{spn:thm:cycle-distances}
For odd $\ell=2k+1$,
\[
 \spnTV(\mu_{\ell,N},\pi_{\ell,N})=2^{\ell-1}(1-q^{-1})^\ell\frac{N^{k+1}}{(N-1)^\ell},\qquad
 \chi^2(\mu_{\ell,N}\Vert\pi_{\ell,N})=(q-1)^\ell\frac N{(N-q)^\ell}.
\]
For even $\ell=2k$,
\begin{gather*}
 \spnTV(\mu_{\ell,N},\pi_{\ell,N})=\frac{N^k(q-1)^{\ell-2}}{q^\ell(N-1)^\ell}
 \Bigl[(N-q)\bigl((q-1)2^{\ell-1}-(q-2)\bigr)+(N-1)(q-2)q^k\Bigr],\\
 \chi^2(\mu_{\ell,N}\Vert\pi_{\ell,N})=(q-1)^{\ell-2}\bigl((N-q)^{2-\ell}+(q-2)(N-1)^{2-\ell}\bigr).
\end{gather*}
\end{theorem}

\begin{proof}
Odd case: \eqref{spn:eq:cycle-odd} and $\sum_a\abs{\kappa(a)}=2(q-1)$ give the distance. With
$b_N(0)=(q-1)N/(N-q)$ and $b_N(a)=-1$ for $a\ne0$ one has $\spnE_{p_N}b_N=0$,
$\spnE_{p_N}b_N^2=(q-1)N/(N-q)$ and $\mu/\pi-1=N^{-k}\prod_ib_N(a_i)$; square and factor.

Even case: write $A_N=N-q$, $C_N=(N-1)q^k$ and $B(a)=(N-1)K_\ell(a)-\prod_i\kappa(a_i)$. If $z\ge1$
coordinates of $a$ vanish, $B(a)=A_N\prod_i\kappa(a_i)/(q-1)$, of sign $(-1)^z$. If none vanishes,
$B(a)=(A_N+(q-2)C_N)/(q-1)$ on $\mathcal R(a)$ and $(A_N-C_N)/(q-1)$ otherwise, positive and negative
respectively when the class is nonempty; there are $(q-1)^{\ell-1}$ resonant and
$(q-2)(q-1)^{\ell-1}$ nonresonant vectors. The vectors with a zero coordinate give
$\sum\abs B=A_N(q-1)^{\ell-1}(2^\ell-1)$, the others $(q-1)^{\ell-2}((3-q)A_N+2(q-2)C_N)$, and
multiplication by $N^k/(2q^\ell(N-1)^\ell)$ gives the distance. For the chi-square,
$\chi^2=\frac1{q^\ell(N-1)^\ell}\sum_a(N/(N-q))^zB(a)^2$; the terms with $z\ge1$ sum to
$(q-1)^{\ell-2}A_N^2[(q(N-1)/A_N)^\ell-1]$ and those with $z=0$ to $(q-1)^{\ell-2}(A_N^2+(q-2)C_N^2)$,
and the $A_N^2$ terms cancel, with $q^{2k}=q^\ell$.
\end{proof}

A cactus is a finite simple graph in which two distinct simple cycles share at most one vertex; its
blocks with edges are bridges and simple cycles.

\begin{theorem}[Cactus factorization; \stU]\label{spn:thm:cactus}
For a cactus, the edge vectors of its blocks are mutually independent; each bridge has law $p_N$
and each cyclically oriented cycle of length $\ell$ the law $\mu_{\ell,N}$.
\end{theorem}

\begin{proof}
For a block meeting the rest in an articulation vertex $v$, condition on all vectors outside the
block other than $X_v$. The other vectors of the block are independent and uniform, and the
conditional law of its labels given $X_v=x$ does not depend on $x\ne0$, because $\Sp(V)$ is transitive
on nonzero vectors, preserves all labels and carries uniform vectors to uniform vectors. So the block
is independent of the rest. Remove leaf blocks of the block-cut tree and iterate, componentwise.
\end{proof}

\begin{corollary}[Independence characterizes forests; \stU]\label{spn:cor:forest-only}
The edge labels of a finite simple graph are independent if and only if the graph is a forest.
\end{corollary}

\begin{proof}
Forests are \cref{spn:thm:forest}. A cycle has an edge marginal at the positive distance of
\cref{spn:thm:cycle-distances}.
\end{proof}

\begin{corollary}[Paired cycle lengths carry equal information; \stU]\label{spn:cor:paired}
For fixed $q$ and $r\ge1$ both lengths $\ell=2r+1$ and $\ell=2r+2$ satisfy
\[
 v_{\ell,N}:=\chi^2(\mu_{\ell,N}\Vert\pi_{\ell,N})=(q-1)^{2r+1}N^{-2r}\bigl(1+O_{q,r}(N^{-1})\bigr),
\]
and uniformly in $a$, $\mu_{\ell,N}(a)/\pi_{\ell,N}(a)-1=O_{q,\ell}(N^{-r})$.
\end{corollary}

\begin{proof}
Expand \cref{spn:thm:cycle-distances}; in the even case the leading factor is
$(q-1)^{2r}(1+q-2)$. The uniform bound follows from \eqref{spn:eq:cycle-odd},
\eqref{spn:eq:cycle-even} and $p_N\to1/q$.
\end{proof}

Adding one edge to an odd cycle therefore leaves its leading information unchanged, although the two
laws differ. A cycle reveals dependence only through all of its edges: every proper subset of its
edges is a forest.

\begin{lemma}[Many small independent likelihood increments; \stU]\label{spn:lem:product-experiments}
For each $n$ let $Q_n=\bigotimes_jQ_{n,j}$ and $P_n=\bigotimes_jP_{n,j}$ be finite products, write
$dP_{n,j}/dQ_{n,j}=1+u_{n,j}$, $v_{n,j}=\spnE_{Q_{n,j}}u_{n,j}^2$, $I_n=\sum_jv_{n,j}$ and
$\delta_n=\max_j\norm{u_{n,j}}_\infty$, and assume $\delta_n\to0$. If $I_n\to\tau\in[0,\infty)$,
then $\spnTV(P_n,Q_n)\to2\Phi(\sqrt\tau/2)-1$, with $\Phi$ the standard normal distribution
function; if $I_n\to\infty$, then $\spnTV(P_n,Q_n)\to1$; and if $0<\tau<\infty$, the log-likelihood
ratio converges in law to $\mathcal N(-\tau/2,\tau)$ under $Q_n$ and to $\mathcal N(\tau/2,\tau)$
under $P_n$.
\end{lemma}

\begin{proof}
Each $u_{n,j}$ has mean zero under $Q_{n,j}$. For $\abs u\le\delta<1/2$,
$\log(1+u)=u-\frac12u^2+O(\delta u^2)$ and $\abs{\log(1+u)}\le2\abs u$. For
$Z_n=\sum_j\log(1+u_{n,j})$ this gives $\spnE_{Q_n}Z_n=-\frac12I_n+O(\delta_nI_n)$ and
$\operatorname{Var}_{Q_n}Z_n=I_n+O(\delta_nI_n)$, the squared means summing to
$O(\delta_n^2I_n)$, and the centred third absolute moments sum to $O(\delta_nI_n)$, so the Lyapunov
central limit theorem gives $Z_n\Rightarrow\mathcal N(-\tau/2,\tau)$ under $Q_n$. Under $P_{n,j}$,
$\spnE_{P_{n,j}}f=\spnE_{Q_{n,j}}(1+u_{n,j})f$, and the same estimates give mean $v_{n,j}/2+O(\delta_nv_{n,j})$,
variance $v_{n,j}+O(\delta_nv_{n,j})$ and the same third-moment control. Since
$\spnTV(P_n,Q_n)=\spnE_{Q_n}(1-e^{Z_n})_+$ with a bounded continuous integrand, and
$\spnE(1-e^Z)_+=\Phi(\sqrt\tau/2)-\Phi(-\sqrt\tau/2)$ for $Z\sim\mathcal N(-\tau/2,\tau)$, the
distance converges as claimed. If $I_n\to0$, Cauchy--Schwarz gives
$\spnTV\le\frac12(\prod_j(1+v_{n,j})-1)^{1/2}\le\frac12(e^{I_n}-1)^{1/2}\to0$. If $I_n\to\infty$, the
Hellinger affinity $A_n=\spnE_{Q_n}(dP_n/dQ_n)^{1/2}$ satisfies, by Taylor's theorem and
$-\frac{d^2}{du^2}\sqrt{1+u}\ge1/(4(1+\delta)^{3/2})$,
$A_n\le\exp(-I_n/(8(1+\delta_n)^{3/2}))\to0$, and $1-\spnTV(P_n,Q_n)\le A_n$.
\end{proof}

\begin{theorem}[Sharp detection on growing cactus graphs; \stU]\label{spn:thm:cactus-detection}
Fix $q$ and $L\ge3$. For each $n$ let $\mathcal G_n$ be any cactus with cycle lengths in
$\set{3,\dots,L}$, any number of vertices and bridges, set of cycles $\mathcal C_n$, law $P_n$ and
product of true marginals $Q_n$, and put $I_n=\sum_{C\in\mathcal C_n}v_{\abs C,q^{2n}}$. If
$I_n\to\tau\in[0,\infty)$, then $\spnTV(P_n,Q_n)\to2\Phi(\sqrt\tau/2)-1$; if $I_n\to\infty$, the
distance tends to one; and for $0<\tau<\infty$ the log-likelihood ratio is asymptotically
$\mathcal N(\mp\tau/2,\tau)$ under $Q_n$ and $P_n$, and the least equal-prior error tends to
$\Phi(-\sqrt\tau/2)$. For finite parameters, with $\delta=\max_C\norm{\mu_{\abs C,N}/\pi_{\abs C,N}-1}_\infty<1$
and $I=\sum_Cv_{\abs C,N}$,
\begin{equation}\label{spn:eq:cactus-finite}
 1-\exp\Bigl(-\frac I{8(1+\delta)^{3/2}}\Bigr)\le\spnTV(P_{\mathcal G},Q_{\mathcal G})
 \le\min\Bigl\{1,\tfrac12\sqrt{e^I-1}\Bigr\}.
\end{equation}
\end{theorem}

\begin{proof}
By \cref{spn:thm:cactus},
$dP_{\mathcal G}/dQ_{\mathcal G}=\prod_C\mu_{\abs C,N}(y_C)/\pi_{\abs C,N}(y_C)$ and
$1+\chi^2(P_{\mathcal G}\Vert Q_{\mathcal G})=\prod_C(1+v_{\abs C,N})$, the bridges having identical
factors. \Cref{spn:cor:paired} gives $\delta_n=O_{q,L}(N^{-1})$; apply
\cref{spn:lem:product-experiments}. The finite bounds are the last two estimates of its proof, and
the least equal-prior error is $(1-\spnTV)/2$.
\end{proof}

Consequently a cactus with $m_n$ cycles, all of lengths $2r+1$ or $2r+2$ in any mixture, has critical
scale $m_n\asymp q^{4nr}/(q-1)^{2r+1}$, and $m_n(q-1)^{2r+1}q^{-4nr}\to\tau$ gives the profile. Also,
when $I_n=O(1)$, both relative entropies $\mathrm{KL}(P_n\Vert Q_n)$ and $\mathrm{KL}(Q_n\Vert P_n)$ are
$\frac12I_n+o(1)$, all cycle likelihoods are eventually positive, and the likelihood-ratio test has
both error probabilities tending to $\Phi(-\sqrt\tau/2)$, by the two expectations of $Z_n$ in the
proof of \cref{spn:lem:product-experiments}.

\begin{proposition}[An asymptotically optimal reduced score; \stU]\label{spn:prop:cactus-score}
Put $r(\ell)=\floor{(\ell-1)/2}$, $g_\ell(a)=\prod_i\kappa(a_i)$ for odd $\ell$ and $g_\ell=K_\ell$ for
even $\ell$. Uniformly in $a$, $\mu_{\ell,N}/\pi_{\ell,N}-1=N^{-r(\ell)}g_\ell+O_{q,\ell}(N^{-r(\ell)-1})$,
and under uniform labels $g_\ell$ has mean zero and variance $(q-1)^{2r(\ell)+1}$. If $I_n\to\tau\in(0,\infty)$
in \cref{spn:thm:cactus-detection} and
$T_n=\sum_CN^{-r(\abs C)}(g_{\abs C}(Y_C)-\spnE_{\pi_{\abs C,N}}g_{\abs C})$, then
$\log(dP_n/dQ_n)=T_n-\frac12I_n+o(1)$ in probability under $Q_n$ and under $P_n$, so the test
$T_n>I_n/2$ is asymptotically optimal.
\end{proposition}

\begin{proof}
The expansion is \cref{spn:cor:paired} with the leading terms of \eqref{spn:eq:cycle-odd},
\eqref{spn:eq:cycle-even}. For odd $\ell$ the variance factors as $(\spnE\kappa^2)^\ell$ with
$\spnE\kappa^2=q-1$; for even $\ell$, Parseval counts the $(q-1)^{\ell-1}$ nonzero frequency vectors
satisfying the product constraint. Let $u_C=\mu/\pi-1$ and $w_C$ its centred leading term. Then
$\spnE_{Q_n}(u_C-w_C)=0$ and $\spnE_{Q_n}(u_C-w_C)^2=O_{q,L}(N^{-2}v_{\abs C,N})$, so by independence
$\sum_C(u_C-w_C)=o_{Q_n}(1)$; also $\sum_Cu_C^2-I_n=o_{Q_n}(1)$, its variance being at most
$\delta_n^2I_n$, and the Taylor remainder has expected absolute value $O(\delta_nI_n)$. Under $P_n$
the mean of the first difference is $\sum_C\spnE_Qu_C(u_C-w_C)=O_{q,L}(I_n/N)$ by Cauchy--Schwarz, its
variance tends to zero by the same estimates and $1+u_C\le1+\delta_n$, and the other two terms are
controlled in the same way.
\end{proof}

\begin{remark}[Over $\F_2$ cycle parity suffices]\label{spn:rem:binary-parity}
For $q=2$ both scores reduce to $g_\ell(a)=(-1)^{a_1+\cdots+a_\ell}$, the single nonzero frequency
vector satisfying the product constraint. So one parity bit per cycle, weighted by $N^{-r(\ell)}$,
is asymptotically sufficient, and for cycles of one adjacent pair of lengths the number of
even-parity cycles suffices after centring. For triangles or squares over $\F_2$ the critical scale
is $m\asymp2^{4n}$; at $m/2^{4n}\to1$ the limiting distance is $2\Phi(1/2)-1=0.3829249225\ldots$
and the least equal-prior error $0.3085375387\ldots$, and at $m/2^{4n}\to4$ the distance is
$2\Phi(1)-1=0.6826894921\ldots$, while every tree-shaped subset of the observations remains exactly
independent.
\end{remark}

The likelihood ratio of a cycle depends only on the number $z$ of zero labels and, for even
cycles without zeros, on the resonance bit, so each cycle reduces exactly to at most $\ell+2$
outcomes, which is how $I$ and $\delta$ are computed in rational arithmetic. Enumeration of
$1{,}745{,}891$ nonzero-vector assignments in $16$ cycle, four cactus and five tree cases over
$\F_2,\F_3,\F_4,\F_5$ in dimensions two and four checks all $1{,}976$ cycle label probabilities,
their distances and divergences, and $11{,}614$ marginal equalities on $608$ proper edge subsets;
$432$ further cases with $q$ up to $25$ check the grouped likelihoods against the closed
formulas.\footnote{\raggedright\artifact{symplectic_dependence_research_2026-11-27/results/independent_enumeration.json},
\artifact{symplectic_dependence_research_2026-11-27/results/grouped_law_checks.json},
\artifact{symplectic_dependence_research_2026-11-27/results/cycle_information.csv},
\artifact{symplectic_dependence_research_2026-11-27/results/example_q4_cycle4.json}}
Cycle counts in finite-field distance and dot-product graphs \cite{IosevichJardineMcDonald2021Cycles},
general nondegenerate bilinear forms \cite{Bright2023Generalized} and symplectic graphs
\cite{TangWan2006Symplectic} are background, and observation-graph-dependent detectability has a
counterpart for real Wishart matrices \cite{BrennanBreslerHuang2021Finetti}.

\paragraph{Overlapping cycles.}
The cactus hypothesis is essential. Let $2\le a\le2n$, assign independent uniform nonzero vectors
$X_1,\dots,X_a$ to one side of the complete bipartite graph $K_{a,m}$ and $Z_1,\dots,Z_m$ to the
other, and observe the columns $Y_j=(\langle X_i,Z_j\rangle)_{i\le a}\in\F_q^a$. Let $P_{a,m,N}$ be
their law and $Q_{a,m,N}=(p_N^{\otimes a})^{\otimes m}$. Define the probability law
$\rho_{a,N}(0)=(N-q^a)/(q^a(N-1))$, $\rho_{a,N}(y)=N/(q^a(N-1))$ for $y\ne0$, and put
$\alpha_{a,N}=\prod_{i=0}^{a-1}(N-q^i)/(N-1)$, $\delta_{a,N}=1-\alpha_{a,N}$ and
$J_{a,q}=q^a-1-a(q-1)>0$.

\begin{proposition}[Independent anchors; \stU]\label{spn:prop:anchors}
The anchors $X_1,\dots,X_a$ are independent with probability $\alpha_{a,N}$, and conditional on this
the columns are independent with law $\rho_{a,N}$. Hence, for every $m$,
\[
 \spnTV(P_{a,m,N},\rho_{a,N}^{\otimes m})\le\delta_{a,N}\le\frac{(q^a-1)/(q-1)-a}{N-1}.
\]
Moreover, uniformly in $y\in\F_q^a$,
$\rho_{a,N}(y)/p_N^{\otimes a}(y)=1+h_a(y)/N+O_{a,q}(N^{-2})$ with
$h_a(y)=\sum_i\kappa(y_i)-(q^a\mathbf 1_{\set{y=0}}-1)$, which under uniform coordinates has mean zero and
second moment $J_{a,q}$; so $\chi^2(\rho_{a,N}\Vert p_N^{\otimes a})=J_{a,q}N^{-2}+O_{a,q}(N^{-3})$.
\end{proposition}

\begin{proof}
After $i$ independent anchors, $q^i-1$ of the $N-1$ candidates lie in their span, which gives
$\alpha_{a,N}$, and $1-\prod_i(1-b_i)\le\sum_ib_i$ gives the bound. For independent anchors the map
$z\mapsto(\langle X_i,z\rangle)_i$ is onto $\F_q^a$ with fibres of size $N/q^a$, and removing $z=0$
gives $\rho_{a,N}$, whatever the anchors; so
$P_{a,m,N}=\alpha_{a,N}\rho_{a,N}^{\otimes m}+\delta_{a,N}\mathcal D$ for a probability law $\mathcal D$.
For $y\ne0$ with $z$ zero coordinates the likelihood is $(1-N^{-1})^{a-1}(1-q/N)^{-z}$, and at $y=0$ it
is $(1-q^a/N)(1-N^{-1})^{a-1}(1-q/N)^{-a}$; expanding gives $h_a$. Since
$q^a\mathbf 1_{\set{y=0}}=\prod_i(1+\kappa(y_i))$, $h_a=-\sum_{\abs S\ge2}\prod_{i\in S}\kappa(y_i)$, an
orthogonal sum with squared norms $\sum_{s\ge2}\binom as(q-1)^s=J_{a,q}$; the reference law differs
from uniform by $O_{a,q}(N^{-1})$ pointwise.
\end{proof}

\begin{theorem}[The fixed-anchor critical profile; \stU]\label{spn:thm:anchors}
Fix $q$ and $a\ge2$, let $n\to\infty$ and let $m=m_n$ be arbitrary. If $mJ_{a,q}/N^2\to\tau\in[0,\infty]$,
then $\spnTV(P_{a,m,N},Q_{a,m,N})\to2\Phi(\sqrt\tau/2)-1$, read as one when $\tau=\infty$.
\end{theorem}

\begin{proof}
Apply \cref{spn:lem:product-experiments} to $\rho_{a,N}^{\otimes m}$ and $(p_N^{\otimes a})^{\otimes m}$
with \cref{spn:prop:anchors}; the actual distance differs by at most $\delta_{a,N}=O_{a,q}(N^{-1})$,
uniformly in $m$.
\end{proof}

For $a=2$ the graph $K_{2,m}$ has $\binom m2$ four-cycles, each with information
$v_{4,N}\sim(q-1)^3N^{-2}$, and $\delta_{2,N}=(q-1)/(N-1)$,
$\chi^2(\rho_{2,N}\Vert p_N^{\otimes2})=(q-1)^2/(N-q)^2$, so
\[
 \spnTV(P_{2,m,N},Q_{2,m,N})\le\frac{q-1}{N-1}+\frac12\Bigl(\Bigl(1+\frac{(q-1)^2}{(N-q)^2}\Bigr)^m-1\Bigr)^{1/2}.
\]
For fixed $1<\beta<2$ and $m=\floor{N^\beta}$ this tends to zero, while
$\binom m2v_{4,N}\sim\frac12(q-1)^3N^{2\beta-2}\to\infty$: adding the information of overlapping cycle
marginals predicts the wrong threshold by a power of $N$. The full chi-square divergence even
diverges: for $c\in\F_q^\times$ the event $Y_{2j}=cY_{1j}$ for all $j$ has probability at least
$1/(N-1)$ under $P_{2,m,N}$, forced by $X_2=cX_1$, and $s_N^m$ under $Q_{2,m,N}$, with
$s_N=1/q+(q-1)/(q(N-1)^2)$, so Cauchy--Schwarz on this event gives
$1+\chi^2(P_{2,m,N}\Vert Q_{2,m,N})\ge1/((N-1)^2s_N^m)$, whose logarithm is at least
$m\log q-2\log(N-1)-O_q(m/N^2)\to\infty$. Rare dependent anchors carry negligible probability but
very large likelihood ratios, so a direct second-moment bound for the mixture is poor, whereas
conditioning on independent anchors gives the correct profile. Twenty complete bipartite cases
enumerating $4{,}094{,}368$ assignments check $7{,}382$ label probabilities and sixteen limiting score
variances.\footnote{\artifact{symplectic_dependence_research_2026-11-27/results/anchor_verification.json}}

%======================================================================

\subsection{Symplectic frames at every congruence precision}\label{spn:sec:padic}
\sources{padic_symplectic_corners; symplectic_compression_2026-11-26; symplectic_compression_2026-11-26_2; exact_symplectic_compressions_2026-11-26; prime_endpoints_and_symplectic_compressions; grh_extensions_2026-11-27_3}

A finite-field compression is the first digit of a compression over a local ring, and local
arithmetic can see the further digits. When the source is a single symplectic pair and the target
a symplectic summand, the law is exact at every precision, and when both frames consist of several
symplectic pairs in a stable range, the distance and the optimal test are the same at every precision. In this subsection $\OO$ is a complete
discrete valuation ring with uniformizer $\varpi$ and finite residue field $\kappa$ of order $q$; its
fraction field may have characteristic zero or $p$, and $q$ may be even. Put
$R_e=\OO/\varpi^e\OO$ for $e\ge1$ and $R_\infty=\OO$, normalize Haar measures to mass one, and
for $t\in R_e$ let $v_e(t)=\max\set{0\le r\le e:t\in\varpi^rR_e}$, so $v_e(0)=e$.

Let $J$ be a unimodular alternating matrix of size $2n$, fix $1\le a\le n$ and
$Y\in\spnMat_{2n,2a}(\OO)$, $X\in\spnMat_{2n,2}(\OO)$ with $Y^{\spnT}JY=J_a$ and
$X^{\spnT}JX=J_1$; these conditions already make their columns span direct summands. For $g$
Haar distributed in $\Sp(J,R_e)$ let $\mu_e$ be the law of
\[
 M=Y^{\spnT}JgX\in\spnMat_{2a,2}(R_e),
\]
let $\nu_e$ be the law of $4a$ independent uniform entries, and for columns $M_1,M_2$ put
$\Omega(M)=M_1^{\spnT}J_aM_2$, the symplectic area; for $a=1$, $\Omega(M)=\det M$. Write
\[
 N=2n-2a,\qquad d_n=1-q^{-2n},\qquad \alpha=\frac{1-q^{-N}}{d_n},
\]
\[
 f_{N,e}(t)=(1-q^{-N})\sum_{r=0}^{\min(v_e(t),e-1)}q^{-(N-1)r}+q^{-(N-1)e}\mathbf 1_{\set{t=0}} .
\]
For $e>k\ge1$, with $e=\infty$ allowed, let $\widetilde\mu_{k,e}$ be the law obtained by sampling
modulo $\varpi^k$ from $\mu_k$ and filling in all further digits independently and uniformly.

\begin{theorem}[Exact laws at every precision; \stU]\label{spn:thm:padic}
Let $a<n$.
\begin{enumerate}
\item The law $\mu_e$ does not depend on the chosen frames. For finite $e$ its likelihood
relative to $\nu_e$ is $f_{N,e}(1-\Omega(M))/d_n$. Over $\OO$, $\mu_\infty$ is absolutely
continuous with respect to $\nu_\infty$, with continuous density
\[
 h(M)=\alpha\sum_{r=0}^{v(1-\Omega(M))}q^{-(N-1)r},
\]
a convergent geometric series when $\Omega(M)=1$. Only the valuation of $1-\Omega(M)$ matters.
\item For every $e\in\set{1,2,\dots,\infty}$,
\[
 \spnTV(\mu_e,\nu_e)=\bigl(1-q^{-1}+q^{-(2a+1)}\bigr)\frac{q^{2a}-1}{q^{2n}-1},
\]
and $\set{\Omega(M)\equiv1\pmod\varpi}$ is an optimal distinguishing event at every precision.
Observing all digits does not improve the optimal equal-prior success probability over
observing the residues.
\item For $e>k\ge1$, with $e=\infty$ allowed,
\[
 \spnTV(\mu_e,\widetilde\mu_{k,e})=(1-q^{-1})(1-\alpha)q^{-Nk},
\]
and $\set{1-\Omega(M)\equiv0\pmod{\varpi^{k+1}}}$ is optimal: one digit beyond the retained
precision detects all the bias discarded by uniform completion. For $a=k=1$ the distance is
$\frac{q-1}q\cdot\frac{q^2-1}{(q^{2n}-1)q^{2n-2}}$, and the second factor is the probability that
the reduced planes $\spnim\overline Y$ and $\spnim(\overline g\,\overline X)$ coincide; so the
bias beyond the first digit is of exact order $q^{-4n}$.
\end{enumerate}
\end{theorem}

\begin{proof}
A pair with pairing one extends to a symplectic basis over $R_e$ and over $\OO$: split off its
unimodular plane and apply symplectic Gram--Schmidt in the orthogonal complement, which needs no
division by two. So the symplectic group is transitive on such pairs, and likewise on frames with
Gram matrix $J_a$; this gives independence of the frames, and $gX$ is uniform among the pairs
with pairing one. Choose a symplectic basis beginning with $Y$ and write
$gX=\left(\begin{smallmatrix}B\\C\end{smallmatrix}\right)$ with $B\in\spnMat_{2a,2}(R_e)$ and
$C\in\spnMat_{N,2}(R_e)$. Then $M=J_aB$, $\Omega(M)=B_1^{\spnT}J_aB_2$, and the frame condition
is
\begin{equation}\label{spn:eq:padic-defect}
 \langle C_1,C_2\rangle=1-\Omega(M)
\end{equation}
for the unimodular form on the complement; pairing one gives a left inverse, so there is no
further injectivity condition.

For independent uniform $u,w\in R_e^N$, the least valuation of the coordinates of $u$ is $r<e$
with probability $(1-q^{-N})q^{-Nr}$, and $u=0$ with probability $q^{-Ne}$. When it is $r<e$,
the functional $w\mapsto\langle u,w\rangle$ is uniform on $\varpi^rR_e$, since $u/\varpi^r$ has a
unit coordinate. Hence $\spnP(\langle u,w\rangle=t)=q^{-e}f_{N,e}(t)$, the last term of
$f_{N,e}$ coming from $u=0$. In $R_e^{2n}$ two independent uniform vectors have pairing one with
probability $q^{-e}d_n$, since the first must be primitive and then the pairing is uniform; so
there are $q^{(4n-1)e}d_n$ frames, while for fixed $B$ the solutions of
\eqref{spn:eq:padic-defect} number $q^{(2N-1)e}f_{N,e}(1-\Omega(M))$. As $2N+4a=4n$, the ratio
to $\nu_e(M)=q^{-4ae}$ is $f_{N,e}(1-\Omega(M))/d_n$. The finite densities, pulled back to
$\spnMat_{2a,2}(\OO)$, agree with $h$ off $\set{1-\Omega\in\varpi^e\OO}$ and differ from it by
$O(q^{-(N-1)e})$ there, so they converge uniformly to $h$; this gives the density of the
projective limit and its continuity. Only cardinalities of ideals and unimodularity are used.

Part (2) is the case $b=1$ of \cref{spn:thm:padic-pairs} below, whose range condition is then
$a<n$: there $H(M)=(1-\Omega(M))J_1$, so its event is $E=\set{\Omega\equiv1\pmod\varpi}$, and its
constant is $\alpha=(1-q^{-N})/d_n$. The density is the constant $\alpha$ off $E$, as (1) also shows
directly. Under $\nu_e$ the first column is nonzero modulo $\varpi$ with probability $1-q^{-2a}$, and
then its pairing with the second is uniform, so $\nu_e(E)=\beta_e=q^{-1}(1-q^{-2a})$ and
$\spnTV(\mu_e,\nu_e)=(1-\beta_e)(1-\alpha)$ with $1-\alpha=(q^{2a}-1)/(q^{2n}-1)$. Put
$\epsilon=q^{-N}\le q^{-2}$.

For (3) the two densities agree where $v(1-\Omega)<k$. On the rest the numerator of the density of
$\widetilde\mu_{k,e}$ is $(1-\epsilon)\sum_{r<k}(q\epsilon)^r+(q\epsilon)^k$, over the same
denominator $d_n$. Where $v(1-\Omega)=k$ the true numerator is smaller by $\epsilon(q\epsilon)^k$;
where $1-\Omega\equiv0\pmod{\varpi^{k+1}}$ it is larger, by $(q\epsilon)^k(q\epsilon-\epsilon)$ if
$e=k+1$ and by at least $\epsilon(q\epsilon)^k(q-1-q\epsilon)$ if $e>k+1$. Conditional on a residue
class in $E$, one coordinate of the first column is a unit, and varying the paired coordinate of
the second column makes $\Omega$ uniform in $1+\varpi R_e$; so $\nu_e(v(1-\Omega)=k)
=\beta_e(1-q^{-1})q^{-(k-1)}$. Integrating the negative part gives
$\beta_e(1-q^{-1})q^{-(k-1)}\epsilon(q\epsilon)^k/d_n=(1-q^{-1})(1-\alpha)q^{-Nk}$. For $a=1$ there
are $(q^{2n}-1)q^{2n-2}/(q^2-1)$ nondegenerate planes in $\kappa^{2n}$, by counting pairs with pairing
one and dividing by $\abs{\Sp_2(\kappa)}=q(q^2-1)$, and the symplectic group is transitive on them.
\end{proof}

\begin{remark}[The endpoint $a=n$]\label{spn:rem:padic-endpoint}
The hypothesis $a<n$ is essential: if $a=n$ there is no complement and $\Omega(M)=1$
identically. For $a=n=1$ the compression is uniform on $\SL_2(R_e)$, so
$\spnTV(\mu_e,\nu_e)=1-q^{-e}(1-q^{-2})$, and over $\OO$ the law is singular with respect to
$\nu_\infty$ and the distance is one; at this endpoint higher digits do improve
distinguishability.
\end{remark}

Over the residue field the theorem with $a=1$ recovers, and the endpoint completes, the law of a
compression between two symplectic planes.

\begin{corollary}[Two symplectic planes over a finite field; \stU]\label{spn:thm:two-plane}
In the setting of \cref{spn:sec:complete-law} let $a=b=2$ and
$H_1=H_0=\left(\begin{smallmatrix}0&1\\-1&0\end{smallmatrix}\right)$. For every prime power $q$,
every $n\ge1$ and every $M\in\spnMat_2(\F_q)$,
\[
 q^4\spnP\bigl(M(g)=M\bigr)=
 \begin{cases}
 \dfrac{1+(q^3-q^2)q^{-2n}}{1-q^{-2n}},&\det M=1,\\[8pt]
 \dfrac{1-q^{2}q^{-2n}}{1-q^{-2n}},&\det M\ne1,
 \end{cases}
\]
and
$\spnTV\bigl(M(g),Z\bigr)=(q^3-q^2+1)(q^2-1)/\bigl(q^3(q^{2n}-1)\bigr)$.
\end{corollary}

\begin{proof}
For $n\ge2$ this is \cref{spn:thm:padic} with $a=1$ and $e=1$, since
$f_{N,1}(t)=1-q^{-N}+q^{1-N}\mathbf 1_{\set{t=0}}$ and $q^{-N}=q^2q^{-2n}$. For $n=1$ the
compression is uniform on $\SL_2(\F_q)$, as it must be for two symplectic bases of one plane, and
both formulas reduce to this. The case also follows from \cref{spn:thm:complete-law}: block
elimination gives $\rank\Gamma_M=2+\rank((1-\det M)H_1)$, and when $\det M=1$ one has
$(T_0,T_1,T_2)=(1,q+1,1)$.
\end{proof}

For $q=3$ and $n=2$, each of the $24$ matrices of determinant one has probability $11/720$, each
of the other $57$ has probability $1/90$, and the distance is $19/270$; in particular $I_2$ and
$\operatorname{diag}(1,2)$ have equal rank and different probabilities, and their least
realization dimensions are two and four. For $q>2$ ordinary rank is therefore not a sufficient
statistic, even among invertible matrices. Exact polynomial arithmetic checks
\cref{spn:thm:first-correction,spn:thm:two-plane} in $388$ matrix
cases.\footnote{\artifact{symplectic_compression_2026-11-26/results/asymptotics.json}}

Over the residue field, the rank-only test for a plane source can also be computed exactly, for
every target summand including $a=n$.

\begin{proposition}[Rank information for a plane source; \stU]\label{spn:prop:plane-rank}
Let $e=1$, $1\le a\le n$, and put $x=q^{-2a}$ and $y=q^{-2(n-a)}$. Then
\[
 \spnTV(\mu_1,\nu_1)=\frac{y(1-x)(q-1+x)}{q(1-xy)},\qquad
 \spnTV\bigl(\rank M,\rank Z\bigr)=\frac{yx(1-x)(q+1-qx)}{1-xy},
\]
where $Z$ has law $\nu_1$, and the optimal rank-only event is full column rank. For fixed $q$ and
$n-a=c$, as $a\to\infty$, the first distance tends to $(1-q^{-1})q^{-2c}=\Delta_{q,2,c}$ while the
second is asymptotic to $(q+1)q^{-2n}$.
\end{proposition}

\begin{proof}
At $e=1$ the likelihood is $(1+(q-1)y)/(1-xy)$ on $\set{\Omega=1}$ and $(1-y)/(1-xy)$ elsewhere:
in a symplectic $2c$-space the ordered pairs with a given nonzero pairing number
$(q^{2c}-1)q^{2c-1}$ and those with pairing zero $q^{2c}+(q^{2c}-1)q^{2c-1}$ (read as $0$ and $1$
when $c=0$), and the hidden pairing must be $1-\Omega(M)$. With $\nu_1(\Omega=1)=q^{-1}(1-x)$ this
gives the first distance, which agrees with \cref{spn:thm:padic} when $a<n$. A rank-deficient
two-column matrix has $\Omega=0$ and hence the lower likelihood; under $\nu_1$ it has probability
$1-(1-x)(1-qx)=x(q+1-qx)$. All deficits of the rank law are at ranks zero and one, and the
compensating excess is at rank two, so the rank distance is $x(q+1-qx)[1-(1-y)/(1-xy)]$. The limit
uses \cref{spn:cor:small-source}.
\end{proof}

For $q=2$, $a=10$ and $n=11$ the formulas give a full distance $0.1250000298022158\ldots$ and a
rank distance $7.152547709667702\ldots\cdot10^{-7}$: the pairing test keeps a discrepancy of about
one eighth while every rank-only test has discrepancy below one millionth.

\paragraph{Cokernels and sampling.} Return to $a=1<n$, so that $\Omega=\det$.

\begin{corollary}[All nontrivial cokernels have one correction factor; \stU]\label{spn:cor:padic-cokernel}
For every event $F\subseteq\set{\det M\equiv0\pmod\varpi}$ and every $e\in\set{1,\dots,\infty}$,
$\mu_e(F)=\alpha\,\nu_e(F)$. In particular every nontrivial cokernel isomorphism class, over $R_e$
or over $\OO$, has $\alpha$ times its probability under independent entries, and conditional on a
nontrivial cokernel its law is exactly the independent-entry law. Over $\OO$ the invariant factors
are almost surely $\varpi^i,\varpi^j$ with $0\le i\le j<\infty$; with $u_q=(1-q^{-1})(1-q^{-2})$,
for $(i,j)\ne(0,0)$
\[
 \spnP_{\mu_\infty}(i,j)=\alpha\begin{cases}u_qq^{-4i},&i=j,\\(1-q^{-2})^2q^{-3i-j},&i<j,\end{cases}
\]
and the cokernel is trivial with probability $1-\alpha(1-u_q)$.
\end{corollary}

\begin{proof}
If $\det M\equiv0\pmod\varpi$, then $1-\det M$ is a unit and the density is $\alpha$; over a local
ring a square matrix has trivial cokernel exactly when its determinant is a unit. Under
independent entries all four entries are divisible by $\varpi^i$ with probability $q^{-4i}$, after
which the quotient is again Haar and invertible with probability $u_q$. For $i<j$, divide by
$\varpi^i$; the reduction must have rank one, which happens for $(q^2-1)^2/(q-1)$ residue matrices,
and given one of them, varying the entry opposite a unit pivot makes the determinant uniform in
$\varpi\OO$, so it has valuation exactly $j-i$ with probability $(1-q^{-1})q^{-(j-i-1)}$. The
product is $(1-q^{-2})^2q^{-3i-j}$. These alternatives exhaust a set of full measure.
\end{proof}

\begin{corollary}[Exact sampling with a universal envelope; \stU]\label{spn:cor:padic-sampler}
For every finite $e$, proposing four independent uniform entries and accepting with probability
$\frac58f_{N,e}(1-\det M)/d_n$ samples $\mu_e$ exactly, with $8/5$ proposals on average, whatever
$q$, $n\ge2$ and $e$. The constant $8/5$ is the least envelope valid for all these experiments and
their complete-ring limits.
\end{corollary}

\begin{proof}
All densities are bounded by the complete-ring maximum $(1-\epsilon)/((1-q\epsilon)(1-\epsilon
q^{-2}))$ with $\epsilon=q^{-N}$, which increases with $\epsilon$, since its logarithmic derivative
$-\frac1{1-\epsilon}+\frac q{1-q\epsilon}+\frac{q^{-2}}{1-\epsilon q^{-2}}$ is positive. So $N=2$
is the worst case, with value $(1+q^{-1})/(1-q^{-4})$, which decreases in $q$ and equals $8/5$ at
$q=2$. The likelihood integrates to one, so the acceptance probability is $5/8$; the acceptance
probabilities are rational, so exact random-bit rejection is possible. At $q=2$, $n=2$ the maxima
approach $8/5$ as $e\to\infty$.
\end{proof}

\begin{example}[The $2$-adic integers in dimension four]\label{spn:ex:padic-binary}
For $\OO=\Z_2$ and $n=2$ one has $N=2$, $\alpha=4/5$ and
$h(M)=\frac45\sum_{r=0}^{v_2(1-\det M)}2^{-r}$; the likelihoods at valuations $0,1,2,3$ are
$4/5,6/5,7/5,3/2$, tending to $8/5$. Under independent entries $\det M$ is odd with probability
$3/8$, and under $\mu_e$ with probability $1/2$, so $\spnTV(\mu_e,\nu_e)=1/8$ whether the matrix is
seen modulo $2$, modulo $2^{1000}$ or exactly. Keeping $k$ exact digits and completing uniformly
leaves the error $\spnTV(\mu_\infty,\widetilde\mu_{k,\infty})=1/(10\cdot4^k)$, which is $1/40$ for
$k=1$. The trivial cokernel has probability $3/8$ under independent entries and $1/2$ under
$\mu_\infty$, while for instance $\spnP(\spncoker M\cong\Z_2/(2))=9/40$ and
$\spnP(\spncoker M\cong(\Z_2/(2))^2)=3/160$. Because the complete-ring likelihood increases with the
valuation, the likelihood-threshold tests are exactly the events
$E_r=\set{1-\det M\in\varpi^r\OO}$, with
$\nu_\infty(E_r)=\beta q^{-(r-1)}$ and $\mu_\infty(E_r)=\nu_\infty(E_r)\,f_{N,r}(0)/d_n$; for $q=2$,
$n=2$ the first three points $(\nu_\infty(E_r),\mu_\infty(E_r))$ are $(3/8,1/2)$,
$(3/16,11/40)$ and $(3/32,23/160)$, and only the first realizes the total variation, so the
single-digit test is optimal for equal costs but not for every decision problem. Further values:
\begin{center}\small
\begin{tabular}{cccc}
\toprule
$q$ & $n$ & $\spnTV(\mu_e,\nu_e)$, all $e$ & $\spnTV(\mu_e,\widetilde\mu_{1,e})$, $e\ge2$\\
\midrule
$2$ & $2$ & $1/8$ & $1/40$\\
$2$ & $3$ & $5/168$ & $1/672$\\
$3$ & $2$ & $19/270$ & $1/135$\\
$3$ & $3$ & $19/2457$ & $2/22113$\\
\bottomrule
\end{tabular}
\end{center}
None of these depends on the ramification index or on the characteristic of the fraction field.
\end{example}

\paragraph{Several symplectic pairs.}
Now let the source be a symplectic frame of any size: $X\in\spnMat_{2n,2b}(\OO)$ with
$X^{\spnT}JX=J_b$, $b\ge1$, while $Y^{\spnT}JY=J_a$ as before, and let $\mu_e$ and $\nu_e$ be the law of
$M=Y^{\spnT}JgX$ and the independent-entry law on $\spnMat_{2a,2b}(R_e)$. Put
$D(v,k)=\prod_{i=v-k+1}^{v}(1-q^{-2i})$, with $D(v,0)=1$, and let
\[
 H(M)=J_b+M^{\spnT}J_a^{-1}M ,
\]
the residual form, alternating in every characteristic. Let $E$ be the event that its reduction
$\overline{H(M)}$ modulo $\varpi$ is singular, and $\beta_{q,a,b}=\nu_1(E^c)$. For $b=1$,
$H(M)=(1-\Omega(M))J_1$ and $E=\set{\Omega(M)\equiv1\pmod\varpi}$.

\begin{theorem}[First-digit optimality for several pairs; \stU]\label{spn:thm:padic-pairs}
Let $a,b\ge1$ and $n\ge a+b+\min(a,b)-1$, and put $\alpha=D(n-a,b)/D(n,b)$. For every
$e\in\set{1,2,\dots,\infty}$,
\[
 \spnTV(\mu_e,\nu_e)=(1-\alpha)\,\beta_{q,a,b}=\spnTV(\mu_1,\nu_1),
\]
and $E$ is an optimal distinguishing event at all these precisions. At each finite precision the
likelihood is exactly $\alpha<1$ on $E^c$ and strictly greater than one on $E$; over $\OO$ the two
measures have a continuous likelihood with the same strict signs.
\end{theorem}

So the rule that declares the symplectic model exactly when the residual form is singular modulo
$\varpi$ achieves the best equal-prior success probability $\frac12(1+\spnTV(\mu_e,\nu_e))$ even when
every digit is observed. This concerns equal-prior testing, not sufficiency: higher digits need not be
uniform and can matter for unequal priors or costs, and the range is sufficient, not claimed to be
sharp. At $e=1$ the statement is contained in \cref{spn:thm:offset-test}(1), with frame dimensions
$2a,2b$ and offset $n-a-b\ge\min(a,b)-1$; the content is the independence of the precision. With
$P_v=\prod_{i=1}^v(1-q^{-2i})$ one has $\alpha=P_{n-a}P_{n-b}/(P_{n-a-b}P_n)=D(n-b,a)/D(n,a)$, and
$\det(I+UV)=\det(I+VU)$ makes the residual events agree after transposition; since interchanging the
frames replaces $M$ by $-M^{\spnT}$ with $g\mapsto g^{-1}$, it suffices to prove the theorem when
$c=n-a\ge2b-1$.

For $s\ge2$ and $C$ uniform in $\spnMat_{2c,s}(R_e)$ put $m=\binom s2$ and
\[
 f_{c,s,e}(H)=q^{me}\,\spnP(C^{\spnT}J_cC=H)\qquad(H\in\spnAlt_s(R_e)),
\]
the density of the Gram matrix relative to the uniform alternating law; at $e=1$ it is the density
$F_c$ of \cref{spn:sec:nondegenerate}.

\begin{lemma}[The Gram differential; \stU]\label{spn:lem:gram-differential}
For $C_0\in\spnMat_{2c,s}(\kappa)$ with kernel $K$ of dimension $\ell$, the image of
$Z\mapsto C_0^{\spnT}J_cZ+Z^{\spnT}J_cC_0$ is exactly the set of alternating forms vanishing on
$K\times K$, of codimension $\binom\ell2$; the map is onto when $\ell\le1$. Consequently, if
$\dim\ker C_0\le1$ and $C_0^{\spnT}J_cC_0=H_0$, every lift $H$ of $H_0$ to $R_e$ is the Gram matrix of
exactly $q^{(2cs-m)(e-1)}$ lifts of $C_0$, so the residue class of $C_0$ contributes the constant
$q^{m-2cs}$ to $f_{c,s,e}(H)$, whatever the higher digits of $H$.
\end{lemma}

\begin{proof}
In a source basis with $C_0=[C'\;0]$ and $C'$ injective, $C'^{\spnT}J_c$ has full row rank, so the
upper left block of the image is an arbitrary matrix minus its transpose, the cross block is arbitrary
and the lower right $\ell\times\ell$ block vanishes; every alternating matrix is a matrix minus its
transpose, by choosing one triangle, with no division by two. Lifting from precision $j$ to $j+1$,
replacing $C$ by $C+\varpi^jZ$ changes the new Gram digit by the differential at $\overline Z$, the
quadratic term being divisible by $\varpi^{2j}$; so every required correction has exactly
$q^{2cs-m}$ solutions, and induction gives the count.
\end{proof}

\begin{lemma}[Likelihood reduction; \stU]\label{spn:lem:padic-reduction}
For finite $e$ and $c=n-a$, $\mu_e(M)/\nu_e(M)=f_{c,2b,e}(H(M))/D(n,b)$.
\end{lemma}

\begin{proof}
Symplectic Gram--Schmidt extends a frame with Gram matrix $J_b$ to a symplectic basis over $R_e$, so
the group acts transitively on such frames. In a symplectic basis beginning with $Y$ write
$gX=(B;C)$; then $M=J_aB$ and $C^{\spnT}J_cC=J_b-B^{\spnT}J_aB=H(M)$, and a matrix with Gram
matrix $J_b$ has a left inverse, so there is no further injectivity condition. By
\cref{spn:lem:gram-frames} over the residue field and \cref{spn:lem:gram-differential} there are
$q^{(4nb-m)e}D(n,b)$ source frames, and for fixed $B$ there are $q^{(4cb-m)e}f_{c,2b,e}(H(M))$
admissible $C$; divide by $\nu_e(M)=q^{-4abe}$.
\end{proof}

\begin{proposition}[A sign that cannot change with the precision; \stU]\label{spn:prop:padic-sign}
Let $s=2b\ge2$ and $c\ge s-1$. If $\overline H$ is singular, then $f_{c,s,e}(H)>1$ for every finite
$e$, with a lower bound independent of $e$; if $\overline H$ is nonsingular, $f_{c,s,e}(H)=D(c,b)$.
\end{proposition}

\begin{proof}
Let $\rank\overline H=2r$ and $t=s-2r\ge2$. The kernel of a matrix with Gram matrix $\overline H$
lies in its radical, which has $(q^t-1)/(q-1)$ lines, and by \cref{spn:lem:gram-differential} the
residue classes of kernel dimension zero and one contribute
\[
 f_{c,s,e}(H)\ge q^{m-2cs}\Bigl(\spnEmb_c(s,r)+\tfrac{q^t-1}{q-1}\spnEmb_c(s-1,r)\Bigr)
 =\prod_{j=u+1}^{c}(1-q^{-2j})\Bigl(1+q^{-2u}\sum_{j=1}^{t-1}q^{-j}\Bigr)
\]
with $u=c-s+r+1\ge0$; for $u=0$ the first count vanishes and the second alone gives the same
expression. With $A=q^{-2u}\le1$ the product exceeds $1-A/(q^2-1)$, and $t\ge2$ gives
$f>(1-A/(q^2-1))(1+A/q)=1+A(q^2-q-1-A)/(q(q^2-1))\ge1$. If $\overline H$ is nonsingular, every matrix
with that Gram matrix is injective, there are $q^{2cs-m}D(c,b)$ of them and all their lifts are flat.
\end{proof}

\begin{proof}[Proof of \cref{spn:thm:padic-pairs}]
Assume $c=n-a\ge2b-1$. Since $a>0$, $D(c,b)<D(n,b)<1$, so by \cref{spn:lem:padic-reduction} and
\cref{spn:prop:padic-sign} the likelihood is $\alpha<1$ on $E^c$ and exceeds one on $E$, uniformly in
the precision. The negative part of $\mu_e-\nu_e$ is $(1-\alpha)\nu_e|_{E^c}$, of mass
$(1-\alpha)\beta_{q,a,b}$, as $E$ depends only on the first digit. Over $\OO$ the cylinder sets
generate the Borel sets and the total variation is the supremum over the finite quotients, so the
same distance holds and the cylinder event $E$ attains it; the continuous likelihood and its strict
signs come from \cref{spn:thm:gram-regularity} below and the precision-independent lower bound.
\end{proof}

For general $s\ge2$, even or odd, the Gram densities have an exact Fourier expansion and a sharp
regularity threshold. Let $\psi_e$ be a primitive additive character of $R_e$ and
$\langle W,H\rangle=\sum_{i<j}W_{ij}H_{ij}$, put
$\mathcal Z_e(c,s)=\sum_{W\in\spnAlt_s(R_e)}q^{-c\,\mathrm{length}_\OO(\spnim W)}$, $\mathcal Z_0(c,s)=1$, and
$\mathcal Z_e(c,s)=1$ for $s\le1$, and let $A_s(r)$ be the number of alternating $s\times s$ matrices of rank
$2r$ over $\kappa$, as in \eqref{spn:eq:alt-rank-count}.

\begin{theorem}[Local Gram densities; \stU]\label{spn:thm:gram-regularity}
\begin{enumerate}
\item For every $H\in\spnAlt_s(R_e)$,
$f_{c,s,e}(H)=\sum_{W\in\spnAlt_s(R_e)}q^{-c\,\mathrm{length}_\OO(\spnim W)}\psi_e(-\langle W,H\rangle)$.
\item For $e\ge1$,
$\mathcal Z_e(c,s)-\mathcal Z_{e-1}(c,s)=\sum_{r=1}^{\floor{s/2}}A_s(r)q^{-2cr}q^{-r(2c-2s+2r+1)(e-1)}\mathcal Z_{e-1}(c,s-2r)$.
\item For $C$ additive-Haar in $\spnMat_{2c,s}(\OO)$, the law of $C^{\spnT}J_cC$ has a bounded Haar
density if and only if $c\ge s-1$. In that range the density $f_{c,s,\infty}$ is continuous, the
uniform limit of the finite densities, with
$\norm{f_{c,s,\infty}-f_{c,s,e}}_\infty=\mathcal Z_\infty(c,s)-\mathcal Z_e(c,s)$, and with $\lambda=2c-2s+3>0$,
\[
 \mathcal Z_\infty(c,s)-\mathcal Z_e(c,s)\sim\frac{A_s(1)q^{-2c}\mathcal Z_\infty(c,s-2)}{1-q^{-\lambda}}\,q^{-\lambda e}.
\]
\end{enumerate}
\end{theorem}

\begin{proof}
(1) Group the rows of $C$ into $c$ independent pairs $x,y$, whose Gram contribution
$xy^{\spnT}-yx^{\spnT}$ pairs with $W$ to $x^{\spnT}Wy$. For fixed $x$ the average of
$\psi_e(x^{\spnT}Wy)$ over $y$ is one if $Wx=0$ and zero otherwise, by primitivity of $\psi_e$, and
averaging over $x$ gives $\abs{\ker W}/\abs{R_e}^s=q^{-\mathrm{length}(\spnim W)}$; independence and
Fourier inversion finish. At $e=1$ this is \cref{spn:lem:alt-gram}(1).

(2) Matrices divisible by $\varpi$ contribute $\mathcal Z_{e-1}(c,s)$. For a residue matrix of rank $2r>0$, an
integral change of basis makes its upper left $2r$-block $A$ invertible, and congruence eliminates the
cross block, leaving the alternating Schur complement $D+B^{\spnT}A^{-1}B=\varpi T$ with
$T\in\spnAlt_{s-2r}(R_{e-1})$; the inverse of a nonsingular alternating matrix is alternating, also in
characteristic two. For each $T$ the choices of $A,B$ number
$q^{(e-1)(\binom s2-\binom{s-2r}2)}$, and the image length is $2re+\mathrm{length}(\spnim T)$.

(3) Choose compatible primitive characters at all levels. The Fourier characters of exact conductor
$e$ have total absolute coefficient mass $\mathcal Z_e-\mathcal Z_{e-1}$. If $c\ge s-1$, then
$r(2c-2s+2r+1)\ge r(2r-1)>0$, and induction on $s$ with (2) shows that $\mathcal Z_e(c,s)$ is bounded, so the
series converges absolutely and uniformly to a continuous function whose integrals over the finite
quotients agree with the Gram law. The omitted tail is at most $\mathcal Z_\infty-\mathcal Z_e$, with equality at
$H=0$. The term $r=1$ of (2) is $A_s(1)q^{-2c}q^{-\lambda(e-1)}\mathcal Z_{e-1}(c,s-2)$, and every term $r\ge2$
decays faster, which gives the asymptotic. If $c\le s-2$, the term $r=1$ alone gives
$\mathcal Z_e-\mathcal Z_{e-1}\ge A_s(1)q^{-2c}q^{-(2c-2s+3)(e-1)}$, which does not tend to zero, so
$f_{c,s,e}(0)=\mathcal Z_e(c,s)$, the average density on $\varpi^e\spnAlt_s(\OO)$, is unbounded.
\end{proof}

The density threshold $c=s-1$ is sharp; it is not claimed that the testing range of
\cref{spn:thm:padic-pairs} is, since a measure with an unbounded density can still have likelihood
above one on the relevant residue classes. The constant $\beta_{q,a,b}$ involves only the residue
field: for uniform $M$ the matrix $G=B^{\spnT}J_aB$ with $B=J_a^{-1}M$ satisfies $H(M)=J_b-G$, so,
with $Z_a(2b,r)$ the count of \cref{spn:sec:nondegenerate} of $2a\times2b$ matrices with a prescribed
Gram matrix of rank $2r$,
\[
 \beta_{q,a,b}=q^{-4ab}\sum_{\substack{G\in\spnAlt_{2b}(\kappa)\\ \det(J_b-G)\ne0}}Z_a\bigl(2b,\tfrac12\rank G\bigr),
\]
a sum over $q^{\binom{2b}2}$ matrices; the half-rank chain after \cref{spn:thm:offset-test}, started
at $b$ and run for $a$ steps, gives it in $O(ab)$ rational operations.

\begin{example}[Two pairs on both sides]\label{spn:ex:padic-pairs}
For $q=2$, $a=b=2$ and $n=5$, a $4\times4$ block in a symplectic $10$-dimensional space,
$\alpha=5376/5797$ and $\beta_{2,2,2}=469/1024$; exactly $30{,}016$ of the $65{,}536$ binary matrices
have a nonsingular residual. So
\[
 \spnTV(\mu_e,\nu_e)=\frac{197449}{5936128}=0.0332622544527\ldots
\]
for every $e$, including $e=\infty$, and the best equal-prior success probability is about
$0.5166311272$ however many digits are observed. Likewise $q=3$, $(a,b,n-a)=(2,2,3)$ gives
$\alpha=597051/605242$, $\beta=38089/59049$ and distance $7609439/871681338$, and $q=3$,
$(a,b,n-a)=(1,2,3)$ gives $81/82$, $163/243$ and $163/19926$.
\end{example}

Integer Fourier transforms of the four-column Gram likelihoods over $\Z/2$, $\Z/3$, $\Z/4$ and
$\F_2[t]/(t^2)$, the last two separating mixed and equal characteristic, an independent enumeration
of $137{,}889$ one-pair Gram inputs, the sign, normalization and distance at all $9{,}714$ residual
matrices of the displayed cases, the enumeration of $72{,}097$ observed matrices for $\beta$, $1{,}728$
rational cases of the strict lower bound and $60$ cases of the recursion against direct image-length
enumeration agree with these
statements.\footnote{\artifact{grh_extensions_2026-11-27_3/data/symplectic_precision_checks.json}}

\paragraph{Arbitrary primitive frames.}
Now let $X\in\spnMat_{2n,b}(\OO)$ and $Y\in\spnMat_{2n,a}(\OO)$ be any matrices whose reductions
have full column rank, with no condition on their Gram matrices, let $M=Y^{\spnT}JgX$, and keep
the notation $\mu_e,\nu_e,\widetilde\mu_{1,e}$ for the corresponding laws on
$\spnMat_{a,b}(R_e)$.

\begin{theorem}[Exact differential codimension; \stU]\label{spn:thm:padic-tangent}
Fix $g_0\in\Sp(J,\kappa)$, put $U=\spnim(g_0\overline X)$, $U'=\spnim\overline Y$ and
$s=\dim_\kappa(U\cap U')$. Conditional on $\overline g=g_0$ the first digit of $M$ is fixed and its
next digit is uniform on an affine subspace of $\spnMat_{a,b}(\kappa)$ of codimension $\binom s2$. The
following are equivalent: $s\le1$; the second digit is uniform on $\spnMat_{a,b}(\kappa)$; at every
finite precision $M$ is uniform on its residue coset; over $\OO$ its conditional law is Haar
measure on that coset. If $s\ge2$, the conditional law of the second digit is at distance exactly
$1-q^{-\binom s2}$ from uniform.
\end{theorem}

\begin{proof}
Reduction $\Sp(J,R_{j+1})\to\Sp(J,R_j)$ is surjective with fibres of equal size, in every residue
characteristic. Lift a symplectic matrix arbitrarily; its symplectic error, divided by $\varpi^j$
and reduced, is an alternating matrix $E$, and correcting the lift by $I+\varpi^jA$ requires
$A^{\spnT}J+JA=-E$, that is $K-K^{\spnT}=-E$ for $K=JA$, which is solved by choosing one triangle of
$K$; the solutions are the symmetric $K$, of which there are $q^{n(2n+1)}$. Choose a symplectic lift
of $g_0$. A first-order perturbation $g\mapsto(I+\varpi^jA)g$ changes the next digit of $M$ by
$\overline Y^{\spnT}JA\,g_0\overline X=\overline Y^{\spnT}K\,g_0\overline X$ with $K=K^{\spnT}$. So
the possible increments are the restrictions to $U'\times U$ of symmetric bilinear forms on
$\kappa^{2n}$. A cross-pairing $U'\times U\to\kappa$ extends to a symmetric form exactly when its
restriction to $(U\cap U')\times(U\cap U')$ is symmetric: necessity is clear, and for sufficiency
split $\kappa^{2n}=(U\cap U')\oplus U_1\oplus U'_1\oplus W$ with $U=(U\cap U')\oplus U_1$ and
$U'=(U\cap U')\oplus U'_1$, prescribe the form on pairs of summands by the cross-pairing and its
transpose, and choose the remaining entries freely. These are $\binom s2$ independent conditions,
diagonal entries imposing none, also in characteristic two. Uniform lifts therefore map to a
uniform affine subspace of codimension $\binom s2$, at distance $1-q^{-\binom s2}$ from uniform. If
$s\le1$ the differential is surjective at every later step, since its reduction does not change,
so every new digit has equally many preimages; induction and passage to the limit give uniformity
on the coset at every precision and Haar measure over $\OO$. Conversely uniformity at any level at
least two forces uniformity at level two, hence $\binom s2=0$.
\end{proof}

For an alternating space $W$ of dimension $d$ with radical of dimension $t$, the numbers of
nondegenerate and of totally isotropic two-dimensional subspaces are
\[
 D(W)=\frac{(q^d-q^t)q^{d-2}}{q^2-1},\qquad I(W)=\spnqbin d2q-D(W),
\]
both zero for $d<2$: choose a first vector outside the radical and a partner with pairing one,
and divide by $q(q^2-1)$.

\begin{theorem}[Precision-uniform lifting error; \stU]\label{spn:thm:padic-lifting}
Let $n\ge2$, $E=\spnim\overline X$, $F=\spnim\overline Y$, and let $p_{\mathrm{bad}}$ be the
probability that $\dim(F\cap\overline gE)\ge2$. For every finite $e$ and for $e=\infty$,
\[
 \spnTV(\mu_e,\widetilde\mu_{1,e})\le p_{\mathrm{bad}}\le\min\set{1,B_n(E,F)},
\]
\[
 B_n(E,F)=\frac{(q^2-1)(q-1)I(E)I(F)}{(q^{2n}-1)(q^{2n-2}-1)}+\frac{(q^2-1)D(E)D(F)}{(q^{2n}-1)q^{2n-2}},
\]
and
$\spnTV(\mu_1,\nu_1)\le\spnTV(\mu_e,\nu_e)\le\spnTV(\mu_1,\nu_1)+B_n(E,F)$. For a fixed residue
field and fixed frame dimensions the difference between the distances at precision $e$ and at the
residue field is $O(q^{-4n})$, uniformly in $e$, and by \cref{spn:thm:padic}(3) this order cannot
be improved in general.
\end{theorem}

\begin{proof}
Disintegrate over the uniform reduction $\overline g$. When the overlap has dimension at most one,
\cref{spn:thm:padic-tangent} says that the conditional law is the uniform lift of its residue; on
the remaining reductions the distance is at most one, and convexity of total variation gives the
first inequality. The number $Z$ of two-planes in $F\cap\overline gE$ dominates the indicator of
$\set{\dim(F\cap\overline gE)\ge2}$. The symplectic group is transitive on isotropic and on
nondegenerate two-planes separately, so $\spnE Z=I(E)I(F)/I(\kappa^{2n})+D(E)D(F)/D(\kappa^{2n})$, with
$I(\kappa^{2n})=(q^{2n}-1)(q^{2n-2}-1)/((q^2-1)(q-1))$ and $D(\kappa^{2n})=(q^{2n}-1)q^{2n-2}/(q^2-1)$.
Reduction modulo $\varpi$ cannot increase total variation, which gives the lower bound;
uniform lifting preserves the distance between the two residue laws, so the triangle inequality
gives the upper bound.
\end{proof}

Together with \cref{spn:thm:first-correction}, applied to the reduced Gram matrices, this shows
that for fixed $a,b$ the first correction $\frac{q^{-2n}}{2q^{ab}}\sum_M\abs{K(M)}$ to the distance
from independent entries holds with the same leading coefficient at every precision, with an error
$O(q^{-4n})$ independent of the precision. When the frame dimensions grow with $n$, the explicit
bound $B_n(E,F)$ is the applicable statement.

The exact laws were compared with exhaustive enumeration of $110{,}894{,}547$ ordered pairs, giving
$14{,}979{,}168$ symplectic frames and $20{,}067$ output probabilities, over $\Z/2$, $\Z/4$, $\Z/8$,
$\Z/3$, $\Z/9$, $\F_4$, $\F_2[t]/(t^2)$, $\F_2[t]/(t^3)$ and $\Z_2[\sqrt2]/((\sqrt2)^3)$ in
ambient dimensions four and six, which covers equal characteristic with nilpotents, a nonprime
residue field and ramification; the normalization, signs, both distance formulas, the precision
hierarchy and the envelope $8/5$ were checked for $704$ parameter triples in rational arithmetic,
the differential codimension for all $4{,}489$ pairs of subspaces of $\F_2^4$, the taller corners
in $7{,}073$ output probabilities over $\F_2$ and $\F_3$, and the likelihoods once more by an
independent implementation over $\Z/4$, $\F_2[t]/(t^2)$ and
$\F_4$.\footnote{\raggedright\artifact{padic_symplectic_corners/results/verification.json},
\artifact{padic_symplectic_corners/results/taller_corner_checks.json},
\artifact{padic_symplectic_corners/results/independent_ring_results.json}}
Lifting random-matrix statistics through congruence levels by the symplectic Lie algebra, with
uniform higher digits on favourable residue classes, is due to
\cite[Sections~8.2.1--8.2.3]{Pirani2024Traces}; the overlap formula $\binom s2$ is the point
specific to compressions. Related exact local laws concern corners of Haar $\GL_N(\OO)$ matrices
\cite{VanPeski2021Limits}, corners of additive-Haar alternating and Hermitian matrices, a
different ensemble \cite{ShenVanPeski2025NonArchimedean}, matrix Hua measures
\cite{Neretin2013Hua}, and cokernels of $I-S$ for full symplectic similitudes
\cite{Garton2015Random}.

%======================================================================

\subsection{Exact tests for a block of a split orthogonal matrix}\label{spn:sec:orthogonal}
\sources{grh_method_transfers_orthogonal_jets_2026-11-26}

The symplectic and orthogonal boundary laws are different, and neither
statement contains the other: in odd characteristic the determinant
component of the orthogonal group leaves a parity signature that has no
symplectic counterpart. Let $q$ be an odd prime power, $m\ge1$, and put
\[
 J=\begin{pmatrix}0&I_m\\I_m&0\end{pmatrix},\qquad
 G=\set{g\in\GL_{2m}(\F_q):g^{\spnT}Jg=J}.
\]
This is the split orthogonal group $O^+(2m,q)$. Its determinant-one subgroup
is denoted by $G^+=SO^+(2m,q)$; this notation does not mean the possibly
smaller commutator subgroup. Write
$g=\left(\begin{smallmatrix}A&B\\C&D\end{smallmatrix}\right)$ in this
hyperbolic basis, and let $U_m$ be a uniform $m\times m$ matrix over $\F_q$.
For laws $P,P'$ on the finite matrix space,
$\spnTV(P,P')=\frac12\sum_M\abs{P(M)-P'(M)}=\max_E\abs{P(E)-P'(E)}$.

\begin{theorem}[Optimal parity and full-rank tests; \stU]\label{spn:thm:orthogonal}
If $g$ is uniform in $G^+$, then $\rank C$ is even and
\begin{equation}\label{spn:eq:SO-TV}
 \spnTV(C,U_m)=\spnP\bigl(\rank U_m\text{ is odd}\bigr).
\end{equation}
Consequently the test ``the block has odd rank'' achieves the full total
variation distance; no test using all its entries improves this
distinguishing advantage. If instead $g$ is uniform in the full group $G$,
then
\begin{equation}\label{spn:eq:O-TV}
 \spnTV(C,U_m)=\prod_{j=1}^{m}(1-q^{-j})
 -\frac{1}{2\prod_{j=1}^{m-1}(1+q^{-j})},
\end{equation}
and the full-rank test achieves the full distance. Empty products equal one,
so the formula also includes $m=1$.
\end{theorem}

\begin{proof}
Let $W=\F_q^m\oplus0$ and $L=gW$. The space $L$ is maximal totally isotropic
for $J$ and $\rank C=m-\dim(L\cap W)$. Write
$T_m=2\prod_{i=1}^{m-1}(1+q^i)$, and let $e_1,\dots,e_m,f_1,\dots,f_m$ be the
hyperbolic basis.

\emph{The determinant and rank parity.} Every element of $G$ stabilizing $W$
has block form
$\left(\begin{smallmatrix}P&*\\0&P^{-\spnT}\end{smallmatrix}\right)$, hence
determinant one. Suppose $\dim(L\cap W)=m-r$. A change of basis in $W$,
extended by the inverse transpose on its complementary isotropic subspace,
puts the intersection in standard position. In its orthogonal quotient, $L$
is transverse to the remaining part of $W$, hence is the graph of a
skew-symmetric matrix. An upper unipotent isometry with a skew-symmetric
upper-right block removes that graph. Both changes stabilize $W$ and have
determinant one. The resulting space is
$\langle e_1,\dots,e_{m-r},f_{m-r+1},\dots,f_m\rangle$. It is the image of
$W$ under the isometry which swaps $e_i$ with $f_i$ in precisely $r$
hyperbolic pairs. That isometry has determinant $(-1)^r$. Multiplying by it
leaves an element stabilizing $W$, proving
\begin{equation}\label{spn:eq:orthogonal-parity}
 \det g=(-1)^{\rank C}.
\end{equation}
This construction also proves transitivity on maximal isotropic subspaces.
Since their stabilizer has determinant one, the two determinant cosets of $G$
give two equally sized families, selected by the parity of $r$.

\emph{The exact rank counts.} Choose $D=L\cap W$ of dimension $m-r$ in
$\spnqbin{m}{r}{q}$ ways. The quotient $D^\perp/D$ is a hyperbolic space of
dimension $2r$. The possible $L/D$ are precisely its maximal isotropic
complements to $W/D$. Writing each such complement as a graph gives a
skew-symmetric $r\times r$ matrix. There are $q^{r(r-1)/2}$ such matrices
because $q$ is odd. Thus the number of possible $L$ at block rank $r$ is
\begin{equation}\label{spn:eq:orthogonal-subspace-count}
 N_m(r)=\spnqbin{m}{r}{q}q^{r(r-1)/2}.
\end{equation}
The $q$-binomial identity
$\sum_{r=0}^{m}\spnqbin{m}{r}{q}q^{r(r-1)/2}=\prod_{i=0}^{m-1}(1+q^i)=T_m$
follows by induction from the Gaussian-binomial recursion. Uniform group
elements induce the uniform law on each orbit. Therefore the full orthogonal
rank law and the special orthogonal rank law are
\[
 o_m(r)=N_m(r)/T_m,\qquad
 s_m(r)=\begin{cases}2o_m(r),&r\text{ even},\\0,&r\text{ odd}.\end{cases}
\]
For an independent-entry matrix, choose its $r$-dimensional image and then a
surjection onto it. This gives
\[
 u_m(r)=\spnP(\rank U_m=r)=\spnqbin{m}{r}{q}F_q(m,r)q^{-m^2},\qquad
 F_q(m,r)=\prod_{i=0}^{r-1}(q^m-q^i).
\]

\emph{Why rank suffices for all matrix tests.} For each $P\in\GL_m(\F_q)$,
the matrix $\operatorname{diag}(P,P^{-\spnT})$ lies in $G^+$. Left and right
multiplication by these matrices transform $C$ into $P^{-\spnT}CP'$ for
independently arbitrary invertible $P,P'$. Uniform measure on either group is
invariant under these operations. Since left and right invertible
transformations act transitively on matrices of a given rank, both block
laws are uniform conditional on rank. The independent-entry law has the same
property. Thus the total variation distance of the full matrices equals that
of their rank laws.

\emph{Comparing the mass at each rank.} Set $R_m(r)=2o_m(r)/u_m(r)$, also for
odd $r$. Cancellation in the exact formulas gives
\[
 R_m(m)=\frac{1}{(1-q^{-m})\prod_{j=1}^{m-1}(1-q^{-2j})}>1,\qquad
 \frac{R_m(r-1)}{R_m(r)}=q^{m-r+1}-1 .
\]
Because $q\ge3$, each displayed ratio is at least two. Hence
$s_m(r)>u_m(r)$ at every even rank, while $s_m(r)=0$ at every odd rank. All
the deficit lies at the odd ranks, proving \eqref{spn:eq:SO-TV}.

For the full orthogonal group, the likelihood ratio is $R_m(r)/2$. At rank
$m$ it is less than one: writing $Q=q^{-1}\le1/3$, its denominator before the
factor two satisfies
\[
 (1-Q^m)\prod_{j=1}^{m-1}(1-Q^{2j})
 \ge(1-Q)\Bigl(1-\sum_{j=1}^{\infty}Q^{2j}\Bigr)\ge\frac{7}{12}>\frac12 .
\]
The ratio at rank $m$ is also greater than $1/2$, because that same product
is less than one. At rank $m-1$ the ratio is therefore greater than
$(q-1)/2\ge1$, and it increases as rank decreases. Thus only full rank has
deficient mass in $o_m$ compared with $u_m$. Their full-rank probabilities
are respectively
\[
 o_m(m)=\frac{1}{2\prod_{j=1}^{m-1}(1+q^{-j})},\qquad
 u_m(m)=\prod_{j=1}^{m}(1-q^{-j}).
\]
Their difference proves \eqref{spn:eq:O-TV}.
\end{proof}

With equally likely hypotheses, the best success probability of a test is
$(1+\spnTV)/2$. The group counts used are classical; for related counts of
isotropic subspaces see \cite[Lemma~3.2.2]{Li2022Arithmetic}, and for
finite-field rank laws \cite{Fulman2015Stein}.

\begin{example}\label{spn:ex:orthogonal-q3}
For $q=3$ the theorem gives the following values, which are confirmed by
exhaustive enumeration of all maximal isotropic subspaces and all square
matrices for $m\le3$, and of all $2$ and $576$ elements of $SO^+(2m,3)$ for
$m=1,2$.\footnote{\artifact{grh_method_transfers_orthogonal_jets_2026-11-26/checks/orthogonal_blocks_checks.json}}
\begin{center}
\begin{tabular}{c|c|c|l}
 $m$ & $\spnTV(C_{SO},U_m)$ & $\spnTV(C_{O},U_m)$ & nonzero $SO$ rank masses\\\hline
 $1$ & $2/3$ & $1/6$ & $r=0\colon 1$\\
 $2$ & $32/81$ & $47/216$ & $r=0\colon 1/4,\quad r=2\colon 3/4$\\
 $3$ & $11570/19683$ & $13597/58320$ & $r=0\colon 1/40,\quad r=2\colon 39/40$
\end{tabular}
\end{center}
\end{example}

For the full orthogonal group, fixed $q$ and $m\to\infty$ in
\eqref{spn:eq:O-TV} give the positive limiting distance
\[
 \prod_{j=1}^{\infty}(1-q^{-j})-\frac{1}{2\prod_{j=1}^{\infty}(1+q^{-j})} .
\]
The products converge to positive limits, and the proof above shows that the
difference is positive for every odd $q$. Even when $q\to\infty$, this
expression tends to $1/2$. The matrix entries therefore do not become
independent at this maximal isotropic boundary: the determinant component
leaves a precise parity signature, and averaging the two components leaves a
precise full-rank signature.

%======================================================================

\subsection{The largest principal kernel of a random binary orthogonal matrix}\label{spn:sec:binary}
\sources{orthogonal_nullity_and_prime_products; large_kernels_and_dirichlet_zeros; grh_method_transfers_orthogonal_jets_2026-11-26}

Let $\spnOrth_n=\set{Q\in\GL_n(\F_2):Q^{\spnT}Q=I}$, the isometries of the dot product, and
for $Q\in\spnOrth_n$ and $S\subseteq[n]$ let $Q[S]$ be the principal submatrix on $S$, the empty
matrix having nullity zero. Vaas \cite{Vaas2026Principal} asks for the typical size of
\[
 \nu(Q)=\max_{S\subseteq[n]}\dim\ker Q[S],\qquad
 Q\text{ uniform in }\spnOrth_n ,
\]
and whether constants $c,\alpha>0$ can make
$\spnP(\nu(Q)\ge cn)\ge1-2^{-\alpha n^2}$. The answer is that $\nu(Q)$ is concentrated on the
two integers next to $\sqrt n$; in particular the proposed linear behaviour fails. Far above
$\sqrt n$, and up to saturation at $n/2$, the probability of a large kernel is determined up to a
factor $1+o(1)$. Everything
below uses the symmetric bilinear dot product over $\F_2$; an orthogonal group attached to a
quadratic form would be a different problem.

Write $\mathbf j_m=(1,\dots,1)^{\spnT}\in\F_2^m$. Every $Q\in\spnOrth_n$ fixes $\mathbf j_n$:
$x\cdot x=x\cdot\mathbf j_n$ for all $x$, so $Qx\cdot\mathbf j_n=Qx\cdot Qx=x\cdot\mathbf j_n$,
and nondegeneracy gives $Q^{-1}\mathbf j_n=\mathbf j_n$. Call an $m\times d$ binary matrix
\emph{regular} if its columns are independent and $\mathbf j_m$ is not in their span.

\begin{lemma}[Comparison on free rows; \stU]\label{spn:lem:fibre}
Let $X$ be a regular $n\times d$ binary matrix, let $Q$ be uniform in $\spnOrth_n$, and let
$Y$ be an $n\times d$ matrix of independent fair bits. Then $QX$ is uniform on the regular
frames with Gram matrix $X^{\spnT}X$. If an event $E$ of positive probability for $n\times d$
matrices leaves $m\ge2d+5$ specified rows completely unrestricted, then
\[
 \Bigl|\frac{\spnP(QX\in E)}{\spnP(Y\in E)}-1\Bigr|\le8\cdot2^{2d-m},
\]
uniformly in $X$, in its Gram matrix and in the conditions $E$ imposes on the other rows.
\end{lemma}

\begin{proof}
\emph{The orbit.} Let $n$ be odd. The projection $v\mapsto v+(v\cdot v)\mathbf j_n$ along
$\mathbf j_n$ maps onto $\mathbf j_n^\perp$, on which the dot product is alternating and
nondegenerate, and is injective on the span of each regular frame; the projected pairings are
$v\cdot w+(v\cdot v)(w\cdot w)$, so two regular frames with the same Gram matrix induce an
isometry of the projected spans. Witt's extension theorem for alternating forms, valid also for
degenerate subspaces (split off a nondegenerate summand, pair a basis of its radical with
suitable vectors in the orthogonal complement of the summand, and complete by symplectic
pairs), extends it to $\mathbf j_n^\perp$, and fixing $\mathbf j_n$ gives an orthogonal map
carrying one frame to the other. For even $n$ adjoin a unit vector $e$, embed the frames with
last coordinate zero and append $e$; the augmented frames are regular because
$\mathbf j_{n+1}=\mathbf j_n+e$, and the odd-dimensional extension fixes $e$, hence restricts to
$e^\perp$. So $QX$ has the law of $Y$ conditioned on regularity and $Y^{\spnT}Y=X^{\spnT}X$.

\emph{Local counts.} Let $Z$ be an $m\times d$ matrix of fair bits, $D=d(d+1)/2$, $H$ any
symmetric $d\times d$ binary matrix and $a=2^{2d-m}$. Then
\begin{equation}\label{spn:eq:fibre-bounds}
 2^D\spnP(Z^{\spnT}Z=H,\ Z\text{ regular})\ge1-\tfrac43a,\qquad
 2^D\spnP(Z^{\spnT}Z=H)\le(1-4a)^{-1}.
\end{equation}
For the first, expose the columns: given $i$ regular columns, the $i$ products and the norm of
the next column, which is its product with $\mathbf j_m$, are $i+1$ independent linear
conditions; of the $2^{m-i-1}$ solutions at most $2^{i+1}$ lie in the span of $\mathbf j_m$ and
the earlier columns. Multiply and use $\prod(1-t_i)\ge1-\sum t_i$ and
$\sum_{i<d}2^{2i+2-m}<4a/3$. For the second, a Fourier character on symmetric matrices,
evaluated at a row contribution $zz^{\spnT}$, has the form $(-1)^{Q(z)}$ with $Q$ a
squarefree polynomial of degree at most two without constant term. If its alternating polar
form $B$ has rank $2s$, then $(\spnE_z(-1)^{Q(z)})^2=2^{-d}\sum_{h\in\operatorname{rad}B}(-1)^{Q(h)}$,
by writing two independent inputs as $z,z+h$. So the bias vanishes unless $Q$ vanishes on the
radical, and otherwise has absolute value $2^{-s}$. For fixed $B$ at most $2^{2s}$ linear
adjustments of $Q$ vanish on the radical, there are at most $2^{2ds}$ forms $B$ of rank $2s$,
and the only character of rank zero with nonzero bias is trivial. Fourier inversion and
independence of the $m$ rows bound the deviation of $2^D\spnP(Z^{\spnT}Z=H)$ from $1$ by
$\sum_{s\ge1}2^{(2d+2-m)s}=4a/(1-4a)$.

\emph{Conclusion.} Condition on the restricted rows of $Y$; they contribute a fixed symmetric
matrix to the Gram matrix, and a regular free-row block makes the whole matrix regular. By
\eqref{spn:eq:fibre-bounds} the conditional probability of a regular frame with a prescribed
Gram matrix lies between $2^{-D}(1-4a/3)$ and $2^{-D}(1-4a)^{-1}$, with or without $E$. Bayes'
rule and the orbit description give the ratio bound, since $a\le1/32$ and
$[(1-4a)(1-4a/3)]^{-1}-1\le8a$.
\end{proof}

For $S\subseteq[n]$ let $V_S$ be the coordinate subspace of vectors supported on $S$. A kernel
vector of $Q[S]$ is a vector of $V_S\cap Q^{-1}V_{S^c}$; applying $Q$ identifies this space
with the kernel of $Q^{-1}[S^c]=Q^{\spnT}[S^c]$, so
\begin{equation}\label{spn:eq:comp}
 \dim\ker Q[S]=\dim\ker Q[S^c],
\end{equation}
and $\nu(Q)\le\floor{n/2}$. Let $G(b,r)$ be the number of $r$-dimensional subspaces of
$\F_2^b$, zero for $r>b$, and put $c_\infty=\prod_{i\ge1}(1-2^{-i})^{-1}$. Counting ordered
bases, $G(b,r)\le c_\infty2^{r(b-r)}$, and $c_\infty<4$, since the first two factors of the
product times $1-\sum_{i\ge3}2^{-i}$ give $9/32>1/4$.

\begin{theorem}[An explicit tail for the maximum principal nullity; \stU]\label{spn:thm:binary}
For every integer $r$ with $1\le r\le(n-2)/2$,
\[
 \spnP\set{\nu(Q)\ge r}\le\tfrac32\,c_\infty\,2^{\,n-r^2+r}.
\]
\end{theorem}

\begin{proof}
By \eqref{spn:eq:comp} it suffices to consider supports $S$ with $b=\abs S\le n/2$. A basis
$X$ of an $r$-plane in $V_S$ is regular, and the plane lies in $\ker Q[S]$ exactly when $QX$ is
supported on $S^c$. Write $g(r)=r(r+1)/2$. The column exposure in the proof of
\eqref{spn:eq:fibre-bounds}, with $m=n$, shows that the orbit of $X$ has at least
$(1-\tfrac43 2^{2r-n})2^{nr-g(r)}\ge\tfrac23 2^{nr-g(r)}$ elements when $r\le(n-2)/2$. Within
the $n-b$ coordinates outside $S$ there are at most $2^{(n-b)r-r(r-1)/2}$ independent frames
with the prescribed Gram matrix, since after $i$ columns the pairings with them are $i$
independent linear conditions on the next one in this nondegenerate space. By the orbit
description of \cref{spn:lem:fibre}, the plane is a kernel plane with probability at most
$\tfrac32 2^{-br+r}$. Summing over the at most $c_\infty2^{r(b-r)}$ planes of $V_S$ and the at
most $2^n$ supports gives the bound.
\end{proof}

At $r=\floor{n/2}$, apply the theorem at $r-1$. Comparing exponents, the resulting bounds imply
$\spnP\set{\nu(Q)\ge r}\le6\cdot2^{\,n-r(r-1)/2}$ for every $1\le r\le\floor{n/2}$; for
$r\le(n-2)/2$ this is immediate from $r^2-r\ge r(r-1)/2$ and $\tfrac32c_\infty<6$, and at
$r=\floor{n/2}$ the bound at $r-1$ is the smaller one when $r\ge4$, while for $r\le5$ the right
side $6\cdot2^{n-r(r-1)/2}$ is at least one.

\begin{theorem}[Threshold and two-point concentration; \stU]\label{spn:thm:threshold}
There are absolute constants $c,C>0$ such that for all sufficiently large $n$ and all integers
$\sqrt n/2\le r\le2\sqrt n$,
\begin{align}
 \spnP(\nu(Q)<r)&\le C\bigl(2^{r^2-n}+2^{-r}\bigr),\label{spn:eq:fail}\\
 c\min\set{1,2^{n-r^2}}\le\spnP(\nu(Q)\ge r)&\le\min\set{1,C2^{n-r^2}}.\label{spn:eq:hit}
\end{align}
If $h$ is the nearest integer to $\sqrt n$, then
$\spnP\set{\nu(Q)\in\set{h-1,h}}=1-O(2^{-c\sqrt n})$. Along every sequence with
$\min_{j\ge0}\abs{n-j^2}\to\infty$ one has $\spnP\set{\nu(Q)=\floor{\sqrt n}}\to1$.
\end{theorem}

The quantity $2^{n-r^2}$ is a first moment: a fixed central principal block contains about a
constant times $2^{-r^2}$ kernel subspaces of dimension $r$ in expectation, and there are about
$2^n$ blocks. A second moment shows that these witnesses do not concentrate on a small
exceptional set of matrices.

\begin{proof}
\emph{First moment.} For a support of size $b\le n/2$ and an $r$-plane $U\le V_S$ with basis
$X$, the event $U\le\ker Q[S]$ asks the rows of $QX$ in $S$ to vanish; for independent bits it
has probability $2^{-br}$, and \cref{spn:lem:fibre} applies with $n-b$ free rows. Summing over
at most $c_\infty2^{r(b-r)}$ planes and over the $2^{n-1}$ complementary pairs of supports gives,
whenever $n/2\ge2r+5$,
\begin{equation}\label{spn:eq:upper}
 \spnP(\nu(Q)\ge r)\le2^{n-1}c_\infty2^{-r^2}\bigl(1+8\cdot2^{2r-n/2}\bigr),
\end{equation}
which is the upper bound in \eqref{spn:eq:hit}.

\emph{Second moment.} Let $W_S(Q)$ be the number of $r$-planes $U\le V_S$ with $(Qx)_i=0$ for
all $x\in U$ and $i\in S$; it equals $G(\dim\ker Q[S],r)$. Put
\[
 Z=\sum_{S\in\mathcal C}W_S(Q),\qquad\mathcal C=\set{S\subseteq[n]:n/3\le\abs S\le2n/3},
\]
so that $\abs{\mathcal C}\ge(1-9/n)2^n$ by Chebyshev's inequality. Constants below are
absolute, $\tfrac12\sqrt n\le r\le2\sqrt n$, and $n$ is large. Let $Q_0$ have independent fair
entries, and $c_r=\prod_{i=1}^r(1-2^{-i})^{-1}$. For $\abs S=b$,
$\mu_b=\spnE W_S(Q_0)=G(b,r)2^{-br}=c_r2^{-r^2}\prod_{i<r}(1-2^{i-b})$, and the last product
differs from $1$ by at most $2^{r-b}$. Since the event ignores $n-b$ rows,
\cref{spn:lem:fibre} gives, uniformly on $\mathcal C$,
\begin{equation}\label{spn:eq:onecentral}
 \spnE W_S(Q)=\bigl(1+O(2^{-cn})\bigr)c_r2^{-r^2}.
\end{equation}
For $S,R\in\mathcal C$ write $b=\abs S$, $b'=\abs R$, $a=\abs{S\cap R}$ and $D=b+b'-2a$, and
draw $U\le V_S$, $V\le V_R$ independently and uniformly among $r$-planes, with
$i=\dim(U\cap V)$. In the independent model each row in $S\cap R$ carries $2r-i$ independent
constraints and each row in $S\triangle R$ carries $r$, so
$\spnE[W_S(Q_0)W_R(Q_0)]/(\mu_b\mu_{b'})=\spnE_{U,V}2^{ai}$. Counting linear maps
$\F_2^a\to U\cap V$ by the dimension of their image,
$2^{ai}=\sum_{j\le\min(a,i)}G(i,j)\prod_{h<j}(2^a-2^h)$. A common $j$-plane lies in
$V_{S\cap R}$, and it lies in $U$ with probability
$\prod_{h<j}(2^r-2^h)/(2^b-2^h)\le2^{-j(b-r)}$, similarly for $V$. With
$G(a,j)\le c_\infty2^{j(a-j)}$ this gives
\begin{equation}\label{spn:eq:overlap}
 0\le\spnE_{U,V}2^{ai}-1\le B_D,\qquad B_D=c_\infty\sum_{j=1}^r2^{j(2r-D-j)}.
\end{equation}
For $S,T\in\mathcal C$, choose $R\in\set{T,T^c}$ minimizing $D=\abs{S\triangle R}$; then
$D\le n/2$, $R\in\mathcal C$ and $n-\abs{S\cup R}\ge n/12$, and $W_T(Q)=W_R(Q)$ by
\eqref{spn:eq:comp}. For any pair of planes on $S,R$, a basis of their sum spans no
$\mathbf j_n$, because $S\cup R$ omits a coordinate, and the joint witness event ignores at
least $n/12$ rows. \Cref{spn:lem:fibre} with at most $2r$ columns therefore compares its
probability with the independent model up to a factor $1+O(2^{-cn})$; a row of fair bits
applied to a matrix of full column rank is uniform on the output rows. Summing over planes,
\[
 \frac{\spnE Z^2}{(\spnE Z)^2}\le\bigl(1+O(2^{-cn})\bigr)\bigl(1+\spnE_{S,T}B_D\bigr),
\]
with $S,T$ independent and uniform in $\mathcal C$.

\emph{Averaging the overlap.} For fixed $S$, the proportion of $T$ with $D=d$ is at most
$2\binom nd/\abs{\mathcal C}\ll2^{-n}\binom nd$. Completing the square gives
$B_d\ll2^{(r-d/2)^2}$, since $\sum_{j\in\Z}2^{-(j-t)^2}$ is bounded uniformly in real $t$.
For $d<r$, $\binom ndB_d\ll2^{r^2+d\log_2n-rd+d^2/4}\le2^{r^2-rd/2}$, and these $d$ contribute
$O(2^{r^2-n})$ to $\spnE B_D$. For $r\le d\le3r$ each exponent in \eqref{spn:eq:overlap} is at
most $r^2/4$, and these distances contribute $O(2^{-n+r^2/4+O(r\log n)})=O(2^{r^2-n})$, since
$r\asymp\sqrt n$. For $d>3r$ the geometric majorant gives $B_d\ll2^{-r}$. Hence
\begin{equation}\label{spn:eq:variance}
 \frac{\operatorname{Var}Z}{(\spnE Z)^2}\ll2^{r^2-n}+2^{-r},
\end{equation}
the factor $1+O(2^{-cn})$ being absorbed. Chebyshev's inequality gives
$\spnP(\nu(Q)<r)\le\spnP(Z=0)\le\operatorname{Var}Z/(\spnE Z)^2$, which is \eqref{spn:eq:fail},
and Cauchy--Schwarz gives $\spnP(\nu(Q)\ge r)\ge\spnP(Z>0)\ge(\spnE Z)^2/\spnE Z^2$, the lower
bound in \eqref{spn:eq:hit}.

\emph{Concentration.} Since $(h-\tfrac12)^2\le n\le(h+\tfrac12)^2$, both $n-(h-1)^2$ and
$(h+1)^2-n$ are at least $h-1$; apply \eqref{spn:eq:fail} at $r=h-1$ and the upper bound of
\eqref{spn:eq:hit} at $r=h+1$. If $k=\floor{\sqrt n}$ and $n$ recedes from the squares, both
$n-k^2$ and $(k+1)^2-n$ tend to infinity, and the bounds at $k$ and $k+1$ give the last
assertion.
\end{proof}

\begin{corollary}[Two values are sometimes necessary; \stU]\label{spn:cor:twovalues}
Along $n=k^2-1$ both $\spnP(\nu(Q)=k-1)$ and $\spnP(\nu(Q)=k)$ have positive lower limits. At
$n=k^2+k$, by contrast, $\nu(Q)=k$ with probability tending to one.
\end{corollary}

\begin{proof}
By \cref{spn:thm:threshold} the maximum lies in $\set{k-1,k}$ with probability tending to one.
The lower bound of \eqref{spn:eq:hit} at $r=k$ is at least $c/2$, and by \eqref{spn:eq:upper}
the probability of $\nu(Q)\ge k$ is at most $c_\infty/4+o(1)<1$. The second assertion is the
last part of \cref{spn:thm:threshold}.
\end{proof}

\begin{proposition}[The leading constant above the threshold; \stU]\label{spn:prop:raretail}
If $r/\sqrt n\to1$ and $r^2-n\to+\infty$, then
\[
 \spnP\bigl(\nu(Q)\ge r\bigr)\sim\tfrac12\,c_\infty\,2^{\,n-r^2}.
\]
The factor $\tfrac12$ reflects the coincidence \eqref{spn:eq:comp} of complementary supports.
\end{proposition}

\begin{proof}
Put $\lambda=2^{n-r^2}$, choose one member of each complementary pair of central supports,
and let $N$ count the chosen supports with nullity at least $r$. For a block of nullity $h$,
counting flags of dimensions $r,r+1$ gives $G(h,r)(2^{h-r}-1)=G(h,r+1)(2^{r+1}-1)$, so on
$h\ge r+1$ the witness count $G(h,r)$ is at most $(2^{r+1}-1)G(h,r+1)$. By
\eqref{spn:eq:onecentral} at $r$ and $r+1$, replacing the witness count by the indicator of
$h\ge r$ costs $O(2^{-r})$ of its mean, so $\spnE N\sim c_\infty\lambda/2$. For the second
factorial moment use $\mathbf 1_{h\ge r}\le G(h,r)$ at two distinct representatives and repeat
the pair comparison; now $D=0$ is impossible, each representative has at most $2\binom nd$
partners at distance $d$, the distances $1\le d<r$ contribute $O(\lambda^{-1}2^{-r/2})$ to the
normalized pair bound, the distances $r\le d\le3r$ contribute $o(1)$ because $r^2/n\to1$, and
larger distances $O(2^{-r})$. Hence $\spnE[N(N-1)]\ll\lambda^2+\lambda2^{-r/2}=o(\lambda)$, and
the first two Bonferroni inequalities give $\spnP(N>0)=\spnE N+o(\lambda)$. Among the remaining
supports take the smaller member of each complementary pair, of size $b<n/3$; those with
$b<r$ cannot witness the event, and for the others the single-plane comparison on $n-b$ free
rows gives $\spnE W_S(Q)\ll2^{-r^2}$. There are $o(2^n)$ noncentral supports, so they
contribute $o(\lambda)$.
\end{proof}

\paragraph{Large kernels far above the threshold.}
For $r/\sqrt n\to\infty$ the rare event is carried by a single complementary pair, and its
probability is known up to a factor $1+o(1)$. Write $\gamma=c_\infty^{-1}=\prod_{j\ge1}(1-2^{-j})$,
$a=\prod_{j\ge1}(1-2^{-2j})$, $b_S=\dim\ker Q[S]$, and
\[
 B_{n,r}=\sum_{s=r}^{n-r}\binom ns ;
\]
by \eqref{spn:eq:comp}, $b_S\ge r$ forces $r\le\abs S\le n-r$, so $B_{n,r}/2$ counts the feasible
complementary pairs. By \cref{spn:lem:fibre}, all vectors of a given parity other than $0$ and
$\mathbf j_d$ form one orbit when present; counting the orbit of a unit vector, which has
$2^{d-1}-\mathbf 1_{d\text{ odd}}$ elements for $d\ge2$, and then its stabilizer gives
\begin{equation}\label{spn:eq:orth-order}
 \abs{\spnOrth_d}=2^{d(d-1)/2}a_d,\qquad a_d=\prod_{1\le i\le\floor{(d-1)/2}}(1-2^{-2i}),
\end{equation}
with $\abs{\spnOrth_0}=\abs{\spnOrth_1}=1$. Call a $j$-plane of $\F_2^v$ \emph{standard} if it is
in the isometry orbit of a coordinate plane. Its stabilizer is $\spnOrth_j\times\spnOrth_{v-j}$,
so by \eqref{spn:eq:orth-order} and $2^{j(v-j)}\le G(v,j)\le\gamma^{-1}2^{j(v-j)}$ the standard
planes form at least an absolute positive fraction of all $j$-planes, uniformly in $v,j$.

\begin{lemma}[One principal submatrix; \stU]\label{spn:lem:onesupport}
Let $\abs S=s$ and $W_j(S)=G(b_S,j)$. Uniformly for $j\le\min(s,n-s)$,
$\spnE W_j(S)\asymp2^{-j^2}$. If $j\to\infty$ and $\min(s-j,n-s-j)\to\infty$, then
$\spnE W_j(S)\sim\abs{\GL_j(\F_2)}^{-1}$, and the same equivalent and comparison hold, as
$j\to\infty$, for $\spnP(b_S\ge j)$ and $\spnP(b_S=j)$.
\end{lemma}

\begin{proof}
The kernel of $Q[S]$ is identified with $QV_S\cap V_{S^c}$. For an unrestricted uniform
$s$-plane the expected number of $j$-planes in its intersection with $V_{S^c}$ is
$G(n-s,j)G(s,j)/G(n,j)$, and domination by standard planes gives the upper estimate. For the
lower estimate restrict to the $\abs{\spnOrth_s}/(\abs{\spnOrth_j}\abs{\spnOrth_{s-j}})$ standard
$j$-planes in $V_S$; fix an orthonormal basis of each and require its image to be an orthonormal
frame in $V_{S^c}$ extendible to an orthonormal basis there. This has probability at least
$(\abs{\spnOrth_{n-s}}/\abs{\spnOrth_{n-s-j}})/(\abs{\spnOrth_n}/\abs{\spnOrth_{n-j}})$, and
\eqref{spn:eq:orth-order} gives the lower bound, also when a complement is zero-dimensional.

For the constant put $D=j(j+1)/2$. By the column exposure of \eqref{spn:eq:fibre-bounds}, for
every symmetric Gram matrix the number of regular $n\times j$ frames with that Gram matrix is
$2^{nj-D}(1+O(2^{2j-n}))$ as $n-2j\to\infty$. Independent frames whose span contains
$\mathbf j_n$ add at most $(2^j-1)2^{n(j-1)-j(j-1)/2}$, by moving the coefficient vector of
$\mathbf j_n$ to the last column; in particular for the zero Gram matrix there are at most
$2\cdot2^{nj-D}$ independent frames whenever $2j\le n$. An ordered kernel witness with its image
is a pair of rank-$j$ frames $X,Y$ supported on $S$ and $S^c$ with $X^{\spnT}X=Y^{\spnT}Y$ and
$QX=Y$; stacking their nonzero rows gives an $n\times j$ matrix $Z$ with $Z^{\spnT}Z=0$. In a
uniform regular such $Z$ every nonzero combination of columns is uniform on the even vectors
other than $0,\mathbf j_n$, so a union bound shows that either restricted rank fails with
probability $O(2^{j-s}+2^{j-(n-s)})$, and the number of eligible pairs is $(1+o(1))2^{nj-D}$. Each
of $X,Y$ is regular, and $\spnP(QX=Y)=(1+o(1))2^{-nj+D}$ by the Gram count. Dividing by the
$\abs{\GL_j(\F_2)}$ ordered bases gives the equivalent. Finally
$0\le W_j-\mathbf 1_{b_S\ge j}\le(2^{j+1}-1)W_{j+1}$ by counting flags, and the upper estimate
bounds the excess by $O(2^{-j^2-j})$.
\end{proof}

A second exceptional pair is controlled through self-dual codes. A self-dual binary code is an
$n$-plane $L=L^\perp$ of $\F_2^{2n}$ for the dot product; it contains $\mathbf j_{2n}$ and is even,
and self-dual codes correspond to the Lagrangians of the symplectic space
$\mathbf j_{2n}^\perp/\langle\mathbf j_{2n}\rangle$ of dimension $2(n-1)$, of which there are
$\prod_{i=1}^{n-1}(2^i+1)$. A uniform self-dual code is the graph of a linear map on the first
$n$ coordinates with probability $p_n=\abs{\spnOrth_n}/\prod_{i=1}^{n-1}(2^i+1)\gg1$, since those
maps are exactly the isometries; conditioning on the graph event costs only a constant in an
upper bound. For $S\subseteq[n]$ put $X_S=V_S\oplus V_{S^c}$, the first summand in the input and
the second in the output, so that $X_S^\perp=X_{S^c}$, and $d_S=\dim(L\cap X_S)$. For a graph
$d_S=b_S$, and for every self-dual code $(\operatorname{proj}_{X_S}L)^{\perp}=L\cap X_S$ inside
$X_S$, so $d_S=d_{S^c}$. For fixed $S$,
\begin{equation}\label{spn:eq:sdone}
 \spnP_{\mathrm{sd}}(d_S\ge r)\ll2^{-r^2}:
\end{equation}
an isotropic $r$-plane in $X_S$ omits $\mathbf j_{2n}$ and lies in a random self-dual code with
probability $\prod_{i=n-r}^{n-1}(2^i+1)^{-1}\le2^{-nr+r(r+1)/2}$, by passing to the quotient by
its image, and the zero-Gram count above bounds the number of such planes by
$2\cdot2^{nr-r(r+1)/2}/\abs{\GL_r(\F_2)}$.

\begin{lemma}[Two complementary pairs; \stU]\label{spn:lem:twopairs}
Let $S,T$ represent distinct complementary pairs, and let
$e=\min(\abs{S\triangle T},n-\abs{S\triangle T})$. There are absolute constants $c,C>0$ such
that, for all sufficiently large $r$, conditional on $d_S=r$ in the self-dual ensemble,
\[
 \spnP_{\mathrm{sd}}(d_T\ge r\mid d_S=r)\le
 \begin{cases}C2^{-e(r-e)/2},&1\le e\le r/2,\\ C2^{-cr^2},&r/2\le e\le n/2.\end{cases}
\]
\end{lemma}

\begin{proof}
Replacing $T$ by $T^c$ if necessary, $X_T$ is the sum of an $(n-e)$-plane of $X=X_S$ and an
$e$-plane of $Y=X_S^\perp$. Conditional on $d_S=r$ the law of $L$ is invariant under
$\spnOrth(X)\times\spnOrth(Y)$, so one may fix $L$ and randomize those two planes over their
standard orbits; up to an absolute constant this reduces to independent uniform planes
$P\subset X$ and $P'\subset Y$ of dimensions $n-e$ and $e$. Put $W=L\cap X$, $U=\operatorname{proj}_XL$,
and $W',U'$ likewise in $Y$; their dimensions are $r,n-r,r,n-r$, and $(x,y)\in L$ induces an
isomorphism $f\colon U/W\to U'/W'$ of spaces of dimension $n-2r$. With
$a=\dim(W\cap P)$, $g=\dim(U\cap P)$, $a'=\dim(W'\cap P')$, $g'=\dim(U'\cap P')$, and $J,J'$ the
images of $U\cap P$, $U'\cap P'$ in the two quotients,
\begin{equation}\label{spn:eq:residual}
 d:=\dim\bigl(L\cap(P\oplus P')\bigr)=a+a'+\dim(fJ\cap J'),
\end{equation}
since sending a pair in the intersection to its common quotient coordinate is onto the last
intersection with kernel $(W\cap P)\oplus(W'\cap P')$. Conditional on the four dimensions, $J$ and
$J'$ are independent uniform planes, because automorphisms of either quotient lift to
automorphisms preserving the flag $W\subset U$. For a uniform $p$-plane and a fixed $v$-plane in
dimension $N$, the intersection dimension exceeds its minimum $(p+v-N)_+$ by at least $t\ge1$
with probability at most
\begin{equation}\label{spn:eq:planetail}
 C2^{-t(t+\abs{N-p-v})},
\end{equation}
by counting $t$-planes in the intersection when $p+v\le N$, and using annihilators otherwise. The
minima of $a,a',g,g'$ are $(r-e)_+$, $0$, $n-r-e$ and $(e-r)_+$. If $e\le r/2$, projection to $Y$
gives $d\le a+g'$, so $d\ge r$ requires $a-(r-e)$ or $g'$ to reach $\ceil{e/2}$, and in both cases
$\abs{N-p-v}=r-e$ in \eqref{spn:eq:planetail}; this gives the first bound. If $e\ge r/2$, apply
\eqref{spn:eq:planetail} to the four dimensions and, conditionally, to the last term of
\eqref{spn:eq:residual}; outside an event of probability $O(2^{-t^2})$ all five excesses are less
than $t$, and then $d\le\max(a+a',g+g'-(n-2r))+t\le(r-e)_++3t$. Taking $t=\floor{r/8}$ gives
$d<r$.
\end{proof}

\begin{theorem}[The probability of a large principal kernel; \stU]\label{spn:thm:rare}
Let $r=r(n)\le\floor{n/2}$ with $r/\sqrt n\to\infty$. Conditional on $\nu(Q)\ge r$, with
probability tending to one exactly one complementary pair has nullity at least $r$, and its
nullity is exactly $r$. If moreover $n-2r\to\infty$, then
\begin{equation}\label{spn:eq:rare}
 \spnP(\nu(Q)\ge r)\sim\frac{B_{n,r}}{2\abs{\GL_r(\F_2)}}\sim\frac{B_{n,r}}{2\gamma}\,2^{-r^2},
\end{equation}
and choosing a side of the exceptional pair by a fair coin gives, in total variation, the uniform
distribution on the $B_{n,r}$ feasible sets.
\end{theorem}

\begin{proof}
For fixed $S$ there are at most $2\binom ne$ oriented sets at folded distance $e$. Summing the
first bound of \cref{spn:lem:twopairs} with $\binom ne\le n^e$ and $r-e\ge r/2$, and the second over
at most $2^n$ sets, gives
$\spnP_{\mathrm{sd}}(\exists\,T\notin\set{S,S^c}:d_T\ge r\mid d_S=r)\ll n2^{-r/4}+2^{n-cr^2}=o(1)$,
since $r/\sqrt n\to\infty$. Multiplying by \eqref{spn:eq:sdone} and dividing by $p_n$ returns to
isometries: uniformly in feasible $S$,
$\spnP(b_S=r,\ \exists\,T\notin\set{S,S^c}:b_T\ge r)=o(2^{-r^2})$, while
$\spnP(b_S>r)=O(2^{-(r+1)^2})$ by \cref{spn:lem:onesupport}. Summing over feasible pairs,
\begin{equation}\label{spn:eq:raresum}
 \spnP(\nu(Q)\ge r)=(1+o(1))\frac12\sum_{r\le\abs S\le n-r}\spnP(b_S=r)\asymp B_{n,r}2^{-r^2},
\end{equation}
the lower bound summing disjoint events on which the indicated pair is unique with nullity $r$,
and the upper bound summing single-pair events; this proves conditional uniqueness. If
$n-2r\to\infty$, remove the $\ell$ layers $\abs S$ nearest either endpoint, with $\ell\to\infty$
and $\ell=o(\min(n-2r,\sqrt n))$; the removed proportion of feasible sets is
$O(\ell/\min(n-2r,\sqrt n))$, since each layer is at most a central binomial coefficient and the
middle $\asymp\min(n-2r,\sqrt n)$ layers are each at least a constant times it. On the retained
layers \cref{spn:lem:onesupport} gives a uniform equivalent $\abs{\GL_r(\F_2)}^{-1}$, and
elsewhere an upper bound of the same order, which gives \eqref{spn:eq:rare} and the
total-variation statement.
\end{proof}

When $(n-2r)/\sqrt n\to t\in(0,\infty)$, $B_{n,r}/2^n$ tends to $2\Phi(t)-1$ with $\Phi$ the
standard normal distribution function, and for $1\ll n-2r\ll\sqrt n$,
$B_{n,r}\sim(n-2r+1)2^n\sqrt{2/(\pi n)}$: the shortening of the feasible interval changes the
prefactor while the exponent stays $-r^2$.

\begin{corollary}[The linear-deviation asymptotic; \stU]\label{spn:cor:no-linear}
For fixed $0<\alpha<1/2$ and $r=\floor{\alpha n}$,
\[
 \log_2\spnP\set{\nu(Q)\ge r}=-r^2+n-1-\log_2\gamma+o(1).
\]
In particular $n^{-2}\log_2\spnP\set{\nu(Q)\ge\floor{\alpha n}}\to-\alpha^2$ for every
$0<\alpha\le1/2$, the endpoint $\alpha=1/2$ following from \cref{spn:thm:saturation}, and the
high-probability linear lower bound proposed in \cite{Vaas2026Principal} is impossible.
\end{corollary}

\begin{proof}
The feasible sets have proportion tending to one, so $B_{n,r}\sim2^n$ in \eqref{spn:eq:rare}.
Keeping the integer $r$ inside the square matters at this precision.
\end{proof}

\begin{theorem}[Saturation; \stU]\label{spn:thm:saturation}
As $m\to\infty$,
\begin{align*}
 n&=2m: &\spnP(\nu(Q)=m)&\sim\frac a2\binom{2m}m2^{-m^2},\\
 n&=2m+1: &\spnP(\nu(Q)=m)&\sim a\binom{2m+1}m2^{-m^2}.
\end{align*}
\end{theorem}

\begin{proof}
The summation \eqref{spn:eq:raresum} holds also for bounded $n-2r$. For $n=2m$ and $\abs S=m$,
$b_S=m$ exactly when $Q$ maps $V_S$ onto $V_{S^c}$, with probability
$\abs{\spnOrth_m}^2/\abs{\spnOrth_{2m}}$; take the $\binom{2m}m/2$ complementary pairs and use
\eqref{spn:eq:orth-order}. For $n=2m+1$ take the side of size $m$: the image of $V_S$ must be a
standard $m$-plane of the $(m+1)$-dimensional space $V_{S^c}$, whose orthogonal complement there is
spanned by a norm-one vector $v$. There are $2^m-\mathbf 1_{m\text{ even}}$ admissible $v$,
excluding $\mathbf j_{m+1}$ when it is odd, since its orthogonal hyperplane is alternating, and
for $m\ge2$ this counts the targets exactly. The probability for one support is
$(2^m-\mathbf 1_{m\text{ even}})\abs{\spnOrth_m}\abs{\spnOrth_{m+1}}/\abs{\spnOrth_{2m+1}}$;
multiply by $\binom{2m+1}m$ and use \eqref{spn:eq:orth-order}.
\end{proof}

\begin{proposition}[Exact count at saturation in even dimension; \stU]\label{spn:prop:saturated-count}
Let $o_j=\abs{\spnOrth_j}$, $o_0=1$, and let $N_{2m}$ be the number of $Q\in\spnOrth_{2m}$ with
$\nu(Q)=m$. Then
\[
 N_{2m}=(2m)!\,[z^m]\Bigl(\sum_{j\ge0}\Bigl(\frac{o_j}{j!}\Bigr)^2z^j\Bigr)^{1/2},
\]
the square root being the formal power series with constant term one. The first values are
$N_2=1$, $N_4=9$, $N_6=225$ and $N_8=71505$.
\end{proposition}

\begin{proof}
Associate to $Q$ the graph on $[2m]$ with an edge $ij$, loops included, whenever $Q_{ij}$ or
$Q_{ji}$ is nonzero. Saturation means a partition into two classes with zero diagonal blocks,
that is a proper two-colouring. On each connected component the two classes have equal size,
since the corresponding diagonal component of $Q$ is invertible and interchanges the two
coordinate spaces, so a properly coloured matrix on $2j$ labels has $\binom{2j}jo_j^2$
possibilities, the two off-diagonal blocks being arbitrary isometries. Each uncoloured bipartite
component has exactly two colourings, so the exponential generating series of coloured
configurations is the square of that of uncoloured ones. Substituting $z=x^2$ gives the formula.
\end{proof}

In dimension six, $225$ of the $23{,}040$ matrices attain the maximum three; this finite count
includes collisions of distinct complementary pairs, which \cref{spn:thm:rare} shows to be
negligible as the dimension grows. The group orders, nullity histograms through dimension six,
the saturated counts through dimension twelve and the quotient-intersection identity
\eqref{spn:eq:residual} are checked by exact
computation.\footnote{\raggedright\artifact{large_kernels_and_dirichlet_zeros/check_orthogonal.txt}}

The proofs do not treat independent-entry matrices as interchangeable with orthogonal ones: the
fixed vector $\mathbf j_n$ and the regularity condition are essential. Exhaustive verification
for $1\le n\le6$ (group orders, the deterministic maximum $\floor{n/2}$, complete nullity
histograms) is
recorded.\footnote{\raggedright\artifact{grh_method_transfers_orthogonal_jets_2026-11-26/checks/binary_orthogonal_checks.json}}
Enumeration of the orthogonal groups in dimensions $2$ to $5$ compares frame orbits with complete
Gram fibres for frames of up to three columns and tests \eqref{spn:eq:comp}; over all $65{,}536$
binary $4\times4$ matrices, the one-dimensional witnesses on the six two-element supports have
$\spnE Z=9/2$ and $\spnE Z^2=201/8$, as the row-wise and overlap counts
predict.\footnote{\raggedright\artifact{orthogonal_nullity_and_prime_products/checks/check_matrices_output.txt}}

%======================================================================
